% =========================================================
% mechlib/mechlib.tex
% ZENTRALER LOADER / OEFFENTLICHER EINSTIEGSPUNKT
% =========================================================
% Diese Datei ist die einzige Datei, die ein Anwender normalerweise direkt
% einbindet. Sie laedt die benoetigten LaTeX-/TikZ-Abhaengigkeiten und danach
% alle Teilmodule der Bibliothek in einer definierten Reihenfolge.
%
% Typische Verwendung bei einem Ordner "mechlib" neben dem Hauptdokument:
%   % =========================================================
% mechlib/mechlib.tex
% ZENTRALER LOADER / OEFFENTLICHER EINSTIEGSPUNKT
% =========================================================
% Diese Datei ist die einzige Datei, die ein Anwender normalerweise direkt
% einbindet. Sie laedt die benoetigten LaTeX-/TikZ-Abhaengigkeiten und danach
% alle Teilmodule der Bibliothek in einer definierten Reihenfolge.
%
% Typische Verwendung bei einem Ordner "mechlib" neben dem Hauptdokument:
%   % =========================================================
% mechlib/mechlib.tex
% ZENTRALER LOADER / OEFFENTLICHER EINSTIEGSPUNKT
% =========================================================
% Diese Datei ist die einzige Datei, die ein Anwender normalerweise direkt
% einbindet. Sie laedt die benoetigten LaTeX-/TikZ-Abhaengigkeiten und danach
% alle Teilmodule der Bibliothek in einer definierten Reihenfolge.
%
% Typische Verwendung bei einem Ordner "mechlib" neben dem Hauptdokument:
%   % =========================================================
% mechlib/mechlib.tex
% ZENTRALER LOADER / OEFFENTLICHER EINSTIEGSPUNKT
% =========================================================
% Diese Datei ist die einzige Datei, die ein Anwender normalerweise direkt
% einbindet. Sie laedt die benoetigten LaTeX-/TikZ-Abhaengigkeiten und danach
% alle Teilmodule der Bibliothek in einer definierten Reihenfolge.
%
% Typische Verwendung bei einem Ordner "mechlib" neben dem Hauptdokument:
%   \input{mechlib/mechlib}
%
% Liegt die Bibliothek an einem anderen Ort, kann vor dem Laden die Variable
% \mechlibpath gesetzt werden:
%   \newcommand{\mechlibpath}{../../../support/mechlib}
%   \input{\mechlibpath/mechlib}
%
% Lade-Reihenfolge:
%   1. core/        Farben, Layer, Stile und globale Geometriekonventionen
%   2. elements/    mechanische Grundelemente (Koerper, Lager, Verbinder, Staebe)
%   3. annotations/ Lasten, Masse, Koordinaten und mechanische Annotationen
%
% WICHTIG:
% Die Modulreihenfolge ist absichtlich festgelegt. Elemente greifen auf Farben,
% Layer und Stile aus core/ zurueck; Annotationen koennen wiederum die bereits
% definierten Grundstile und Pics verwenden.
% =========================================================
% -------------------------------------------------------------------------
% Externe Pakete und TikZ-Bibliotheken
% -------------------------------------------------------------------------
% \RequirePackage ist bewusst statt \usepackage gewaehlt, da mechlib.tex selbst
% wie eine Bibliotheks-/Paketdatei verwendet wird.
\RequirePackage{tikz}
\RequirePackage{amsmath}
\RequirePackage{siunitx}
\RequirePackage{xcolor}
\usetikzlibrary{
  arrows.meta,
  calc,
  patterns,
  patterns.meta,
  intersections,
  angles,
  quotes,
  decorations.pathmorphing,
  backgrounds
}
% Base path: override before loading when mechlib is stored elsewhere.
% -------------------------------------------------------------------------
% Aufloesung des Bibliothekspfades
% -------------------------------------------------------------------------
% \providecommand respektiert eine bereits im Hauptdokument gesetzte Definition.
\providecommand{\mechlibpath}{mechlib}
% -------------------------------------------------------------------------
% Kernmodule: muessen vor allen Elementen geladen werden.
% -------------------------------------------------------------------------
\input{\mechlibpath/core/colors}
\input{\mechlibpath/core/layers}
\input{\mechlibpath/core/styles}
\input{\mechlibpath/core/defaults}
% -------------------------------------------------------------------------
% Atomare mechanische Elemente.
% -------------------------------------------------------------------------
\input{\mechlibpath/elements/bodies}
\input{\mechlibpath/elements/supports}
\input{\mechlibpath/elements/connectors}
\input{\mechlibpath/elements/trusses}
% -------------------------------------------------------------------------
% Annotationen und mechanische Symbolgruppen.
% -------------------------------------------------------------------------
\input{\mechlibpath/annotations/loads}
\input{\mechlibpath/annotations/dimensions}
\input{\mechlibpath/annotations/coordinates}
\input{\mechlibpath/annotations/mechanics}

%
% Liegt die Bibliothek an einem anderen Ort, kann vor dem Laden die Variable
% \mechlibpath gesetzt werden:
%   \newcommand{\mechlibpath}{../../../support/mechlib}
%   % =========================================================
% mechlib/mechlib.tex
% ZENTRALER LOADER / OEFFENTLICHER EINSTIEGSPUNKT
% =========================================================
% Diese Datei ist die einzige Datei, die ein Anwender normalerweise direkt
% einbindet. Sie laedt die benoetigten LaTeX-/TikZ-Abhaengigkeiten und danach
% alle Teilmodule der Bibliothek in einer definierten Reihenfolge.
%
% Typische Verwendung bei einem Ordner "mechlib" neben dem Hauptdokument:
%   \input{mechlib/mechlib}
%
% Liegt die Bibliothek an einem anderen Ort, kann vor dem Laden die Variable
% \mechlibpath gesetzt werden:
%   \newcommand{\mechlibpath}{../../../support/mechlib}
%   \input{\mechlibpath/mechlib}
%
% Lade-Reihenfolge:
%   1. core/        Farben, Layer, Stile und globale Geometriekonventionen
%   2. elements/    mechanische Grundelemente (Koerper, Lager, Verbinder, Staebe)
%   3. annotations/ Lasten, Masse, Koordinaten und mechanische Annotationen
%
% WICHTIG:
% Die Modulreihenfolge ist absichtlich festgelegt. Elemente greifen auf Farben,
% Layer und Stile aus core/ zurueck; Annotationen koennen wiederum die bereits
% definierten Grundstile und Pics verwenden.
% =========================================================
% -------------------------------------------------------------------------
% Externe Pakete und TikZ-Bibliotheken
% -------------------------------------------------------------------------
% \RequirePackage ist bewusst statt \usepackage gewaehlt, da mechlib.tex selbst
% wie eine Bibliotheks-/Paketdatei verwendet wird.
\RequirePackage{tikz}
\RequirePackage{amsmath}
\RequirePackage{siunitx}
\RequirePackage{xcolor}
\usetikzlibrary{
  arrows.meta,
  calc,
  patterns,
  patterns.meta,
  intersections,
  angles,
  quotes,
  decorations.pathmorphing,
  backgrounds
}
% Base path: override before loading when mechlib is stored elsewhere.
% -------------------------------------------------------------------------
% Aufloesung des Bibliothekspfades
% -------------------------------------------------------------------------
% \providecommand respektiert eine bereits im Hauptdokument gesetzte Definition.
\providecommand{\mechlibpath}{mechlib}
% -------------------------------------------------------------------------
% Kernmodule: muessen vor allen Elementen geladen werden.
% -------------------------------------------------------------------------
\input{\mechlibpath/core/colors}
\input{\mechlibpath/core/layers}
\input{\mechlibpath/core/styles}
\input{\mechlibpath/core/defaults}
% -------------------------------------------------------------------------
% Atomare mechanische Elemente.
% -------------------------------------------------------------------------
\input{\mechlibpath/elements/bodies}
\input{\mechlibpath/elements/supports}
\input{\mechlibpath/elements/connectors}
\input{\mechlibpath/elements/trusses}
% -------------------------------------------------------------------------
% Annotationen und mechanische Symbolgruppen.
% -------------------------------------------------------------------------
\input{\mechlibpath/annotations/loads}
\input{\mechlibpath/annotations/dimensions}
\input{\mechlibpath/annotations/coordinates}
\input{\mechlibpath/annotations/mechanics}

%
% Lade-Reihenfolge:
%   1. core/        Farben, Layer, Stile und globale Geometriekonventionen
%   2. elements/    mechanische Grundelemente (Koerper, Lager, Verbinder, Staebe)
%   3. annotations/ Lasten, Masse, Koordinaten und mechanische Annotationen
%
% WICHTIG:
% Die Modulreihenfolge ist absichtlich festgelegt. Elemente greifen auf Farben,
% Layer und Stile aus core/ zurueck; Annotationen koennen wiederum die bereits
% definierten Grundstile und Pics verwenden.
% =========================================================
% -------------------------------------------------------------------------
% Externe Pakete und TikZ-Bibliotheken
% -------------------------------------------------------------------------
% \RequirePackage ist bewusst statt \usepackage gewaehlt, da mechlib.tex selbst
% wie eine Bibliotheks-/Paketdatei verwendet wird.
\RequirePackage{tikz}
\RequirePackage{amsmath}
\RequirePackage{siunitx}
\RequirePackage{xcolor}
\usetikzlibrary{
  arrows.meta,
  calc,
  patterns,
  patterns.meta,
  intersections,
  angles,
  quotes,
  decorations.pathmorphing,
  backgrounds
}
% Base path: override before loading when mechlib is stored elsewhere.
% -------------------------------------------------------------------------
% Aufloesung des Bibliothekspfades
% -------------------------------------------------------------------------
% \providecommand respektiert eine bereits im Hauptdokument gesetzte Definition.
\providecommand{\mechlibpath}{mechlib}
% -------------------------------------------------------------------------
% Kernmodule: muessen vor allen Elementen geladen werden.
% -------------------------------------------------------------------------
% =========================================================
% mechlib/core/colors.tex
% ZENTRALE FARBPALETTE
% =========================================================
% Alle bibliotheksweiten Farben werden ausschliesslich hier definiert. Dadurch
% kann der visuelle Charakter der gesamten mechlib an einer Stelle angepasst
% werden, ohne einzelne Elemente zu veraendern.
%
% Konvention:
%   mechBodyColor    Fuellung allgemeiner Koerper und Fachwerkgelenke
%   mechSupportColor Fuellung von Lagerkoerpern
%   mechLoadColor    Kraefte, Momente und Streckenlasten
%   mechAxisColor    Achsen und Koordinatensysteme
%   mechRodColor     historische Stabfarbe; aktuelle Fachwerkstaebe sind schwarz
%
% Die Namen sind Teil der internen visuellen API und sollten bevorzugt statt
% harter RGB-Werte in den Elementdateien verwendet werden.
% =========================================================
% -------------------------------------------------------------------------
% Farbdefinitionen
% -------------------------------------------------------------------------
% RGB wird bewusst explizit verwendet, damit die Darstellung unabhaengig von
% Dokumentklassen und externen Farbmodellen reproduzierbar bleibt.
\definecolor{mechBodyColor}{RGB}{180,180,180}
\definecolor{mechSupportColor}{RGB}{205,205,205}
\definecolor{mechLoadColor}{RGB}{178,25,25}
\definecolor{mechAxisColor}{RGB}{6,13,141}
\definecolor{mechRodColor}{RGB}{130,130,130}

% =========================================================
% mechlib/core/layers.tex
% ZEICHENEBENEN / PGF-LAYER
% =========================================================
% Fachwerkstaebe und andere hinterlegte Geometrien koennen auf dem Layer
% "mech background" gezeichnet werden. Der normale TikZ-Layer "main" liegt
% darueber. Dadurch koennen z. B. gefuellte Gelenke sauber vor den Staeben
% erscheinen, ohne dass die Zeichenreihenfolge im Anwendungsdokument aufwendig
% kontrolliert werden muss.
% =========================================================
% Zuerst wird der benannte Hintergrundlayer deklariert. Anschliessend wird die
% globale Layerreihenfolge explizit gesetzt: Hintergrund vor Hauptlayer.
\pgfdeclarelayer{mech background}
\pgfsetlayers{mech background,main}

% =========================================================
% mechlib/core/styles.tex
% ZENTRALE, NAMENSRAUMGESCHUETZTE TIKZ-STILE
% =========================================================
% Diese Datei definiert ausschliesslich visuelle Grundstile. Die eigentliche
% Geometrie eines Elements gehoert in elements/ bzw. annotations/.
%
% Alle oeffentlichen Bibliotheksstile beginnen mit "mech/". Das reduziert
% Namenskollisionen mit dem Hauptdokument und macht im TikZ-Code sichtbar,
% welche Formatierung aus der mechlib stammt.
%
% Designregel:
%   - Farben werden aus core/colors.tex bezogen.
%   - Geometrien werden NICHT hier definiert.
%   - Linienstarken werden hier zentral gehalten, soweit sie elementuebergreifend
%     konsistent sein sollen.
% =========================================================
% -------------------------------------------------------------------------
% Visuelle Basisstile
% -------------------------------------------------------------------------
% Einzelne Pics duerfen diese Stile ueber ihre jeweiligen "style"-Keys ersetzen.
\tikzset{
  mech/body/.style={
    draw=black,
    fill=mechBodyColor,
    line width=.6pt
  },
  mech/support/.style={
    draw=black,
    fill=mechSupportColor,
    line width=.6pt
  },
  mech/connector/.style={
    draw=black,
    line width=.6pt
  },
  mech/rod/.style={
    draw=black,
    line width=.85pt
  },
  mech/load/.style={
    draw=mechLoadColor,
    text=mechLoadColor,
    line width=1pt
  },
  mech/load fill/.style={
    fill=mechLoadColor,
    fill opacity=.16,
    draw=none
  },
  mech/dimension/.style={
    {Latex[length=2mm]}-{Latex[length=2mm]},
    line width=.5pt
  },
  mech/axis/.style={
    draw=mechAxisColor,
    text=mechAxisColor,
    line width=.6pt
  },
  mech/label/.style={
    inner sep=1pt
  }
}

% =========================================================
% mechlib/core/defaults.tex
% GLOBALE STANDARDWERTE UND GEMEINSAME GEOMETRIEKONVENTIONEN
% =========================================================
% Hier stehen Werte, die von mehreren Modulen gemeinsam benutzt werden. Solche
% Konstanten verhindern, dass geometrisch zusammengehoerige Elemente lokal mit
% leicht unterschiedlichen Zahlen definiert werden.
%
% Beispiel: \mechConnectionOffset legt die vertikale Lage paralleler
% Anschlussanker fest. Masse und Einspannung verwenden exakt denselben Wert;
% dadurch liegen obere/mittlere/untere Anschluesse bei gleicher Mittelpunktlage
% garantiert auf gleicher Hoehe.
% =========================================================
% Globale Schriftvorgabe fuer jedes TikZ-Bild. "append style" ergaenzt
% vorhandene Bildstile, statt sie zu ueberschreiben.
\tikzset{
  every picture/.append style={font=\fontsize{10}{10}\selectfont},
}

% Standardwinkel fuer technische Bodenschraffuren.
\def\mechHatchAngle{45}

% Shared vertical distance of parallel connector anchors from the middle anchor.
% Wall and mass use this exact value so upper/middle/lower are aligned whenever
% their center points lie on the same horizontal line.
% Gemeinsamer vertikaler Versatz fuer parallele Anschlussanker.
\def\mechConnectionOffset{2.5mm}

% -------------------------------------------------------------------------
% Atomare mechanische Elemente.
% -------------------------------------------------------------------------
% =========================================================
% mechlib/elements/bodies.tex
% ATOMARE KOERPER: MASSE, PUNKTMASSE, WAGEN UND BALKEN
% =========================================================
% Die hier definierten TikZ-Pics stellen mechanische Grundkoerper bereit. Jedes
% Pic besitzt eine eigene pgfkeys-Familie unter /mech/... und kann dadurch mit
% benannten Optionen konfiguriert werden.
%
% Grundprinzip der API:
%   \pic (name) at (<Ort>) {mech/<element>={<key>=<wert>,...}};
%
% Benannte Pics exportieren Koordinaten bzw. werden bei node-basierten Koerpern
% auf den internen Shape-Knoten aliasiert. So koennen regulaere TikZ-Anker wie
% (m1.east), (m1.north west) usw. direkt benutzt werden.
%
% Oeffentliche Pics in dieser Datei:
%   mech/mass       rechteckiger Starrkoerper bzw. konzentrierte Masse
%   mech/point-mass kreisfoermige Punktmasse
%   mech/cart       Wagen mit zwei Raedern und optionalem Boden
%   mech/beam       Balken zwischen zwei beliebigen Punkten
% =========================================================
\pgfkeys{
  /mech/mass/.is family,
  /mech/mass,
  width/.initial=1.5cm,
  height/.initial=1cm,
  label/.initial=,
  style/.initial=mech/body,
}
\tikzset{
  pics/mech/mass/.style args={#1}{
    code={
      \pgfkeys{/mech/mass/.cd,#1}
      \node[
        \pgfkeysvalueof{/mech/mass/style},
        minimum width=\pgfkeysvalueof{/mech/mass/width},
        minimum height=\pgfkeysvalueof{/mech/mass/height},
        inner sep=0pt
      ] (-shape) at (0,0) {\pgfkeysvalueof{/mech/mass/label}};

      % Named body pics behave like regular TikZ nodes.
      % Example: \pic (m1) ... then (m1.east), (m1.south east), ...
      \edef\mechPicName{\pgfkeysvalueof{/tikz/name prefix}}
      \ifx\mechPicName\pgfutil@empty\else
        \pgfnodealias{\mechPicName}{\mechPicName-shape}
      \fi

      \coordinate (-center) at (-shape.center);

      % Three standardized connection levels on each vertical side.
      % The shared \mechConnectionOffset is also used by fixed supports.
      \coordinate (-west-upper)  at ($(-shape.west)+(0,\mechConnectionOffset)$);
      \coordinate (-west-middle) at (-shape.west);
      \coordinate (-west-lower)  at ($(-shape.west)+(0,-\mechConnectionOffset)$);

      \coordinate (-east-upper)  at ($(-shape.east)+(0,\mechConnectionOffset)$);
      \coordinate (-east-middle) at (-shape.east);
      \coordinate (-east-lower)  at ($(-shape.east)+(0,-\mechConnectionOffset)$);

      % Generic connection aliases
      \coordinate (-connection-upper)  at (-east-upper);
      \coordinate (-connection-middle) at (-east-middle);
      \coordinate (-connection-lower)  at (-east-lower);

      \coordinate (-north) at (-shape.north);
      \coordinate (-south) at (-shape.south);
    }
  },
  pics/mech/mass/.default={},
}
\pgfkeys{
  /mech/point-mass/.is family,
  /mech/point-mass,
  radius/.initial=2.2mm,
  label/.initial=,
  style/.initial=mech/body,
}
\tikzset{
  pics/mech/point-mass/.style args={#1}{
    code={
      \pgfkeys{/mech/point-mass/.cd,#1}
      \edef\mechPointMassRadius{\pgfkeysvalueof{/mech/point-mass/radius}}
      \edef\mechPointMassLabel{\pgfkeysvalueof{/mech/point-mass/label}}
      \edef\mechPointMassStyle{\pgfkeysvalueof{/mech/point-mass/style}}
      \node[
        \mechPointMassStyle,
        circle,
        minimum size=2\mechPointMassRadius,
        inner sep=0pt
      ] (-shape) at (0,0) {\mechPointMassLabel};

      \edef\mechPicName{\pgfkeysvalueof{/tikz/name prefix}}
      \ifx\mechPicName\pgfutil@empty\else
        \pgfnodealias{\mechPicName}{\mechPicName-shape}
      \fi

      \coordinate (-center) at (-shape.center);
    }
  },
  pics/mech/point-mass/.default={},
}
% -------------------------------------------------------------------------
% Wagen
% -------------------------------------------------------------------------
% Der Boolean "ground" wird als TeX-if gespeichert. Dadurch kann die optionale
% Bodendarstellung ohne Stringvergleich ein- bzw. ausgeschaltet werden.
\newif\ifmechCartGround
\pgfkeys{
  /mech/cart/.is family,
  /mech/cart,
  width/.initial=2cm,
  height/.initial=1cm,
  wheel radius/.initial=2mm,
  label/.initial=,
  ground/.is if=mechCartGround,
  ground/.default=true,
  ground=true,
}
\tikzset{
  pics/mech/cart/.style args={#1}{
    code={
      \pgfkeys{/mech/cart/.cd,#1}
      \edef\mechCartWidth{\pgfkeysvalueof{/mech/cart/width}}
      \edef\mechCartHeight{\pgfkeysvalueof{/mech/cart/height}}
      \edef\mechCartWheelRadius{\pgfkeysvalueof{/mech/cart/wheel radius}}
      \node[
        mech/body,
        minimum width=\mechCartWidth,
        minimum height=\mechCartHeight,
        inner sep=0pt
      ] (-shape) at (0,0) {\pgfkeysvalueof{/mech/cart/label}};

      \edef\mechPicName{\pgfkeysvalueof{/tikz/name prefix}}
      \ifx\mechPicName\pgfutil@empty\else
        \pgfnodealias{\mechPicName}{\mechPicName-shape}
      \fi

      \coordinate (-center) at (-shape.center);
      \coordinate (-west) at (-shape.west);
      \coordinate (-east) at (-shape.east);
      \coordinate (-north) at (-shape.north);
      \coordinate (-south) at (-shape.south);
      \draw[mech/body]
        (-.3\mechCartWidth,-.5\mechCartHeight-\mechCartWheelRadius)
        circle[radius=\mechCartWheelRadius];
      \draw[mech/body]
        (.3\mechCartWidth,-.5\mechCartHeight-\mechCartWheelRadius)
        circle[radius=\mechCartWheelRadius];
      \ifmechCartGround
        \draw
          (-.5\mechCartWidth,-.5\mechCartHeight-2\mechCartWheelRadius)
          --
          (.5\mechCartWidth,-.5\mechCartHeight-2\mechCartWheelRadius);
        \fill[
          pattern={Lines[angle=\mechHatchAngle,distance=1pt]},
          pattern color=black
        ]
          (-.5\mechCartWidth,-.5\mechCartHeight-2\mechCartWheelRadius)
          rectangle
          (.5\mechCartWidth,-.5\mechCartHeight-2\mechCartWheelRadius-3pt);
      \fi
    }
  },
  pics/mech/cart/.default={},
}
\pgfkeys{
  /mech/beam/.is family,
  /mech/beam,
  from/.initial={(0,0)},
  to/.initial={(1,0)},
  height/.initial=3mm,
  style/.initial=mech/body,
}
\tikzset{
  pics/mech/beam/.style args={#1}{
    code={
      \pgfkeys{/mech/beam/.cd,#1}
      \edef\mechBeamFrom{\pgfkeysvalueof{/mech/beam/from}}
      \edef\mechBeamTo{\pgfkeysvalueof{/mech/beam/to}}
      \edef\mechBeamHeight{\pgfkeysvalueof{/mech/beam/height}}
      \edef\mechBeamStyle{\pgfkeysvalueof{/mech/beam/style}}
      \path let
        \p1=\mechBeamFrom,
        \p2=\mechBeamTo,
        \n1={veclen(\x2-\x1,\y2-\y1)},
        \n2={(\y2-\y1)/\n1*\mechBeamHeight/2},
        \n3={-(\x2-\x1)/\n1*\mechBeamHeight/2}
      in
        coordinate (-A) at (\x1+\n2,\y1+\n3)
        coordinate (-B) at (\x2+\n2,\y2+\n3)
        coordinate (-C) at (\x2-\n2,\y2-\n3)
        coordinate (-D) at (\x1-\n2,\y1-\n3)
        coordinate (-start) at (\x1,\y1)
        coordinate (-end) at (\x2,\y2)
        coordinate (-center) at ({(\x1+\x2)/2},{(\y1+\y2)/2});
      \draw[\mechBeamStyle] (-A)--(-B)--(-C)--(-D)--cycle;
    }
  },
  pics/mech/beam/.default={},
}

% =========================================================
% mechlib/elements/supports.tex
% ATOMARE LAGER: FESTLAGER, LOSLAGER UND EINSPANNUNG
% =========================================================
% Diese Datei kapselt die Lagergeometrie. Fest- und Loslager teilen sich die
% Key-Familie /mech/support-common, damit Breite, Hoehe, Schraffur und Label
% konsistent parametrisiert werden koennen.
%
% Oeffentliche Pics:
%   mech/pinned-support Festlager / gelenkiges Lager
%   mech/roller-support Loslager mit sichtbarem Abstand zum Boden
%   mech/fixed-support  Einspannung, optional gedreht
%
% Aktuelle Fachwerk-Konvention:
%   - Die Dreiecksspitze liegt exakt im Mittelpunkt des Gelenks (0,0).
%   - Das Dreieck ist gegenueber der frueheren Version kompakt skaliert.
%   - Bodenlinie und Schraffur verwenden dieselben x-Grenzen wie die
%     Dreiecksgrundseite; beim Loslager wird die sichtbare Linienbreite explizit
%     beruecksichtigt.
%   - Das Gelenk selbst ist mit mechBodyColor gefuellt.
%
% Die pgfmathsetlengthmacro-Aufrufe sind absichtlich verwendet: LaTeX-Laengen
% duerfen nicht durch Textersatz wie ".5\\macro" skaliert werden, da dies bei
% expandierenden Dimensionsmakros zu sichtbarem Resttext (z. B. ".8mm") fuehren
% kann. Alle abgeleiteten Abmessungen werden deshalb numerisch berechnet.
% =========================================================
\pgfkeys{
  /mech/support-common/.is family,
  /mech/support-common,
  width/.initial=2.72mm,
  height/.initial=2.24mm,
  hatch angle/.initial=45,
  label/.initial=,
  label position/.initial=above,
}
\tikzset{
  % -----------------------------------------------------------------------
  % Festlager
  % -----------------------------------------------------------------------
  pics/mech/pinned-support/.style args={#1}{
    code={
      \pgfkeys{/mech/support-common/.cd,#1}
      \edef\mechSupportWidth{\pgfkeysvalueof{/mech/support-common/width}}
      \edef\mechSupportHeight{\pgfkeysvalueof{/mech/support-common/height}}
      \edef\mechSupportHatch{\pgfkeysvalueof{/mech/support-common/hatch angle}}
      \pgfmathsetlengthmacro{\mechSupportHalfWidth}{0.5*\mechSupportWidth}
      % Die Dreieckskontur ragt an der Grundseite durch die halbe Linienbreite
      % minimal ueber die geometrische Breite hinaus. Fuer identische sichtbare
      % Breite wird die Bodenlinie daher um 0.3pt je Seite erweitert.
      \pgfmathsetlengthmacro{\mechSupportGroundHalfWidth}{\mechSupportHalfWidth+0.3pt}
      \pgfmathsetlengthmacro{\mechSupportApexY}{0mm}
      \pgfmathsetlengthmacro{\mechSupportHatchDepth}{0.256*\mechSupportHeight}

      \coordinate (-connection) at (0,0);
      \coordinate (-center) at (0,0);

      \draw[mech/support]
        (-\mechSupportHalfWidth,-\mechSupportHeight)
        -- (0,\mechSupportApexY)
        -- (\mechSupportHalfWidth,-\mechSupportHeight)
        -- cycle;
      \draw[mech/body] (0,0) circle[radius=.6mm];
      \draw[black,line width=.6pt,line cap=butt]
        (-\mechSupportGroundHalfWidth,-\mechSupportHeight)
        -- (\mechSupportGroundHalfWidth,-\mechSupportHeight);
      \fill[
        pattern={Lines[angle=\mechSupportHatch,distance=1pt]},
        pattern color=black
      ]
        (-\mechSupportGroundHalfWidth,-\mechSupportHeight)
        rectangle
        (\mechSupportGroundHalfWidth,{-(\mechSupportHeight+\mechSupportHatchDepth)});

      \node[\pgfkeysvalueof{/mech/support-common/label position}]
        at (0,0) {\pgfkeysvalueof{/mech/support-common/label}};
    }
  },
  pics/mech/pinned-support/.default={},
  % -----------------------------------------------------------------------
  % Loslager
  % -----------------------------------------------------------------------
  pics/mech/roller-support/.style args={#1}{
    code={
      \pgfkeys{/mech/support-common/.cd,#1}
      \edef\mechSupportWidth{\pgfkeysvalueof{/mech/support-common/width}}
      \edef\mechSupportHeight{\pgfkeysvalueof{/mech/support-common/height}}
      \edef\mechSupportHatch{\pgfkeysvalueof{/mech/support-common/hatch angle}}
      \pgfmathsetlengthmacro{\mechSupportHalfWidth}{0.5*\mechSupportWidth}
      % The triangle outline extends by half its line width at the base corners.
      % Add 0.3pt so the roller ground has the same visible total width.
      \pgfmathsetlengthmacro{\mechSupportGroundHalfWidth}{\mechSupportHalfWidth+0.3pt}
      \pgfmathsetlengthmacro{\mechSupportApexY}{0mm}
      \pgfmathsetlengthmacro{\mechSupportRollerGap}{0.272*\mechSupportHeight}
      \pgfmathsetlengthmacro{\mechSupportHatchDepth}{0.256*\mechSupportHeight}
      \pgfmathsetlengthmacro{\mechSupportGroundY}{\mechSupportHeight+\mechSupportRollerGap}

      \coordinate (-connection) at (0,0);
      \coordinate (-center) at (0,0);

      \draw[mech/support]
        (-\mechSupportHalfWidth,-\mechSupportHeight)
        -- (0,\mechSupportApexY)
        -- (\mechSupportHalfWidth,-\mechSupportHeight)
        -- cycle;
      \draw[mech/body] (0,0) circle[radius=.6mm];
      \draw[black,line width=.6pt,line cap=butt]
        (-\mechSupportGroundHalfWidth,-\mechSupportGroundY)
        -- (\mechSupportGroundHalfWidth,-\mechSupportGroundY);
      \fill[
        pattern={Lines[angle=\mechSupportHatch,distance=1pt]},
        pattern color=black
      ]
        (-\mechSupportGroundHalfWidth,-\mechSupportGroundY)
        rectangle
        (\mechSupportGroundHalfWidth,{-(\mechSupportGroundY+\mechSupportHatchDepth)});

      \node[\pgfkeysvalueof{/mech/support-common/label position}]
        at (0,0) {\pgfkeysvalueof{/mech/support-common/label}};
    }
  },
  pics/mech/roller-support/.default={},
}
\pgfkeys{
  /mech/fixed-support/.is family,
  /mech/fixed-support,
  width/.initial=1cm,
  depth/.initial=1.5mm,
  hatch angle/.initial=45,
  angle/.initial=0,
  label/.initial=,
}
\tikzset{
  pics/mech/fixed-support/.style args={#1}{
    code={
      \pgfkeys{/mech/fixed-support/.cd,#1}
      \edef\mechFixedWidth{\pgfkeysvalueof{/mech/fixed-support/width}}
      \edef\mechFixedDepth{\pgfkeysvalueof{/mech/fixed-support/depth}}
      \edef\mechFixedHatch{\pgfkeysvalueof{/mech/fixed-support/hatch angle}}
      \edef\mechFixedAngle{\pgfkeysvalueof{/mech/fixed-support/angle}}

      \coordinate (-center) at (0,0);
      \coordinate (-connection-upper)  at (0,\mechConnectionOffset);
      \coordinate (-connection-middle) at (0,0);
      \coordinate (-connection-lower)  at (0,-\mechConnectionOffset);

      \begin{scope}[rotate=\mechFixedAngle]
        \draw[mech/body,fill=none] (-.5\mechFixedWidth,0)--(.5\mechFixedWidth,0);
        \fill[
          pattern={Lines[angle=\mechFixedHatch,distance={3pt/sqrt(2)}]},
          pattern color=black
        ]
          (-.5\mechFixedWidth,0)
          rectangle
          (.5\mechFixedWidth,\mechFixedDepth);
      \end{scope}

      \node[above] at (0,0) {\pgfkeysvalueof{/mech/fixed-support/label}};
    }
  },
  pics/mech/fixed-support/.default={},
}

% =========================================================
% mechlib/elements/connectors.tex
% ATOMARE VERBINDUNGSELEMENTE: FEDER UND DAEMPFER
% =========================================================
% Beide Elemente folgen derselben grundlegenden API: "from", "to" und "label".
% Die Endpunkte koennen beliebige TikZ-Koordinaten oder exportierte Anker aus
% anderen mechlib-Pics sein.
%
% Oeffentliche Pics:
%   mech/spring  Zickzack-Feder entlang der Verbindungslinie
%   mech/damper  Daempfer mit rechteckigem Daempferkoerper in Linienmitte
%
% Die Beschriftung wird "sloped" gezeichnet und folgt daher der Orientierung
% des Verbindungselements. Ueber "label side" kann sie ober- oder unterhalb
% der lokalen Verbindungslinie platziert werden.
% =========================================================
\pgfkeys{
  /mech/spring/.is family,
  /mech/spring,
  from/.initial={(0,0)},
  to/.initial={(1,0)},
  label/.initial=,
  label side/.initial=above,
  amplitude/.initial=2mm,
  segment length/.initial=4mm,
  style/.initial=mech/connector,
}
\tikzset{
  pics/mech/spring/.style args={#1}{
    code={
      \pgfkeys{/mech/spring/.cd,#1}
      \edef\mechSpringFrom{\pgfkeysvalueof{/mech/spring/from}}
      \edef\mechSpringTo{\pgfkeysvalueof{/mech/spring/to}}
      \draw[\pgfkeysvalueof{/mech/spring/style},
        decorate,
        decoration={zigzag,
          segment length=\pgfkeysvalueof{/mech/spring/segment length},
          amplitude=\pgfkeysvalueof{/mech/spring/amplitude}}]
        \mechSpringFrom -- \mechSpringTo
        node[midway,sloped,mech/label,outer sep=2mm,\pgfkeysvalueof{/mech/spring/label side}]
        {\pgfkeysvalueof{/mech/spring/label}};
    }
  },
  pics/mech/spring/.default={},
}
\pgfkeys{
  /mech/damper/.is family,
  /mech/damper,
  from/.initial={(0,0)},
  to/.initial={(1,0)},
  label/.initial=,
  label side/.initial=below,
  body width/.initial=8mm,
  body height/.initial=3mm,
  style/.initial=mech/connector,
}
\tikzset{
  pics/mech/damper/.style args={#1}{
    code={
      \pgfkeys{/mech/damper/.cd,#1}
      \edef\mechDamperFrom{\pgfkeysvalueof{/mech/damper/from}}
      \edef\mechDamperTo{\pgfkeysvalueof{/mech/damper/to}}
      \draw[\pgfkeysvalueof{/mech/damper/style}]
        \mechDamperFrom -- \mechDamperTo
        node[midway,sloped,draw,fill=white,
          minimum width=\pgfkeysvalueof{/mech/damper/body width},
          minimum height=\pgfkeysvalueof{/mech/damper/body height},
          inner sep=0pt] {};
      \path \mechDamperFrom -- \mechDamperTo
        node[midway,sloped,mech/label,outer sep=2mm,\pgfkeysvalueof{/mech/damper/label side}]
        {\pgfkeysvalueof{/mech/damper/label}};
    }
  },
  pics/mech/damper/.default={},
}

% =========================================================
% mechlib/elements/trusses.tex
% FACHWERKE: STAB UND GEMEINSAME GELENKDARSTELLUNG
% =========================================================
% Dieses Modul enthaelt die atomaren Fachwerkelemente. Die Topologie eines
% Fachwerks bleibt bewusst im Anwendungsdokument: Knoten werden als TikZ-
% Koordinaten definiert und durch mech/rod bzw. \mechJoints verbunden.
%
% Oeffentliche API:
%   mech/rod     Stab zwischen "from" und "to" mit optionalem Label
%   \mechJoints{A,B,C,...} zeichnet alle angegebenen Gelenke
%
% Darstellungsregeln:
%   - Staebe verwenden standardmaessig den Stil mech/rod: schwarz und duenn.
%   - Staebe werden auf dem Layer "mech background" gezeichnet.
%   - Stabnummern besitzen keine Fuellung (fill=none) und bleiben transparent.
%   - Gelenke liegen optisch vor den Staeben und sind mit mechBodyColor gefuellt.
%
% Beispiel:
%   \pic {mech/rod={from=(A),to=(B),label=1,label side=above}};
%   \mechJoints{A,B}
% =========================================================
% -------------------------------------------------------------------------
% Einzelner Fachwerkstab
% -------------------------------------------------------------------------
% "style" erlaubt lokale Abweichungen, ohne den globalen mech/rod-Stil zu
% veraendern. Die Standardbeschriftung liegt an der Stabmitte.
\pgfkeys{
  /mech/rod/.is family,
  /mech/rod,
  from/.initial={(0,0)},
  to/.initial={(1,0)},
  label/.initial=,
  label side/.initial=above,
  label style/.initial={mech/label,fill=none,text=black},
  style/.initial=mech/rod,
}
\tikzset{
  pics/mech/rod/.style args={#1}{
    code={
      \pgfkeys{/mech/rod/.cd,#1}
      \edef\mechRodFrom{\pgfkeysvalueof{/mech/rod/from}}
      \edef\mechRodTo{\pgfkeysvalueof{/mech/rod/to}}
      \coordinate (-start) at \mechRodFrom;
      \coordinate (-end) at \mechRodTo;
      \coordinate (-center) at ($(-start)!.5!(-end)$);
      \begin{pgfonlayer}{mech background}
        \draw[\pgfkeysvalueof{/mech/rod/style}] (-start)--(-end);
      \end{pgfonlayer}
      \node[\pgfkeysvalueof{/mech/rod/label style},\pgfkeysvalueof{/mech/rod/label side}]
        at (-center) {\pgfkeysvalueof{/mech/rod/label}};
    }
  },
  pics/mech/rod/.default={},
}
% -------------------------------------------------------------------------
% Gelenkliste
% -------------------------------------------------------------------------
% Das Makro erwartet eine kommaseparierte Liste bereits definierter TikZ-
% Koordinaten. Jedes Gelenk wird identisch gezeichnet. Die Gelenke sollten nach
% den Staeben aufgerufen werden, damit sie die Stabenden optisch ueberdecken.
\newcommand{\mechJoints}[1]{%
  \foreach \mechJoint in {#1}{%
    \draw[fill=mechBodyColor,draw=black,line width=.2mm]
      (\mechJoint) circle[radius=.6mm];%
  }%
}

% -------------------------------------------------------------------------
% Annotationen und mechanische Symbolgruppen.
% -------------------------------------------------------------------------
% =========================================================
% mechlib/annotations/loads.tex
% LASTANNOTATIONEN: KRAFT, MOMENT UND STRECKENLAST
% =========================================================
% Diese Datei stellt die Lastsymbole der mechlib bereit. Alle Lasten verwenden
% standardmaessig den zentralen Stil mech/load und damit mechLoadColor.
%
% Oeffentliche Pics:
%   mech/force            Einzelkraft als Winkel/Laenge oder from/to
%   mech/moment           Kreisbogensymbol im oder gegen den Uhrzeigersinn
%   mech/distributed-load konstante, dreieckige oder linear veraenderliche Last
%
% Besonderheit der Streckenlast:
% Die Grundlinie wird lokal auf die x-Achse gedreht. Dadurch kann die gesamte
% Geometrie in einem einfachen lokalen Koordinatensystem berechnet werden. Sehr
% kurze Pfeile, deren Laenge kleiner als ihre Pfeilspitze ist, werden bewusst
% unterdrueckt; so entstehen am Nullende einer Dreieckslast keine Artefakte.
% =========================================================
\pgfkeys{
  /mech/force/.is family,
  /mech/force,
  from/.initial={(0,0)},
  to/.initial=,
  angle/.initial=0,
  length/.initial=1.5cm,
  label/.initial=,
  label position/.initial=above,
  direction/.initial=away,
  style/.initial=mech/load,
}
\tikzset{
  pics/mech/force/.style args={#1}{
    code={
      \pgfkeys{/mech/force/.cd,#1}
      \edef\mechForceFrom{\pgfkeysvalueof{/mech/force/from}}
      \edef\mechForceTo{\pgfkeysvalueof{/mech/force/to}}
      \edef\mechForceDirection{\pgfkeysvalueof{/mech/force/direction}}
      \def\mechForceToward{toward}
      \ifx\mechForceTo\empty
        \ifx\mechForceDirection\mechForceToward
          \draw[\pgfkeysvalueof{/mech/force/style},-{Latex[length=2.5mm]}]
            (\pgfkeysvalueof{/mech/force/angle}:\pgfkeysvalueof{/mech/force/length}) -- (0,0)
            node[midway,\pgfkeysvalueof{/mech/force/label position}]
            {\pgfkeysvalueof{/mech/force/label}};
        \else
          \draw[\pgfkeysvalueof{/mech/force/style},-{Latex[length=2.5mm]}]
            (0,0) -- (\pgfkeysvalueof{/mech/force/angle}:\pgfkeysvalueof{/mech/force/length})
            node[midway,\pgfkeysvalueof{/mech/force/label position}]
            {\pgfkeysvalueof{/mech/force/label}};
        \fi
      \else
        \ifx\mechForceDirection\mechForceToward
          \draw[\pgfkeysvalueof{/mech/force/style},-{Latex[length=2.5mm]}]
            \mechForceTo -- \mechForceFrom
            node[midway,\pgfkeysvalueof{/mech/force/label position}]
            {\pgfkeysvalueof{/mech/force/label}};
        \else
          \draw[\pgfkeysvalueof{/mech/force/style},-{Latex[length=2.5mm]}]
            \mechForceFrom -- \mechForceTo
            node[midway,\pgfkeysvalueof{/mech/force/label position}]
            {\pgfkeysvalueof{/mech/force/label}};
        \fi
      \fi
    }
  },
  pics/mech/force/.default={},
}
\pgfkeys{
  /mech/moment/.is family,
  /mech/moment,
  radius/.initial=6mm,
  label/.initial=,
  direction/.is choice,
  direction/ccw/.code={\def\mechMomentStart{-60}\def\mechMomentEnd{220}},
  direction/cw/.code={\def\mechMomentStart{240}\def\mechMomentEnd{-40}},
  direction=ccw,
  style/.initial=mech/load,
}
\tikzset{
  pics/mech/moment/.style args={#1}{
    code={
      \pgfkeys{/mech/moment/.cd,direction=ccw,#1}
      \draw[\pgfkeysvalueof{/mech/moment/style},-{Latex[length=2.5mm]}]
        (\mechMomentStart:\pgfkeysvalueof{/mech/moment/radius})
        arc[start angle=\mechMomentStart,end angle=\mechMomentEnd,
            radius=\pgfkeysvalueof{/mech/moment/radius}];
      \node[mech/label,above] at (0,\pgfkeysvalueof{/mech/moment/radius})
        {\pgfkeysvalueof{/mech/moment/label}};
    }
  },
  pics/mech/moment/.default={},
}
\pgfkeys{
  /mech/distributed-load/.is family,
  /mech/distributed-load,
  from/.initial={(0,0)},
  to/.initial={(1,0)},
  start value/.initial=1,
  end value/.initial=1,
  height/.initial=8mm,
  offset/.initial=0mm,
  arrows/.initial=6,
  arrowhead length/.initial=2mm,
  label/.initial=,
  label position/.initial=above,
  style/.initial=mech/load,
}
\tikzset{
  pics/mech/distributed-load/.style args={#1}{
    code={
      \pgfkeys{/mech/distributed-load/.cd,#1}
      \edef\mechDLFrom{\pgfkeysvalueof{/mech/distributed-load/from}}
      \edef\mechDLTo{\pgfkeysvalueof{/mech/distributed-load/to}}
      \path let
        \p1=\mechDLFrom,\p2=\mechDLTo,
        \n1={veclen(\x2-\x1,\y2-\y1)},
        \n2={atan2(\y2-\y1,\x2-\x1)}
      in
        coordinate (-start) at (\x1,\y1)
        coordinate (-end) at (\x2,\y2)
        coordinate (-helper) at (\n2,\n1);
      \path (-helper); \pgfgetlastxy{\mechDLAngle}{\mechDLLength};
      \pgfmathsetmacro{\mechDLStartValue}{\pgfkeysvalueof{/mech/distributed-load/start value}}
      \pgfmathsetmacro{\mechDLEndValue}{\pgfkeysvalueof{/mech/distributed-load/end value}}
      \pgfmathsetmacro{\mechDLArrowCount}{\pgfkeysvalueof{/mech/distributed-load/arrows}}
      \begin{scope}[shift={(-start)},rotate=\mechDLAngle]
        \pgfmathsetlengthmacro{\mechDLOffset}{\pgfkeysvalueof{/mech/distributed-load/offset}}
        \pgfmathsetlengthmacro{\mechDLArrowheadLength}{\pgfkeysvalueof{/mech/distributed-load/arrowhead length}}
        \pgfmathsetlengthmacro{\mechDLStartArrowLength}{abs(\mechDLStartValue)*\pgfkeysvalueof{/mech/distributed-load/height}}

        \fill[mech/load fill]
          (0,\mechDLOffset)
          --(0,{\mechDLOffset+\mechDLStartValue*\pgfkeysvalueof{/mech/distributed-load/height}})
          --(\mechDLLength,{\mechDLOffset+\mechDLEndValue*\pgfkeysvalueof{/mech/distributed-load/height}})
          --(\mechDLLength,\mechDLOffset)
          --cycle;

        % Keep contour and first arrow continuous whenever the first arrow
        % is long enough to accommodate its arrowhead. Otherwise draw only
        % the contour and suppress the too-short arrow.
        \ifdim\mechDLStartArrowLength<\mechDLArrowheadLength
          \draw[\pgfkeysvalueof{/mech/distributed-load/style}]
            (\mechDLLength,{\mechDLOffset+\mechDLEndValue*\pgfkeysvalueof{/mech/distributed-load/height}})
            --(0,{\mechDLOffset+\mechDLStartValue*\pgfkeysvalueof{/mech/distributed-load/height}});
        \else
          \draw[
            \pgfkeysvalueof{/mech/distributed-load/style},
            -{Latex[length=\pgfkeysvalueof{/mech/distributed-load/arrowhead length}]}
          ]
            (\mechDLLength,{\mechDLOffset+\mechDLEndValue*\pgfkeysvalueof{/mech/distributed-load/height}})
            --(0,{\mechDLOffset+\mechDLStartValue*\pgfkeysvalueof{/mech/distributed-load/height}})
            --(0,\mechDLOffset);
        \fi

        \foreach \mechDLIndex in {1,...,\mechDLArrowCount}{
          \pgfmathsetmacro{\mechDLT}{\mechDLIndex/\mechDLArrowCount}
          \pgfmathsetmacro{\mechDLValue}{\mechDLStartValue+(\mechDLEndValue-\mechDLStartValue)*\mechDLT}
          \pgfmathsetlengthmacro{\mechDLArrowLength}{abs(\mechDLValue)*\pgfkeysvalueof{/mech/distributed-load/height}}

          \ifdim\mechDLArrowLength<\mechDLArrowheadLength
            % Suppress arrows that are shorter than their arrowhead.
          \else
            \draw[
              \pgfkeysvalueof{/mech/distributed-load/style},
              -{Latex[length=\pgfkeysvalueof{/mech/distributed-load/arrowhead length}]}
            ]
              ({\mechDLT*\mechDLLength},{\mechDLOffset+\mechDLValue*\pgfkeysvalueof{/mech/distributed-load/height}})
              --({\mechDLT*\mechDLLength},\mechDLOffset);
          \fi
        }

        \node[mech/label,\pgfkeysvalueof{/mech/distributed-load/label position}]
          at (\mechDLLength,{\mechDLOffset+\mechDLEndValue*\pgfkeysvalueof{/mech/distributed-load/height}})
          {\pgfkeysvalueof{/mech/distributed-load/label}};
      \end{scope}
    }
  },
  pics/mech/distributed-load/.default={},
}

% =========================================================
% mechlib/annotations/dimensions.tex
% BEMASSUNGEN: LINEARES MASS UND WINKELMASS
% =========================================================
% Lineare Bemassungen werden geometrisch aus zwei Endpunkten abgeleitet. Der
% Offset wirkt senkrecht auf die Verbindung von "from" nach "to"; damit
% funktioniert dasselbe Pic fuer horizontale, vertikale und schraege Masse.
%
% Oeffentliche Pics:
%   mech/dimension       lineare Bemassung mit Doppelpfeil
%   mech/angle-dimension Kreisbogen zwischen zwei Winkeln
%
% Interne Koordinaten -A und -B sind die senkrecht versetzten Endpunkte der
% Masslinie. Die Vektorrechnung erfolgt in TikZ "let"-Syntax.
% =========================================================
\pgfkeys{
  /mech/dimension/.is family,
  /mech/dimension,
  from/.initial={(0,0)},
  to/.initial={(1,0)},
  label/.initial={$l$},
  offset/.initial=3mm,
  style/.initial=mech/dimension,
}
\tikzset{
  pics/mech/dimension/.style args={#1}{
    code={
      \pgfkeys{/mech/dimension/.cd,#1}
      \edef\mechDimFrom{\pgfkeysvalueof{/mech/dimension/from}}
      \edef\mechDimTo{\pgfkeysvalueof{/mech/dimension/to}}
      \path let
        \p1=\mechDimFrom,\p2=\mechDimTo,
        \n1={veclen(\x2-\x1,\y2-\y1)},
        \n2={-(\y2-\y1)/\n1*\pgfkeysvalueof{/mech/dimension/offset}},
        \n3={(\x2-\x1)/\n1*\pgfkeysvalueof{/mech/dimension/offset}}
      in
        coordinate (-A) at (\x1+\n2,\y1+\n3)
        coordinate (-B) at (\x2+\n2,\y2+\n3);
      \draw[\pgfkeysvalueof{/mech/dimension/style}]
        (-A)--(-B)
        node[midway,fill=white,inner sep=1pt]
        {\pgfkeysvalueof{/mech/dimension/label}};
    }
  },
  pics/mech/dimension/.default={},
}
\pgfkeys{
  /mech/angle-dimension/.is family,
  /mech/angle-dimension,
  start angle/.initial=0,
  end angle/.initial=45,
  radius/.initial=7mm,
  label/.initial={$\alpha$},
}
\tikzset{
  pics/mech/angle-dimension/.style args={#1}{
    code={
      \pgfkeys{/mech/angle-dimension/.cd,#1}
      \draw
        (\pgfkeysvalueof{/mech/angle-dimension/start angle}:\pgfkeysvalueof{/mech/angle-dimension/radius})
        arc[start angle=\pgfkeysvalueof{/mech/angle-dimension/start angle},
            end angle=\pgfkeysvalueof{/mech/angle-dimension/end angle},
            radius=\pgfkeysvalueof{/mech/angle-dimension/radius}];
      \pgfmathsetmacro{\mechAngleMid}{.5*(\pgfkeysvalueof{/mech/angle-dimension/start angle}+\pgfkeysvalueof{/mech/angle-dimension/end angle})}
      \node at (\mechAngleMid:{.7*\pgfkeysvalueof{/mech/angle-dimension/radius}})
        {\pgfkeysvalueof{/mech/angle-dimension/label}};
    }
  },
  pics/mech/angle-dimension/.default={},
}

% =========================================================
% mechlib/annotations/coordinates.tex
% KOORDINATENANNOTATIONEN: EINZELACHSE UND 2D-KOORDINATENSYSTEM
% =========================================================
% Die Pics dieser Datei zeichnen mathematische Bezugsachsen im Stil mech/axis.
% Winkel werden in TikZ-Grad angegeben; Laengen sind regulaere TeX-Dimensionen.
%
% Oeffentliche Pics:
%   mech/axis              einzelne gerichtete Achse
%   mech/coordinate-system zwei unabhaengig orientierbare Achsen
%
% Beide Pics exportieren den lokalen Ursprung als "-origin". Bei einem benannten
% Pic (cs) kann dieser z. B. als (cs-origin) referenziert werden.
% =========================================================
\pgfkeys{
  /mech/axis/.is family,
  /mech/axis,
  label/.initial={$x$},
  angle/.initial=0,
  length/.initial=1cm,
  direction/.initial=positive,
}
\tikzset{
  pics/mech/axis/.style args={#1}{
    code={
      \pgfkeys{/mech/axis/.cd,#1}
      \edef\mechAxisDirection{\pgfkeysvalueof{/mech/axis/direction}}
      \def\mechAxisNegative{negative}
      \ifx\mechAxisDirection\mechAxisNegative
        \draw[mech/axis,{Latex[length=2mm]}-]
          (0,0)--(\pgfkeysvalueof{/mech/axis/angle}:\pgfkeysvalueof{/mech/axis/length})
          node[pos=1,anchor=west]{\pgfkeysvalueof{/mech/axis/label}};
      \else
        \draw[mech/axis,-{Latex[length=2mm]}]
          (0,0)--(\pgfkeysvalueof{/mech/axis/angle}:\pgfkeysvalueof{/mech/axis/length})
          node[pos=1,anchor=west]{\pgfkeysvalueof{/mech/axis/label}};
      \fi
      \coordinate (-origin) at (0,0);
    }
  },
  pics/mech/axis/.default={},
}
\pgfkeys{
  /mech/coordinate-system/.is family,
  /mech/coordinate-system,
  x label/.initial={$x$},
  y label/.initial={$y$},
  x angle/.initial=0,
  y angle/.initial=90,
  x length/.initial=1cm,
  y length/.initial=1cm,
}
\tikzset{
  pics/mech/coordinate-system/.style args={#1}{
    code={
      \pgfkeys{/mech/coordinate-system/.cd,#1}
      \draw[mech/axis,-{Latex[length=2mm]}]
        (0,0)--(\pgfkeysvalueof{/mech/coordinate-system/x angle}:\pgfkeysvalueof{/mech/coordinate-system/x length})
        node[pos=1,anchor=west]{\pgfkeysvalueof{/mech/coordinate-system/x label}};
      \draw[mech/axis,-{Latex[length=2mm]}]
        (0,0)--(\pgfkeysvalueof{/mech/coordinate-system/y angle}:\pgfkeysvalueof{/mech/coordinate-system/y length})
        node[pos=1,anchor=east]{\pgfkeysvalueof{/mech/coordinate-system/y label}};
      \coordinate (-origin) at (0,0);
    }
  },
  pics/mech/coordinate-system/.default={},
}

% =========================================================
% mechlib/annotations/mechanics.tex
% MECHANISCHE ANNOTATIONEN: SCHNITTGROESSEN UND LAGERREAKTIONEN
% =========================================================
% Die Pics kombinieren mehrere elementare Lastsymbole zu haeufig benoetigten
% mechanischen Symbolgruppen. Farben und Pfeilgestaltung folgen mech/load.
%
% Oeffentliche Pics:
%   mech/internal-forces  Normalkraft N, Querkraft Q und Moment M am Schnitt
%   mech/fixed-reactions  horizontale/vertikale Reaktion und Einspannmoment
%
% Bei mech/internal-forces steuert "side=right|left" die Vorzeichenrichtung
% aller drei Schnittgroessen konsistent. Das Moment wird dazu intern als
% mech/moment mit passender Drehrichtung wiederverwendet.
% =========================================================
\pgfkeys{
  /mech/internal-forces/.is family,
  /mech/internal-forces,
  normal/.initial={$N$},
  shear/.initial={$Q$},
  moment/.initial={$M$},
  length/.initial=7mm,
  side/.is choice,
  side/right/.code={\def\mechIFSign{1}},
  side/left/.code={\def\mechIFSign{-1}},
  side=right,
}
\tikzset{
  pics/mech/internal-forces/.style args={#1}{
    code={
      \pgfkeys{/mech/internal-forces/.cd,side=right,#1}
      \edef\mechIFLength{\pgfkeysvalueof{/mech/internal-forces/length}}
      \draw[mech/load,-{Latex[length=2mm]}]
        (0,0)--({\mechIFSign*\mechIFLength},0)
        node[pos=1,anchor=west]{\pgfkeysvalueof{/mech/internal-forces/normal}};
      \draw[mech/load,-{Latex[length=2mm]}]
        (0,0)--(0,{-.8*\mechIFSign*\mechIFLength})
        node[pos=1,anchor=west]{\pgfkeysvalueof{/mech/internal-forces/shear}};
      \ifnum\mechIFSign=1
        \pic {mech/moment={direction=ccw,radius=.55\mechIFLength,label=\pgfkeysvalueof{/mech/internal-forces/moment}}};
      \else
        \pic {mech/moment={direction=cw,radius=.55\mechIFLength,label=\pgfkeysvalueof{/mech/internal-forces/moment}}};
      \fi
    }
  },
  pics/mech/internal-forces/.default={},
}
\pgfkeys{
  /mech/fixed-reactions/.is family,
  /mech/fixed-reactions,
  x/.initial={$A_x$},
  y/.initial={$A_y$},
  moment/.initial={$M_A$},
  length/.initial=7mm,
}
\tikzset{
  pics/mech/fixed-reactions/.style args={#1}{
    code={
      \pgfkeys{/mech/fixed-reactions/.cd,#1}
      \edef\mechReactionLength{\pgfkeysvalueof{/mech/fixed-reactions/length}}
      \draw[mech/load,-{Latex[length=2mm]}] (-\mechReactionLength,0)--(0,0)
        node[pos=0,anchor=east]{\pgfkeysvalueof{/mech/fixed-reactions/x}};
      \draw[mech/load,-{Latex[length=2mm]}] (0,.7\mechReactionLength)--(0,-.3\mechReactionLength)
        node[pos=0,anchor=south]{\pgfkeysvalueof{/mech/fixed-reactions/y}};
      \pic {mech/moment={direction=ccw,radius=.45\mechReactionLength,label=\pgfkeysvalueof{/mech/fixed-reactions/moment}}};
    }
  },
  pics/mech/fixed-reactions/.default={},
}


%
% Liegt die Bibliothek an einem anderen Ort, kann vor dem Laden die Variable
% \mechlibpath gesetzt werden:
%   \newcommand{\mechlibpath}{../../../support/mechlib}
%   % =========================================================
% mechlib/mechlib.tex
% ZENTRALER LOADER / OEFFENTLICHER EINSTIEGSPUNKT
% =========================================================
% Diese Datei ist die einzige Datei, die ein Anwender normalerweise direkt
% einbindet. Sie laedt die benoetigten LaTeX-/TikZ-Abhaengigkeiten und danach
% alle Teilmodule der Bibliothek in einer definierten Reihenfolge.
%
% Typische Verwendung bei einem Ordner "mechlib" neben dem Hauptdokument:
%   % =========================================================
% mechlib/mechlib.tex
% ZENTRALER LOADER / OEFFENTLICHER EINSTIEGSPUNKT
% =========================================================
% Diese Datei ist die einzige Datei, die ein Anwender normalerweise direkt
% einbindet. Sie laedt die benoetigten LaTeX-/TikZ-Abhaengigkeiten und danach
% alle Teilmodule der Bibliothek in einer definierten Reihenfolge.
%
% Typische Verwendung bei einem Ordner "mechlib" neben dem Hauptdokument:
%   \input{mechlib/mechlib}
%
% Liegt die Bibliothek an einem anderen Ort, kann vor dem Laden die Variable
% \mechlibpath gesetzt werden:
%   \newcommand{\mechlibpath}{../../../support/mechlib}
%   \input{\mechlibpath/mechlib}
%
% Lade-Reihenfolge:
%   1. core/        Farben, Layer, Stile und globale Geometriekonventionen
%   2. elements/    mechanische Grundelemente (Koerper, Lager, Verbinder, Staebe)
%   3. annotations/ Lasten, Masse, Koordinaten und mechanische Annotationen
%
% WICHTIG:
% Die Modulreihenfolge ist absichtlich festgelegt. Elemente greifen auf Farben,
% Layer und Stile aus core/ zurueck; Annotationen koennen wiederum die bereits
% definierten Grundstile und Pics verwenden.
% =========================================================
% -------------------------------------------------------------------------
% Externe Pakete und TikZ-Bibliotheken
% -------------------------------------------------------------------------
% \RequirePackage ist bewusst statt \usepackage gewaehlt, da mechlib.tex selbst
% wie eine Bibliotheks-/Paketdatei verwendet wird.
\RequirePackage{tikz}
\RequirePackage{amsmath}
\RequirePackage{siunitx}
\RequirePackage{xcolor}
\usetikzlibrary{
  arrows.meta,
  calc,
  patterns,
  patterns.meta,
  intersections,
  angles,
  quotes,
  decorations.pathmorphing,
  backgrounds
}
% Base path: override before loading when mechlib is stored elsewhere.
% -------------------------------------------------------------------------
% Aufloesung des Bibliothekspfades
% -------------------------------------------------------------------------
% \providecommand respektiert eine bereits im Hauptdokument gesetzte Definition.
\providecommand{\mechlibpath}{mechlib}
% -------------------------------------------------------------------------
% Kernmodule: muessen vor allen Elementen geladen werden.
% -------------------------------------------------------------------------
\input{\mechlibpath/core/colors}
\input{\mechlibpath/core/layers}
\input{\mechlibpath/core/styles}
\input{\mechlibpath/core/defaults}
% -------------------------------------------------------------------------
% Atomare mechanische Elemente.
% -------------------------------------------------------------------------
\input{\mechlibpath/elements/bodies}
\input{\mechlibpath/elements/supports}
\input{\mechlibpath/elements/connectors}
\input{\mechlibpath/elements/trusses}
% -------------------------------------------------------------------------
% Annotationen und mechanische Symbolgruppen.
% -------------------------------------------------------------------------
\input{\mechlibpath/annotations/loads}
\input{\mechlibpath/annotations/dimensions}
\input{\mechlibpath/annotations/coordinates}
\input{\mechlibpath/annotations/mechanics}

%
% Liegt die Bibliothek an einem anderen Ort, kann vor dem Laden die Variable
% \mechlibpath gesetzt werden:
%   \newcommand{\mechlibpath}{../../../support/mechlib}
%   % =========================================================
% mechlib/mechlib.tex
% ZENTRALER LOADER / OEFFENTLICHER EINSTIEGSPUNKT
% =========================================================
% Diese Datei ist die einzige Datei, die ein Anwender normalerweise direkt
% einbindet. Sie laedt die benoetigten LaTeX-/TikZ-Abhaengigkeiten und danach
% alle Teilmodule der Bibliothek in einer definierten Reihenfolge.
%
% Typische Verwendung bei einem Ordner "mechlib" neben dem Hauptdokument:
%   \input{mechlib/mechlib}
%
% Liegt die Bibliothek an einem anderen Ort, kann vor dem Laden die Variable
% \mechlibpath gesetzt werden:
%   \newcommand{\mechlibpath}{../../../support/mechlib}
%   \input{\mechlibpath/mechlib}
%
% Lade-Reihenfolge:
%   1. core/        Farben, Layer, Stile und globale Geometriekonventionen
%   2. elements/    mechanische Grundelemente (Koerper, Lager, Verbinder, Staebe)
%   3. annotations/ Lasten, Masse, Koordinaten und mechanische Annotationen
%
% WICHTIG:
% Die Modulreihenfolge ist absichtlich festgelegt. Elemente greifen auf Farben,
% Layer und Stile aus core/ zurueck; Annotationen koennen wiederum die bereits
% definierten Grundstile und Pics verwenden.
% =========================================================
% -------------------------------------------------------------------------
% Externe Pakete und TikZ-Bibliotheken
% -------------------------------------------------------------------------
% \RequirePackage ist bewusst statt \usepackage gewaehlt, da mechlib.tex selbst
% wie eine Bibliotheks-/Paketdatei verwendet wird.
\RequirePackage{tikz}
\RequirePackage{amsmath}
\RequirePackage{siunitx}
\RequirePackage{xcolor}
\usetikzlibrary{
  arrows.meta,
  calc,
  patterns,
  patterns.meta,
  intersections,
  angles,
  quotes,
  decorations.pathmorphing,
  backgrounds
}
% Base path: override before loading when mechlib is stored elsewhere.
% -------------------------------------------------------------------------
% Aufloesung des Bibliothekspfades
% -------------------------------------------------------------------------
% \providecommand respektiert eine bereits im Hauptdokument gesetzte Definition.
\providecommand{\mechlibpath}{mechlib}
% -------------------------------------------------------------------------
% Kernmodule: muessen vor allen Elementen geladen werden.
% -------------------------------------------------------------------------
\input{\mechlibpath/core/colors}
\input{\mechlibpath/core/layers}
\input{\mechlibpath/core/styles}
\input{\mechlibpath/core/defaults}
% -------------------------------------------------------------------------
% Atomare mechanische Elemente.
% -------------------------------------------------------------------------
\input{\mechlibpath/elements/bodies}
\input{\mechlibpath/elements/supports}
\input{\mechlibpath/elements/connectors}
\input{\mechlibpath/elements/trusses}
% -------------------------------------------------------------------------
% Annotationen und mechanische Symbolgruppen.
% -------------------------------------------------------------------------
\input{\mechlibpath/annotations/loads}
\input{\mechlibpath/annotations/dimensions}
\input{\mechlibpath/annotations/coordinates}
\input{\mechlibpath/annotations/mechanics}

%
% Lade-Reihenfolge:
%   1. core/        Farben, Layer, Stile und globale Geometriekonventionen
%   2. elements/    mechanische Grundelemente (Koerper, Lager, Verbinder, Staebe)
%   3. annotations/ Lasten, Masse, Koordinaten und mechanische Annotationen
%
% WICHTIG:
% Die Modulreihenfolge ist absichtlich festgelegt. Elemente greifen auf Farben,
% Layer und Stile aus core/ zurueck; Annotationen koennen wiederum die bereits
% definierten Grundstile und Pics verwenden.
% =========================================================
% -------------------------------------------------------------------------
% Externe Pakete und TikZ-Bibliotheken
% -------------------------------------------------------------------------
% \RequirePackage ist bewusst statt \usepackage gewaehlt, da mechlib.tex selbst
% wie eine Bibliotheks-/Paketdatei verwendet wird.
\RequirePackage{tikz}
\RequirePackage{amsmath}
\RequirePackage{siunitx}
\RequirePackage{xcolor}
\usetikzlibrary{
  arrows.meta,
  calc,
  patterns,
  patterns.meta,
  intersections,
  angles,
  quotes,
  decorations.pathmorphing,
  backgrounds
}
% Base path: override before loading when mechlib is stored elsewhere.
% -------------------------------------------------------------------------
% Aufloesung des Bibliothekspfades
% -------------------------------------------------------------------------
% \providecommand respektiert eine bereits im Hauptdokument gesetzte Definition.
\providecommand{\mechlibpath}{mechlib}
% -------------------------------------------------------------------------
% Kernmodule: muessen vor allen Elementen geladen werden.
% -------------------------------------------------------------------------
% =========================================================
% mechlib/core/colors.tex
% ZENTRALE FARBPALETTE
% =========================================================
% Alle bibliotheksweiten Farben werden ausschliesslich hier definiert. Dadurch
% kann der visuelle Charakter der gesamten mechlib an einer Stelle angepasst
% werden, ohne einzelne Elemente zu veraendern.
%
% Konvention:
%   mechBodyColor    Fuellung allgemeiner Koerper und Fachwerkgelenke
%   mechSupportColor Fuellung von Lagerkoerpern
%   mechLoadColor    Kraefte, Momente und Streckenlasten
%   mechAxisColor    Achsen und Koordinatensysteme
%   mechRodColor     historische Stabfarbe; aktuelle Fachwerkstaebe sind schwarz
%
% Die Namen sind Teil der internen visuellen API und sollten bevorzugt statt
% harter RGB-Werte in den Elementdateien verwendet werden.
% =========================================================
% -------------------------------------------------------------------------
% Farbdefinitionen
% -------------------------------------------------------------------------
% RGB wird bewusst explizit verwendet, damit die Darstellung unabhaengig von
% Dokumentklassen und externen Farbmodellen reproduzierbar bleibt.
\definecolor{mechBodyColor}{RGB}{180,180,180}
\definecolor{mechSupportColor}{RGB}{205,205,205}
\definecolor{mechLoadColor}{RGB}{178,25,25}
\definecolor{mechAxisColor}{RGB}{6,13,141}
\definecolor{mechRodColor}{RGB}{130,130,130}

% =========================================================
% mechlib/core/layers.tex
% ZEICHENEBENEN / PGF-LAYER
% =========================================================
% Fachwerkstaebe und andere hinterlegte Geometrien koennen auf dem Layer
% "mech background" gezeichnet werden. Der normale TikZ-Layer "main" liegt
% darueber. Dadurch koennen z. B. gefuellte Gelenke sauber vor den Staeben
% erscheinen, ohne dass die Zeichenreihenfolge im Anwendungsdokument aufwendig
% kontrolliert werden muss.
% =========================================================
% Zuerst wird der benannte Hintergrundlayer deklariert. Anschliessend wird die
% globale Layerreihenfolge explizit gesetzt: Hintergrund vor Hauptlayer.
\pgfdeclarelayer{mech background}
\pgfsetlayers{mech background,main}

% =========================================================
% mechlib/core/styles.tex
% ZENTRALE, NAMENSRAUMGESCHUETZTE TIKZ-STILE
% =========================================================
% Diese Datei definiert ausschliesslich visuelle Grundstile. Die eigentliche
% Geometrie eines Elements gehoert in elements/ bzw. annotations/.
%
% Alle oeffentlichen Bibliotheksstile beginnen mit "mech/". Das reduziert
% Namenskollisionen mit dem Hauptdokument und macht im TikZ-Code sichtbar,
% welche Formatierung aus der mechlib stammt.
%
% Designregel:
%   - Farben werden aus core/colors.tex bezogen.
%   - Geometrien werden NICHT hier definiert.
%   - Linienstarken werden hier zentral gehalten, soweit sie elementuebergreifend
%     konsistent sein sollen.
% =========================================================
% -------------------------------------------------------------------------
% Visuelle Basisstile
% -------------------------------------------------------------------------
% Einzelne Pics duerfen diese Stile ueber ihre jeweiligen "style"-Keys ersetzen.
\tikzset{
  mech/body/.style={
    draw=black,
    fill=mechBodyColor,
    line width=.6pt
  },
  mech/support/.style={
    draw=black,
    fill=mechSupportColor,
    line width=.6pt
  },
  mech/connector/.style={
    draw=black,
    line width=.6pt
  },
  mech/rod/.style={
    draw=black,
    line width=.85pt
  },
  mech/load/.style={
    draw=mechLoadColor,
    text=mechLoadColor,
    line width=1pt
  },
  mech/load fill/.style={
    fill=mechLoadColor,
    fill opacity=.16,
    draw=none
  },
  mech/dimension/.style={
    {Latex[length=2mm]}-{Latex[length=2mm]},
    line width=.5pt
  },
  mech/axis/.style={
    draw=mechAxisColor,
    text=mechAxisColor,
    line width=.6pt
  },
  mech/label/.style={
    inner sep=1pt
  }
}

% =========================================================
% mechlib/core/defaults.tex
% GLOBALE STANDARDWERTE UND GEMEINSAME GEOMETRIEKONVENTIONEN
% =========================================================
% Hier stehen Werte, die von mehreren Modulen gemeinsam benutzt werden. Solche
% Konstanten verhindern, dass geometrisch zusammengehoerige Elemente lokal mit
% leicht unterschiedlichen Zahlen definiert werden.
%
% Beispiel: \mechConnectionOffset legt die vertikale Lage paralleler
% Anschlussanker fest. Masse und Einspannung verwenden exakt denselben Wert;
% dadurch liegen obere/mittlere/untere Anschluesse bei gleicher Mittelpunktlage
% garantiert auf gleicher Hoehe.
% =========================================================
% Globale Schriftvorgabe fuer jedes TikZ-Bild. "append style" ergaenzt
% vorhandene Bildstile, statt sie zu ueberschreiben.
\tikzset{
  every picture/.append style={font=\fontsize{10}{10}\selectfont},
}

% Standardwinkel fuer technische Bodenschraffuren.
\def\mechHatchAngle{45}

% Shared vertical distance of parallel connector anchors from the middle anchor.
% Wall and mass use this exact value so upper/middle/lower are aligned whenever
% their center points lie on the same horizontal line.
% Gemeinsamer vertikaler Versatz fuer parallele Anschlussanker.
\def\mechConnectionOffset{2.5mm}

% -------------------------------------------------------------------------
% Atomare mechanische Elemente.
% -------------------------------------------------------------------------
% =========================================================
% mechlib/elements/bodies.tex
% ATOMARE KOERPER: MASSE, PUNKTMASSE, WAGEN UND BALKEN
% =========================================================
% Die hier definierten TikZ-Pics stellen mechanische Grundkoerper bereit. Jedes
% Pic besitzt eine eigene pgfkeys-Familie unter /mech/... und kann dadurch mit
% benannten Optionen konfiguriert werden.
%
% Grundprinzip der API:
%   \pic (name) at (<Ort>) {mech/<element>={<key>=<wert>,...}};
%
% Benannte Pics exportieren Koordinaten bzw. werden bei node-basierten Koerpern
% auf den internen Shape-Knoten aliasiert. So koennen regulaere TikZ-Anker wie
% (m1.east), (m1.north west) usw. direkt benutzt werden.
%
% Oeffentliche Pics in dieser Datei:
%   mech/mass       rechteckiger Starrkoerper bzw. konzentrierte Masse
%   mech/point-mass kreisfoermige Punktmasse
%   mech/cart       Wagen mit zwei Raedern und optionalem Boden
%   mech/beam       Balken zwischen zwei beliebigen Punkten
% =========================================================
\pgfkeys{
  /mech/mass/.is family,
  /mech/mass,
  width/.initial=1.5cm,
  height/.initial=1cm,
  label/.initial=,
  style/.initial=mech/body,
}
\tikzset{
  pics/mech/mass/.style args={#1}{
    code={
      \pgfkeys{/mech/mass/.cd,#1}
      \node[
        \pgfkeysvalueof{/mech/mass/style},
        minimum width=\pgfkeysvalueof{/mech/mass/width},
        minimum height=\pgfkeysvalueof{/mech/mass/height},
        inner sep=0pt
      ] (-shape) at (0,0) {\pgfkeysvalueof{/mech/mass/label}};

      % Named body pics behave like regular TikZ nodes.
      % Example: \pic (m1) ... then (m1.east), (m1.south east), ...
      \edef\mechPicName{\pgfkeysvalueof{/tikz/name prefix}}
      \ifx\mechPicName\pgfutil@empty\else
        \pgfnodealias{\mechPicName}{\mechPicName-shape}
      \fi

      \coordinate (-center) at (-shape.center);

      % Three standardized connection levels on each vertical side.
      % The shared \mechConnectionOffset is also used by fixed supports.
      \coordinate (-west-upper)  at ($(-shape.west)+(0,\mechConnectionOffset)$);
      \coordinate (-west-middle) at (-shape.west);
      \coordinate (-west-lower)  at ($(-shape.west)+(0,-\mechConnectionOffset)$);

      \coordinate (-east-upper)  at ($(-shape.east)+(0,\mechConnectionOffset)$);
      \coordinate (-east-middle) at (-shape.east);
      \coordinate (-east-lower)  at ($(-shape.east)+(0,-\mechConnectionOffset)$);

      % Generic connection aliases
      \coordinate (-connection-upper)  at (-east-upper);
      \coordinate (-connection-middle) at (-east-middle);
      \coordinate (-connection-lower)  at (-east-lower);

      \coordinate (-north) at (-shape.north);
      \coordinate (-south) at (-shape.south);
    }
  },
  pics/mech/mass/.default={},
}
\pgfkeys{
  /mech/point-mass/.is family,
  /mech/point-mass,
  radius/.initial=2.2mm,
  label/.initial=,
  style/.initial=mech/body,
}
\tikzset{
  pics/mech/point-mass/.style args={#1}{
    code={
      \pgfkeys{/mech/point-mass/.cd,#1}
      \edef\mechPointMassRadius{\pgfkeysvalueof{/mech/point-mass/radius}}
      \edef\mechPointMassLabel{\pgfkeysvalueof{/mech/point-mass/label}}
      \edef\mechPointMassStyle{\pgfkeysvalueof{/mech/point-mass/style}}
      \node[
        \mechPointMassStyle,
        circle,
        minimum size=2\mechPointMassRadius,
        inner sep=0pt
      ] (-shape) at (0,0) {\mechPointMassLabel};

      \edef\mechPicName{\pgfkeysvalueof{/tikz/name prefix}}
      \ifx\mechPicName\pgfutil@empty\else
        \pgfnodealias{\mechPicName}{\mechPicName-shape}
      \fi

      \coordinate (-center) at (-shape.center);
    }
  },
  pics/mech/point-mass/.default={},
}
% -------------------------------------------------------------------------
% Wagen
% -------------------------------------------------------------------------
% Der Boolean "ground" wird als TeX-if gespeichert. Dadurch kann die optionale
% Bodendarstellung ohne Stringvergleich ein- bzw. ausgeschaltet werden.
\newif\ifmechCartGround
\pgfkeys{
  /mech/cart/.is family,
  /mech/cart,
  width/.initial=2cm,
  height/.initial=1cm,
  wheel radius/.initial=2mm,
  label/.initial=,
  ground/.is if=mechCartGround,
  ground/.default=true,
  ground=true,
}
\tikzset{
  pics/mech/cart/.style args={#1}{
    code={
      \pgfkeys{/mech/cart/.cd,#1}
      \edef\mechCartWidth{\pgfkeysvalueof{/mech/cart/width}}
      \edef\mechCartHeight{\pgfkeysvalueof{/mech/cart/height}}
      \edef\mechCartWheelRadius{\pgfkeysvalueof{/mech/cart/wheel radius}}
      \node[
        mech/body,
        minimum width=\mechCartWidth,
        minimum height=\mechCartHeight,
        inner sep=0pt
      ] (-shape) at (0,0) {\pgfkeysvalueof{/mech/cart/label}};

      \edef\mechPicName{\pgfkeysvalueof{/tikz/name prefix}}
      \ifx\mechPicName\pgfutil@empty\else
        \pgfnodealias{\mechPicName}{\mechPicName-shape}
      \fi

      \coordinate (-center) at (-shape.center);
      \coordinate (-west) at (-shape.west);
      \coordinate (-east) at (-shape.east);
      \coordinate (-north) at (-shape.north);
      \coordinate (-south) at (-shape.south);
      \draw[mech/body]
        (-.3\mechCartWidth,-.5\mechCartHeight-\mechCartWheelRadius)
        circle[radius=\mechCartWheelRadius];
      \draw[mech/body]
        (.3\mechCartWidth,-.5\mechCartHeight-\mechCartWheelRadius)
        circle[radius=\mechCartWheelRadius];
      \ifmechCartGround
        \draw
          (-.5\mechCartWidth,-.5\mechCartHeight-2\mechCartWheelRadius)
          --
          (.5\mechCartWidth,-.5\mechCartHeight-2\mechCartWheelRadius);
        \fill[
          pattern={Lines[angle=\mechHatchAngle,distance=1pt]},
          pattern color=black
        ]
          (-.5\mechCartWidth,-.5\mechCartHeight-2\mechCartWheelRadius)
          rectangle
          (.5\mechCartWidth,-.5\mechCartHeight-2\mechCartWheelRadius-3pt);
      \fi
    }
  },
  pics/mech/cart/.default={},
}
\pgfkeys{
  /mech/beam/.is family,
  /mech/beam,
  from/.initial={(0,0)},
  to/.initial={(1,0)},
  height/.initial=3mm,
  style/.initial=mech/body,
}
\tikzset{
  pics/mech/beam/.style args={#1}{
    code={
      \pgfkeys{/mech/beam/.cd,#1}
      \edef\mechBeamFrom{\pgfkeysvalueof{/mech/beam/from}}
      \edef\mechBeamTo{\pgfkeysvalueof{/mech/beam/to}}
      \edef\mechBeamHeight{\pgfkeysvalueof{/mech/beam/height}}
      \edef\mechBeamStyle{\pgfkeysvalueof{/mech/beam/style}}
      \path let
        \p1=\mechBeamFrom,
        \p2=\mechBeamTo,
        \n1={veclen(\x2-\x1,\y2-\y1)},
        \n2={(\y2-\y1)/\n1*\mechBeamHeight/2},
        \n3={-(\x2-\x1)/\n1*\mechBeamHeight/2}
      in
        coordinate (-A) at (\x1+\n2,\y1+\n3)
        coordinate (-B) at (\x2+\n2,\y2+\n3)
        coordinate (-C) at (\x2-\n2,\y2-\n3)
        coordinate (-D) at (\x1-\n2,\y1-\n3)
        coordinate (-start) at (\x1,\y1)
        coordinate (-end) at (\x2,\y2)
        coordinate (-center) at ({(\x1+\x2)/2},{(\y1+\y2)/2});
      \draw[\mechBeamStyle] (-A)--(-B)--(-C)--(-D)--cycle;
    }
  },
  pics/mech/beam/.default={},
}

% =========================================================
% mechlib/elements/supports.tex
% ATOMARE LAGER: FESTLAGER, LOSLAGER UND EINSPANNUNG
% =========================================================
% Diese Datei kapselt die Lagergeometrie. Fest- und Loslager teilen sich die
% Key-Familie /mech/support-common, damit Breite, Hoehe, Schraffur und Label
% konsistent parametrisiert werden koennen.
%
% Oeffentliche Pics:
%   mech/pinned-support Festlager / gelenkiges Lager
%   mech/roller-support Loslager mit sichtbarem Abstand zum Boden
%   mech/fixed-support  Einspannung, optional gedreht
%
% Aktuelle Fachwerk-Konvention:
%   - Die Dreiecksspitze liegt exakt im Mittelpunkt des Gelenks (0,0).
%   - Das Dreieck ist gegenueber der frueheren Version kompakt skaliert.
%   - Bodenlinie und Schraffur verwenden dieselben x-Grenzen wie die
%     Dreiecksgrundseite; beim Loslager wird die sichtbare Linienbreite explizit
%     beruecksichtigt.
%   - Das Gelenk selbst ist mit mechBodyColor gefuellt.
%
% Die pgfmathsetlengthmacro-Aufrufe sind absichtlich verwendet: LaTeX-Laengen
% duerfen nicht durch Textersatz wie ".5\\macro" skaliert werden, da dies bei
% expandierenden Dimensionsmakros zu sichtbarem Resttext (z. B. ".8mm") fuehren
% kann. Alle abgeleiteten Abmessungen werden deshalb numerisch berechnet.
% =========================================================
\pgfkeys{
  /mech/support-common/.is family,
  /mech/support-common,
  width/.initial=2.72mm,
  height/.initial=2.24mm,
  hatch angle/.initial=45,
  label/.initial=,
  label position/.initial=above,
}
\tikzset{
  % -----------------------------------------------------------------------
  % Festlager
  % -----------------------------------------------------------------------
  pics/mech/pinned-support/.style args={#1}{
    code={
      \pgfkeys{/mech/support-common/.cd,#1}
      \edef\mechSupportWidth{\pgfkeysvalueof{/mech/support-common/width}}
      \edef\mechSupportHeight{\pgfkeysvalueof{/mech/support-common/height}}
      \edef\mechSupportHatch{\pgfkeysvalueof{/mech/support-common/hatch angle}}
      \pgfmathsetlengthmacro{\mechSupportHalfWidth}{0.5*\mechSupportWidth}
      % Die Dreieckskontur ragt an der Grundseite durch die halbe Linienbreite
      % minimal ueber die geometrische Breite hinaus. Fuer identische sichtbare
      % Breite wird die Bodenlinie daher um 0.3pt je Seite erweitert.
      \pgfmathsetlengthmacro{\mechSupportGroundHalfWidth}{\mechSupportHalfWidth+0.3pt}
      \pgfmathsetlengthmacro{\mechSupportApexY}{0mm}
      \pgfmathsetlengthmacro{\mechSupportHatchDepth}{0.256*\mechSupportHeight}

      \coordinate (-connection) at (0,0);
      \coordinate (-center) at (0,0);

      \draw[mech/support]
        (-\mechSupportHalfWidth,-\mechSupportHeight)
        -- (0,\mechSupportApexY)
        -- (\mechSupportHalfWidth,-\mechSupportHeight)
        -- cycle;
      \draw[mech/body] (0,0) circle[radius=.6mm];
      \draw[black,line width=.6pt,line cap=butt]
        (-\mechSupportGroundHalfWidth,-\mechSupportHeight)
        -- (\mechSupportGroundHalfWidth,-\mechSupportHeight);
      \fill[
        pattern={Lines[angle=\mechSupportHatch,distance=1pt]},
        pattern color=black
      ]
        (-\mechSupportGroundHalfWidth,-\mechSupportHeight)
        rectangle
        (\mechSupportGroundHalfWidth,{-(\mechSupportHeight+\mechSupportHatchDepth)});

      \node[\pgfkeysvalueof{/mech/support-common/label position}]
        at (0,0) {\pgfkeysvalueof{/mech/support-common/label}};
    }
  },
  pics/mech/pinned-support/.default={},
  % -----------------------------------------------------------------------
  % Loslager
  % -----------------------------------------------------------------------
  pics/mech/roller-support/.style args={#1}{
    code={
      \pgfkeys{/mech/support-common/.cd,#1}
      \edef\mechSupportWidth{\pgfkeysvalueof{/mech/support-common/width}}
      \edef\mechSupportHeight{\pgfkeysvalueof{/mech/support-common/height}}
      \edef\mechSupportHatch{\pgfkeysvalueof{/mech/support-common/hatch angle}}
      \pgfmathsetlengthmacro{\mechSupportHalfWidth}{0.5*\mechSupportWidth}
      % The triangle outline extends by half its line width at the base corners.
      % Add 0.3pt so the roller ground has the same visible total width.
      \pgfmathsetlengthmacro{\mechSupportGroundHalfWidth}{\mechSupportHalfWidth+0.3pt}
      \pgfmathsetlengthmacro{\mechSupportApexY}{0mm}
      \pgfmathsetlengthmacro{\mechSupportRollerGap}{0.272*\mechSupportHeight}
      \pgfmathsetlengthmacro{\mechSupportHatchDepth}{0.256*\mechSupportHeight}
      \pgfmathsetlengthmacro{\mechSupportGroundY}{\mechSupportHeight+\mechSupportRollerGap}

      \coordinate (-connection) at (0,0);
      \coordinate (-center) at (0,0);

      \draw[mech/support]
        (-\mechSupportHalfWidth,-\mechSupportHeight)
        -- (0,\mechSupportApexY)
        -- (\mechSupportHalfWidth,-\mechSupportHeight)
        -- cycle;
      \draw[mech/body] (0,0) circle[radius=.6mm];
      \draw[black,line width=.6pt,line cap=butt]
        (-\mechSupportGroundHalfWidth,-\mechSupportGroundY)
        -- (\mechSupportGroundHalfWidth,-\mechSupportGroundY);
      \fill[
        pattern={Lines[angle=\mechSupportHatch,distance=1pt]},
        pattern color=black
      ]
        (-\mechSupportGroundHalfWidth,-\mechSupportGroundY)
        rectangle
        (\mechSupportGroundHalfWidth,{-(\mechSupportGroundY+\mechSupportHatchDepth)});

      \node[\pgfkeysvalueof{/mech/support-common/label position}]
        at (0,0) {\pgfkeysvalueof{/mech/support-common/label}};
    }
  },
  pics/mech/roller-support/.default={},
}
\pgfkeys{
  /mech/fixed-support/.is family,
  /mech/fixed-support,
  width/.initial=1cm,
  depth/.initial=1.5mm,
  hatch angle/.initial=45,
  angle/.initial=0,
  label/.initial=,
}
\tikzset{
  pics/mech/fixed-support/.style args={#1}{
    code={
      \pgfkeys{/mech/fixed-support/.cd,#1}
      \edef\mechFixedWidth{\pgfkeysvalueof{/mech/fixed-support/width}}
      \edef\mechFixedDepth{\pgfkeysvalueof{/mech/fixed-support/depth}}
      \edef\mechFixedHatch{\pgfkeysvalueof{/mech/fixed-support/hatch angle}}
      \edef\mechFixedAngle{\pgfkeysvalueof{/mech/fixed-support/angle}}

      \coordinate (-center) at (0,0);
      \coordinate (-connection-upper)  at (0,\mechConnectionOffset);
      \coordinate (-connection-middle) at (0,0);
      \coordinate (-connection-lower)  at (0,-\mechConnectionOffset);

      \begin{scope}[rotate=\mechFixedAngle]
        \draw[mech/body,fill=none] (-.5\mechFixedWidth,0)--(.5\mechFixedWidth,0);
        \fill[
          pattern={Lines[angle=\mechFixedHatch,distance={3pt/sqrt(2)}]},
          pattern color=black
        ]
          (-.5\mechFixedWidth,0)
          rectangle
          (.5\mechFixedWidth,\mechFixedDepth);
      \end{scope}

      \node[above] at (0,0) {\pgfkeysvalueof{/mech/fixed-support/label}};
    }
  },
  pics/mech/fixed-support/.default={},
}

% =========================================================
% mechlib/elements/connectors.tex
% ATOMARE VERBINDUNGSELEMENTE: FEDER UND DAEMPFER
% =========================================================
% Beide Elemente folgen derselben grundlegenden API: "from", "to" und "label".
% Die Endpunkte koennen beliebige TikZ-Koordinaten oder exportierte Anker aus
% anderen mechlib-Pics sein.
%
% Oeffentliche Pics:
%   mech/spring  Zickzack-Feder entlang der Verbindungslinie
%   mech/damper  Daempfer mit rechteckigem Daempferkoerper in Linienmitte
%
% Die Beschriftung wird "sloped" gezeichnet und folgt daher der Orientierung
% des Verbindungselements. Ueber "label side" kann sie ober- oder unterhalb
% der lokalen Verbindungslinie platziert werden.
% =========================================================
\pgfkeys{
  /mech/spring/.is family,
  /mech/spring,
  from/.initial={(0,0)},
  to/.initial={(1,0)},
  label/.initial=,
  label side/.initial=above,
  amplitude/.initial=2mm,
  segment length/.initial=4mm,
  style/.initial=mech/connector,
}
\tikzset{
  pics/mech/spring/.style args={#1}{
    code={
      \pgfkeys{/mech/spring/.cd,#1}
      \edef\mechSpringFrom{\pgfkeysvalueof{/mech/spring/from}}
      \edef\mechSpringTo{\pgfkeysvalueof{/mech/spring/to}}
      \draw[\pgfkeysvalueof{/mech/spring/style},
        decorate,
        decoration={zigzag,
          segment length=\pgfkeysvalueof{/mech/spring/segment length},
          amplitude=\pgfkeysvalueof{/mech/spring/amplitude}}]
        \mechSpringFrom -- \mechSpringTo
        node[midway,sloped,mech/label,outer sep=2mm,\pgfkeysvalueof{/mech/spring/label side}]
        {\pgfkeysvalueof{/mech/spring/label}};
    }
  },
  pics/mech/spring/.default={},
}
\pgfkeys{
  /mech/damper/.is family,
  /mech/damper,
  from/.initial={(0,0)},
  to/.initial={(1,0)},
  label/.initial=,
  label side/.initial=below,
  body width/.initial=8mm,
  body height/.initial=3mm,
  style/.initial=mech/connector,
}
\tikzset{
  pics/mech/damper/.style args={#1}{
    code={
      \pgfkeys{/mech/damper/.cd,#1}
      \edef\mechDamperFrom{\pgfkeysvalueof{/mech/damper/from}}
      \edef\mechDamperTo{\pgfkeysvalueof{/mech/damper/to}}
      \draw[\pgfkeysvalueof{/mech/damper/style}]
        \mechDamperFrom -- \mechDamperTo
        node[midway,sloped,draw,fill=white,
          minimum width=\pgfkeysvalueof{/mech/damper/body width},
          minimum height=\pgfkeysvalueof{/mech/damper/body height},
          inner sep=0pt] {};
      \path \mechDamperFrom -- \mechDamperTo
        node[midway,sloped,mech/label,outer sep=2mm,\pgfkeysvalueof{/mech/damper/label side}]
        {\pgfkeysvalueof{/mech/damper/label}};
    }
  },
  pics/mech/damper/.default={},
}

% =========================================================
% mechlib/elements/trusses.tex
% FACHWERKE: STAB UND GEMEINSAME GELENKDARSTELLUNG
% =========================================================
% Dieses Modul enthaelt die atomaren Fachwerkelemente. Die Topologie eines
% Fachwerks bleibt bewusst im Anwendungsdokument: Knoten werden als TikZ-
% Koordinaten definiert und durch mech/rod bzw. \mechJoints verbunden.
%
% Oeffentliche API:
%   mech/rod     Stab zwischen "from" und "to" mit optionalem Label
%   \mechJoints{A,B,C,...} zeichnet alle angegebenen Gelenke
%
% Darstellungsregeln:
%   - Staebe verwenden standardmaessig den Stil mech/rod: schwarz und duenn.
%   - Staebe werden auf dem Layer "mech background" gezeichnet.
%   - Stabnummern besitzen keine Fuellung (fill=none) und bleiben transparent.
%   - Gelenke liegen optisch vor den Staeben und sind mit mechBodyColor gefuellt.
%
% Beispiel:
%   \pic {mech/rod={from=(A),to=(B),label=1,label side=above}};
%   \mechJoints{A,B}
% =========================================================
% -------------------------------------------------------------------------
% Einzelner Fachwerkstab
% -------------------------------------------------------------------------
% "style" erlaubt lokale Abweichungen, ohne den globalen mech/rod-Stil zu
% veraendern. Die Standardbeschriftung liegt an der Stabmitte.
\pgfkeys{
  /mech/rod/.is family,
  /mech/rod,
  from/.initial={(0,0)},
  to/.initial={(1,0)},
  label/.initial=,
  label side/.initial=above,
  label style/.initial={mech/label,fill=none,text=black},
  style/.initial=mech/rod,
}
\tikzset{
  pics/mech/rod/.style args={#1}{
    code={
      \pgfkeys{/mech/rod/.cd,#1}
      \edef\mechRodFrom{\pgfkeysvalueof{/mech/rod/from}}
      \edef\mechRodTo{\pgfkeysvalueof{/mech/rod/to}}
      \coordinate (-start) at \mechRodFrom;
      \coordinate (-end) at \mechRodTo;
      \coordinate (-center) at ($(-start)!.5!(-end)$);
      \begin{pgfonlayer}{mech background}
        \draw[\pgfkeysvalueof{/mech/rod/style}] (-start)--(-end);
      \end{pgfonlayer}
      \node[\pgfkeysvalueof{/mech/rod/label style},\pgfkeysvalueof{/mech/rod/label side}]
        at (-center) {\pgfkeysvalueof{/mech/rod/label}};
    }
  },
  pics/mech/rod/.default={},
}
% -------------------------------------------------------------------------
% Gelenkliste
% -------------------------------------------------------------------------
% Das Makro erwartet eine kommaseparierte Liste bereits definierter TikZ-
% Koordinaten. Jedes Gelenk wird identisch gezeichnet. Die Gelenke sollten nach
% den Staeben aufgerufen werden, damit sie die Stabenden optisch ueberdecken.
\newcommand{\mechJoints}[1]{%
  \foreach \mechJoint in {#1}{%
    \draw[fill=mechBodyColor,draw=black,line width=.2mm]
      (\mechJoint) circle[radius=.6mm];%
  }%
}

% -------------------------------------------------------------------------
% Annotationen und mechanische Symbolgruppen.
% -------------------------------------------------------------------------
% =========================================================
% mechlib/annotations/loads.tex
% LASTANNOTATIONEN: KRAFT, MOMENT UND STRECKENLAST
% =========================================================
% Diese Datei stellt die Lastsymbole der mechlib bereit. Alle Lasten verwenden
% standardmaessig den zentralen Stil mech/load und damit mechLoadColor.
%
% Oeffentliche Pics:
%   mech/force            Einzelkraft als Winkel/Laenge oder from/to
%   mech/moment           Kreisbogensymbol im oder gegen den Uhrzeigersinn
%   mech/distributed-load konstante, dreieckige oder linear veraenderliche Last
%
% Besonderheit der Streckenlast:
% Die Grundlinie wird lokal auf die x-Achse gedreht. Dadurch kann die gesamte
% Geometrie in einem einfachen lokalen Koordinatensystem berechnet werden. Sehr
% kurze Pfeile, deren Laenge kleiner als ihre Pfeilspitze ist, werden bewusst
% unterdrueckt; so entstehen am Nullende einer Dreieckslast keine Artefakte.
% =========================================================
\pgfkeys{
  /mech/force/.is family,
  /mech/force,
  from/.initial={(0,0)},
  to/.initial=,
  angle/.initial=0,
  length/.initial=1.5cm,
  label/.initial=,
  label position/.initial=above,
  direction/.initial=away,
  style/.initial=mech/load,
}
\tikzset{
  pics/mech/force/.style args={#1}{
    code={
      \pgfkeys{/mech/force/.cd,#1}
      \edef\mechForceFrom{\pgfkeysvalueof{/mech/force/from}}
      \edef\mechForceTo{\pgfkeysvalueof{/mech/force/to}}
      \edef\mechForceDirection{\pgfkeysvalueof{/mech/force/direction}}
      \def\mechForceToward{toward}
      \ifx\mechForceTo\empty
        \ifx\mechForceDirection\mechForceToward
          \draw[\pgfkeysvalueof{/mech/force/style},-{Latex[length=2.5mm]}]
            (\pgfkeysvalueof{/mech/force/angle}:\pgfkeysvalueof{/mech/force/length}) -- (0,0)
            node[midway,\pgfkeysvalueof{/mech/force/label position}]
            {\pgfkeysvalueof{/mech/force/label}};
        \else
          \draw[\pgfkeysvalueof{/mech/force/style},-{Latex[length=2.5mm]}]
            (0,0) -- (\pgfkeysvalueof{/mech/force/angle}:\pgfkeysvalueof{/mech/force/length})
            node[midway,\pgfkeysvalueof{/mech/force/label position}]
            {\pgfkeysvalueof{/mech/force/label}};
        \fi
      \else
        \ifx\mechForceDirection\mechForceToward
          \draw[\pgfkeysvalueof{/mech/force/style},-{Latex[length=2.5mm]}]
            \mechForceTo -- \mechForceFrom
            node[midway,\pgfkeysvalueof{/mech/force/label position}]
            {\pgfkeysvalueof{/mech/force/label}};
        \else
          \draw[\pgfkeysvalueof{/mech/force/style},-{Latex[length=2.5mm]}]
            \mechForceFrom -- \mechForceTo
            node[midway,\pgfkeysvalueof{/mech/force/label position}]
            {\pgfkeysvalueof{/mech/force/label}};
        \fi
      \fi
    }
  },
  pics/mech/force/.default={},
}
\pgfkeys{
  /mech/moment/.is family,
  /mech/moment,
  radius/.initial=6mm,
  label/.initial=,
  direction/.is choice,
  direction/ccw/.code={\def\mechMomentStart{-60}\def\mechMomentEnd{220}},
  direction/cw/.code={\def\mechMomentStart{240}\def\mechMomentEnd{-40}},
  direction=ccw,
  style/.initial=mech/load,
}
\tikzset{
  pics/mech/moment/.style args={#1}{
    code={
      \pgfkeys{/mech/moment/.cd,direction=ccw,#1}
      \draw[\pgfkeysvalueof{/mech/moment/style},-{Latex[length=2.5mm]}]
        (\mechMomentStart:\pgfkeysvalueof{/mech/moment/radius})
        arc[start angle=\mechMomentStart,end angle=\mechMomentEnd,
            radius=\pgfkeysvalueof{/mech/moment/radius}];
      \node[mech/label,above] at (0,\pgfkeysvalueof{/mech/moment/radius})
        {\pgfkeysvalueof{/mech/moment/label}};
    }
  },
  pics/mech/moment/.default={},
}
\pgfkeys{
  /mech/distributed-load/.is family,
  /mech/distributed-load,
  from/.initial={(0,0)},
  to/.initial={(1,0)},
  start value/.initial=1,
  end value/.initial=1,
  height/.initial=8mm,
  offset/.initial=0mm,
  arrows/.initial=6,
  arrowhead length/.initial=2mm,
  label/.initial=,
  label position/.initial=above,
  style/.initial=mech/load,
}
\tikzset{
  pics/mech/distributed-load/.style args={#1}{
    code={
      \pgfkeys{/mech/distributed-load/.cd,#1}
      \edef\mechDLFrom{\pgfkeysvalueof{/mech/distributed-load/from}}
      \edef\mechDLTo{\pgfkeysvalueof{/mech/distributed-load/to}}
      \path let
        \p1=\mechDLFrom,\p2=\mechDLTo,
        \n1={veclen(\x2-\x1,\y2-\y1)},
        \n2={atan2(\y2-\y1,\x2-\x1)}
      in
        coordinate (-start) at (\x1,\y1)
        coordinate (-end) at (\x2,\y2)
        coordinate (-helper) at (\n2,\n1);
      \path (-helper); \pgfgetlastxy{\mechDLAngle}{\mechDLLength};
      \pgfmathsetmacro{\mechDLStartValue}{\pgfkeysvalueof{/mech/distributed-load/start value}}
      \pgfmathsetmacro{\mechDLEndValue}{\pgfkeysvalueof{/mech/distributed-load/end value}}
      \pgfmathsetmacro{\mechDLArrowCount}{\pgfkeysvalueof{/mech/distributed-load/arrows}}
      \begin{scope}[shift={(-start)},rotate=\mechDLAngle]
        \pgfmathsetlengthmacro{\mechDLOffset}{\pgfkeysvalueof{/mech/distributed-load/offset}}
        \pgfmathsetlengthmacro{\mechDLArrowheadLength}{\pgfkeysvalueof{/mech/distributed-load/arrowhead length}}
        \pgfmathsetlengthmacro{\mechDLStartArrowLength}{abs(\mechDLStartValue)*\pgfkeysvalueof{/mech/distributed-load/height}}

        \fill[mech/load fill]
          (0,\mechDLOffset)
          --(0,{\mechDLOffset+\mechDLStartValue*\pgfkeysvalueof{/mech/distributed-load/height}})
          --(\mechDLLength,{\mechDLOffset+\mechDLEndValue*\pgfkeysvalueof{/mech/distributed-load/height}})
          --(\mechDLLength,\mechDLOffset)
          --cycle;

        % Keep contour and first arrow continuous whenever the first arrow
        % is long enough to accommodate its arrowhead. Otherwise draw only
        % the contour and suppress the too-short arrow.
        \ifdim\mechDLStartArrowLength<\mechDLArrowheadLength
          \draw[\pgfkeysvalueof{/mech/distributed-load/style}]
            (\mechDLLength,{\mechDLOffset+\mechDLEndValue*\pgfkeysvalueof{/mech/distributed-load/height}})
            --(0,{\mechDLOffset+\mechDLStartValue*\pgfkeysvalueof{/mech/distributed-load/height}});
        \else
          \draw[
            \pgfkeysvalueof{/mech/distributed-load/style},
            -{Latex[length=\pgfkeysvalueof{/mech/distributed-load/arrowhead length}]}
          ]
            (\mechDLLength,{\mechDLOffset+\mechDLEndValue*\pgfkeysvalueof{/mech/distributed-load/height}})
            --(0,{\mechDLOffset+\mechDLStartValue*\pgfkeysvalueof{/mech/distributed-load/height}})
            --(0,\mechDLOffset);
        \fi

        \foreach \mechDLIndex in {1,...,\mechDLArrowCount}{
          \pgfmathsetmacro{\mechDLT}{\mechDLIndex/\mechDLArrowCount}
          \pgfmathsetmacro{\mechDLValue}{\mechDLStartValue+(\mechDLEndValue-\mechDLStartValue)*\mechDLT}
          \pgfmathsetlengthmacro{\mechDLArrowLength}{abs(\mechDLValue)*\pgfkeysvalueof{/mech/distributed-load/height}}

          \ifdim\mechDLArrowLength<\mechDLArrowheadLength
            % Suppress arrows that are shorter than their arrowhead.
          \else
            \draw[
              \pgfkeysvalueof{/mech/distributed-load/style},
              -{Latex[length=\pgfkeysvalueof{/mech/distributed-load/arrowhead length}]}
            ]
              ({\mechDLT*\mechDLLength},{\mechDLOffset+\mechDLValue*\pgfkeysvalueof{/mech/distributed-load/height}})
              --({\mechDLT*\mechDLLength},\mechDLOffset);
          \fi
        }

        \node[mech/label,\pgfkeysvalueof{/mech/distributed-load/label position}]
          at (\mechDLLength,{\mechDLOffset+\mechDLEndValue*\pgfkeysvalueof{/mech/distributed-load/height}})
          {\pgfkeysvalueof{/mech/distributed-load/label}};
      \end{scope}
    }
  },
  pics/mech/distributed-load/.default={},
}

% =========================================================
% mechlib/annotations/dimensions.tex
% BEMASSUNGEN: LINEARES MASS UND WINKELMASS
% =========================================================
% Lineare Bemassungen werden geometrisch aus zwei Endpunkten abgeleitet. Der
% Offset wirkt senkrecht auf die Verbindung von "from" nach "to"; damit
% funktioniert dasselbe Pic fuer horizontale, vertikale und schraege Masse.
%
% Oeffentliche Pics:
%   mech/dimension       lineare Bemassung mit Doppelpfeil
%   mech/angle-dimension Kreisbogen zwischen zwei Winkeln
%
% Interne Koordinaten -A und -B sind die senkrecht versetzten Endpunkte der
% Masslinie. Die Vektorrechnung erfolgt in TikZ "let"-Syntax.
% =========================================================
\pgfkeys{
  /mech/dimension/.is family,
  /mech/dimension,
  from/.initial={(0,0)},
  to/.initial={(1,0)},
  label/.initial={$l$},
  offset/.initial=3mm,
  style/.initial=mech/dimension,
}
\tikzset{
  pics/mech/dimension/.style args={#1}{
    code={
      \pgfkeys{/mech/dimension/.cd,#1}
      \edef\mechDimFrom{\pgfkeysvalueof{/mech/dimension/from}}
      \edef\mechDimTo{\pgfkeysvalueof{/mech/dimension/to}}
      \path let
        \p1=\mechDimFrom,\p2=\mechDimTo,
        \n1={veclen(\x2-\x1,\y2-\y1)},
        \n2={-(\y2-\y1)/\n1*\pgfkeysvalueof{/mech/dimension/offset}},
        \n3={(\x2-\x1)/\n1*\pgfkeysvalueof{/mech/dimension/offset}}
      in
        coordinate (-A) at (\x1+\n2,\y1+\n3)
        coordinate (-B) at (\x2+\n2,\y2+\n3);
      \draw[\pgfkeysvalueof{/mech/dimension/style}]
        (-A)--(-B)
        node[midway,fill=white,inner sep=1pt]
        {\pgfkeysvalueof{/mech/dimension/label}};
    }
  },
  pics/mech/dimension/.default={},
}
\pgfkeys{
  /mech/angle-dimension/.is family,
  /mech/angle-dimension,
  start angle/.initial=0,
  end angle/.initial=45,
  radius/.initial=7mm,
  label/.initial={$\alpha$},
}
\tikzset{
  pics/mech/angle-dimension/.style args={#1}{
    code={
      \pgfkeys{/mech/angle-dimension/.cd,#1}
      \draw
        (\pgfkeysvalueof{/mech/angle-dimension/start angle}:\pgfkeysvalueof{/mech/angle-dimension/radius})
        arc[start angle=\pgfkeysvalueof{/mech/angle-dimension/start angle},
            end angle=\pgfkeysvalueof{/mech/angle-dimension/end angle},
            radius=\pgfkeysvalueof{/mech/angle-dimension/radius}];
      \pgfmathsetmacro{\mechAngleMid}{.5*(\pgfkeysvalueof{/mech/angle-dimension/start angle}+\pgfkeysvalueof{/mech/angle-dimension/end angle})}
      \node at (\mechAngleMid:{.7*\pgfkeysvalueof{/mech/angle-dimension/radius}})
        {\pgfkeysvalueof{/mech/angle-dimension/label}};
    }
  },
  pics/mech/angle-dimension/.default={},
}

% =========================================================
% mechlib/annotations/coordinates.tex
% KOORDINATENANNOTATIONEN: EINZELACHSE UND 2D-KOORDINATENSYSTEM
% =========================================================
% Die Pics dieser Datei zeichnen mathematische Bezugsachsen im Stil mech/axis.
% Winkel werden in TikZ-Grad angegeben; Laengen sind regulaere TeX-Dimensionen.
%
% Oeffentliche Pics:
%   mech/axis              einzelne gerichtete Achse
%   mech/coordinate-system zwei unabhaengig orientierbare Achsen
%
% Beide Pics exportieren den lokalen Ursprung als "-origin". Bei einem benannten
% Pic (cs) kann dieser z. B. als (cs-origin) referenziert werden.
% =========================================================
\pgfkeys{
  /mech/axis/.is family,
  /mech/axis,
  label/.initial={$x$},
  angle/.initial=0,
  length/.initial=1cm,
  direction/.initial=positive,
}
\tikzset{
  pics/mech/axis/.style args={#1}{
    code={
      \pgfkeys{/mech/axis/.cd,#1}
      \edef\mechAxisDirection{\pgfkeysvalueof{/mech/axis/direction}}
      \def\mechAxisNegative{negative}
      \ifx\mechAxisDirection\mechAxisNegative
        \draw[mech/axis,{Latex[length=2mm]}-]
          (0,0)--(\pgfkeysvalueof{/mech/axis/angle}:\pgfkeysvalueof{/mech/axis/length})
          node[pos=1,anchor=west]{\pgfkeysvalueof{/mech/axis/label}};
      \else
        \draw[mech/axis,-{Latex[length=2mm]}]
          (0,0)--(\pgfkeysvalueof{/mech/axis/angle}:\pgfkeysvalueof{/mech/axis/length})
          node[pos=1,anchor=west]{\pgfkeysvalueof{/mech/axis/label}};
      \fi
      \coordinate (-origin) at (0,0);
    }
  },
  pics/mech/axis/.default={},
}
\pgfkeys{
  /mech/coordinate-system/.is family,
  /mech/coordinate-system,
  x label/.initial={$x$},
  y label/.initial={$y$},
  x angle/.initial=0,
  y angle/.initial=90,
  x length/.initial=1cm,
  y length/.initial=1cm,
}
\tikzset{
  pics/mech/coordinate-system/.style args={#1}{
    code={
      \pgfkeys{/mech/coordinate-system/.cd,#1}
      \draw[mech/axis,-{Latex[length=2mm]}]
        (0,0)--(\pgfkeysvalueof{/mech/coordinate-system/x angle}:\pgfkeysvalueof{/mech/coordinate-system/x length})
        node[pos=1,anchor=west]{\pgfkeysvalueof{/mech/coordinate-system/x label}};
      \draw[mech/axis,-{Latex[length=2mm]}]
        (0,0)--(\pgfkeysvalueof{/mech/coordinate-system/y angle}:\pgfkeysvalueof{/mech/coordinate-system/y length})
        node[pos=1,anchor=east]{\pgfkeysvalueof{/mech/coordinate-system/y label}};
      \coordinate (-origin) at (0,0);
    }
  },
  pics/mech/coordinate-system/.default={},
}

% =========================================================
% mechlib/annotations/mechanics.tex
% MECHANISCHE ANNOTATIONEN: SCHNITTGROESSEN UND LAGERREAKTIONEN
% =========================================================
% Die Pics kombinieren mehrere elementare Lastsymbole zu haeufig benoetigten
% mechanischen Symbolgruppen. Farben und Pfeilgestaltung folgen mech/load.
%
% Oeffentliche Pics:
%   mech/internal-forces  Normalkraft N, Querkraft Q und Moment M am Schnitt
%   mech/fixed-reactions  horizontale/vertikale Reaktion und Einspannmoment
%
% Bei mech/internal-forces steuert "side=right|left" die Vorzeichenrichtung
% aller drei Schnittgroessen konsistent. Das Moment wird dazu intern als
% mech/moment mit passender Drehrichtung wiederverwendet.
% =========================================================
\pgfkeys{
  /mech/internal-forces/.is family,
  /mech/internal-forces,
  normal/.initial={$N$},
  shear/.initial={$Q$},
  moment/.initial={$M$},
  length/.initial=7mm,
  side/.is choice,
  side/right/.code={\def\mechIFSign{1}},
  side/left/.code={\def\mechIFSign{-1}},
  side=right,
}
\tikzset{
  pics/mech/internal-forces/.style args={#1}{
    code={
      \pgfkeys{/mech/internal-forces/.cd,side=right,#1}
      \edef\mechIFLength{\pgfkeysvalueof{/mech/internal-forces/length}}
      \draw[mech/load,-{Latex[length=2mm]}]
        (0,0)--({\mechIFSign*\mechIFLength},0)
        node[pos=1,anchor=west]{\pgfkeysvalueof{/mech/internal-forces/normal}};
      \draw[mech/load,-{Latex[length=2mm]}]
        (0,0)--(0,{-.8*\mechIFSign*\mechIFLength})
        node[pos=1,anchor=west]{\pgfkeysvalueof{/mech/internal-forces/shear}};
      \ifnum\mechIFSign=1
        \pic {mech/moment={direction=ccw,radius=.55\mechIFLength,label=\pgfkeysvalueof{/mech/internal-forces/moment}}};
      \else
        \pic {mech/moment={direction=cw,radius=.55\mechIFLength,label=\pgfkeysvalueof{/mech/internal-forces/moment}}};
      \fi
    }
  },
  pics/mech/internal-forces/.default={},
}
\pgfkeys{
  /mech/fixed-reactions/.is family,
  /mech/fixed-reactions,
  x/.initial={$A_x$},
  y/.initial={$A_y$},
  moment/.initial={$M_A$},
  length/.initial=7mm,
}
\tikzset{
  pics/mech/fixed-reactions/.style args={#1}{
    code={
      \pgfkeys{/mech/fixed-reactions/.cd,#1}
      \edef\mechReactionLength{\pgfkeysvalueof{/mech/fixed-reactions/length}}
      \draw[mech/load,-{Latex[length=2mm]}] (-\mechReactionLength,0)--(0,0)
        node[pos=0,anchor=east]{\pgfkeysvalueof{/mech/fixed-reactions/x}};
      \draw[mech/load,-{Latex[length=2mm]}] (0,.7\mechReactionLength)--(0,-.3\mechReactionLength)
        node[pos=0,anchor=south]{\pgfkeysvalueof{/mech/fixed-reactions/y}};
      \pic {mech/moment={direction=ccw,radius=.45\mechReactionLength,label=\pgfkeysvalueof{/mech/fixed-reactions/moment}}};
    }
  },
  pics/mech/fixed-reactions/.default={},
}


%
% Lade-Reihenfolge:
%   1. core/        Farben, Layer, Stile und globale Geometriekonventionen
%   2. elements/    mechanische Grundelemente (Koerper, Lager, Verbinder, Staebe)
%   3. annotations/ Lasten, Masse, Koordinaten und mechanische Annotationen
%
% WICHTIG:
% Die Modulreihenfolge ist absichtlich festgelegt. Elemente greifen auf Farben,
% Layer und Stile aus core/ zurueck; Annotationen koennen wiederum die bereits
% definierten Grundstile und Pics verwenden.
% =========================================================
% -------------------------------------------------------------------------
% Externe Pakete und TikZ-Bibliotheken
% -------------------------------------------------------------------------
% \RequirePackage ist bewusst statt \usepackage gewaehlt, da mechlib.tex selbst
% wie eine Bibliotheks-/Paketdatei verwendet wird.
\RequirePackage{tikz}
\RequirePackage{amsmath}
\RequirePackage{siunitx}
\RequirePackage{xcolor}
\usetikzlibrary{
  arrows.meta,
  calc,
  patterns,
  patterns.meta,
  intersections,
  angles,
  quotes,
  decorations.pathmorphing,
  backgrounds
}
% Base path: override before loading when mechlib is stored elsewhere.
% -------------------------------------------------------------------------
% Aufloesung des Bibliothekspfades
% -------------------------------------------------------------------------
% \providecommand respektiert eine bereits im Hauptdokument gesetzte Definition.
\providecommand{\mechlibpath}{mechlib}
% -------------------------------------------------------------------------
% Kernmodule: muessen vor allen Elementen geladen werden.
% -------------------------------------------------------------------------
% =========================================================
% mechlib/core/colors.tex
% ZENTRALE FARBPALETTE
% =========================================================
% Alle bibliotheksweiten Farben werden ausschliesslich hier definiert. Dadurch
% kann der visuelle Charakter der gesamten mechlib an einer Stelle angepasst
% werden, ohne einzelne Elemente zu veraendern.
%
% Konvention:
%   mechBodyColor    Fuellung allgemeiner Koerper und Fachwerkgelenke
%   mechSupportColor Fuellung von Lagerkoerpern
%   mechLoadColor    Kraefte, Momente und Streckenlasten
%   mechAxisColor    Achsen und Koordinatensysteme
%   mechRodColor     historische Stabfarbe; aktuelle Fachwerkstaebe sind schwarz
%
% Die Namen sind Teil der internen visuellen API und sollten bevorzugt statt
% harter RGB-Werte in den Elementdateien verwendet werden.
% =========================================================
% -------------------------------------------------------------------------
% Farbdefinitionen
% -------------------------------------------------------------------------
% RGB wird bewusst explizit verwendet, damit die Darstellung unabhaengig von
% Dokumentklassen und externen Farbmodellen reproduzierbar bleibt.
\definecolor{mechBodyColor}{RGB}{180,180,180}
\definecolor{mechSupportColor}{RGB}{205,205,205}
\definecolor{mechLoadColor}{RGB}{178,25,25}
\definecolor{mechAxisColor}{RGB}{6,13,141}
\definecolor{mechRodColor}{RGB}{130,130,130}

% =========================================================
% mechlib/core/layers.tex
% ZEICHENEBENEN / PGF-LAYER
% =========================================================
% Fachwerkstaebe und andere hinterlegte Geometrien koennen auf dem Layer
% "mech background" gezeichnet werden. Der normale TikZ-Layer "main" liegt
% darueber. Dadurch koennen z. B. gefuellte Gelenke sauber vor den Staeben
% erscheinen, ohne dass die Zeichenreihenfolge im Anwendungsdokument aufwendig
% kontrolliert werden muss.
% =========================================================
% Zuerst wird der benannte Hintergrundlayer deklariert. Anschliessend wird die
% globale Layerreihenfolge explizit gesetzt: Hintergrund vor Hauptlayer.
\pgfdeclarelayer{mech background}
\pgfsetlayers{mech background,main}

% =========================================================
% mechlib/core/styles.tex
% ZENTRALE, NAMENSRAUMGESCHUETZTE TIKZ-STILE
% =========================================================
% Diese Datei definiert ausschliesslich visuelle Grundstile. Die eigentliche
% Geometrie eines Elements gehoert in elements/ bzw. annotations/.
%
% Alle oeffentlichen Bibliotheksstile beginnen mit "mech/". Das reduziert
% Namenskollisionen mit dem Hauptdokument und macht im TikZ-Code sichtbar,
% welche Formatierung aus der mechlib stammt.
%
% Designregel:
%   - Farben werden aus core/colors.tex bezogen.
%   - Geometrien werden NICHT hier definiert.
%   - Linienstarken werden hier zentral gehalten, soweit sie elementuebergreifend
%     konsistent sein sollen.
% =========================================================
% -------------------------------------------------------------------------
% Visuelle Basisstile
% -------------------------------------------------------------------------
% Einzelne Pics duerfen diese Stile ueber ihre jeweiligen "style"-Keys ersetzen.
\tikzset{
  mech/body/.style={
    draw=black,
    fill=mechBodyColor,
    line width=.6pt
  },
  mech/support/.style={
    draw=black,
    fill=mechSupportColor,
    line width=.6pt
  },
  mech/connector/.style={
    draw=black,
    line width=.6pt
  },
  mech/rod/.style={
    draw=black,
    line width=.85pt
  },
  mech/load/.style={
    draw=mechLoadColor,
    text=mechLoadColor,
    line width=1pt
  },
  mech/load fill/.style={
    fill=mechLoadColor,
    fill opacity=.16,
    draw=none
  },
  mech/dimension/.style={
    {Latex[length=2mm]}-{Latex[length=2mm]},
    line width=.5pt
  },
  mech/axis/.style={
    draw=mechAxisColor,
    text=mechAxisColor,
    line width=.6pt
  },
  mech/label/.style={
    inner sep=1pt
  }
}

% =========================================================
% mechlib/core/defaults.tex
% GLOBALE STANDARDWERTE UND GEMEINSAME GEOMETRIEKONVENTIONEN
% =========================================================
% Hier stehen Werte, die von mehreren Modulen gemeinsam benutzt werden. Solche
% Konstanten verhindern, dass geometrisch zusammengehoerige Elemente lokal mit
% leicht unterschiedlichen Zahlen definiert werden.
%
% Beispiel: \mechConnectionOffset legt die vertikale Lage paralleler
% Anschlussanker fest. Masse und Einspannung verwenden exakt denselben Wert;
% dadurch liegen obere/mittlere/untere Anschluesse bei gleicher Mittelpunktlage
% garantiert auf gleicher Hoehe.
% =========================================================
% Globale Schriftvorgabe fuer jedes TikZ-Bild. "append style" ergaenzt
% vorhandene Bildstile, statt sie zu ueberschreiben.
\tikzset{
  every picture/.append style={font=\fontsize{10}{10}\selectfont},
}

% Standardwinkel fuer technische Bodenschraffuren.
\def\mechHatchAngle{45}

% Shared vertical distance of parallel connector anchors from the middle anchor.
% Wall and mass use this exact value so upper/middle/lower are aligned whenever
% their center points lie on the same horizontal line.
% Gemeinsamer vertikaler Versatz fuer parallele Anschlussanker.
\def\mechConnectionOffset{2.5mm}

% -------------------------------------------------------------------------
% Atomare mechanische Elemente.
% -------------------------------------------------------------------------
% =========================================================
% mechlib/elements/bodies.tex
% ATOMARE KOERPER: MASSE, PUNKTMASSE, WAGEN UND BALKEN
% =========================================================
% Die hier definierten TikZ-Pics stellen mechanische Grundkoerper bereit. Jedes
% Pic besitzt eine eigene pgfkeys-Familie unter /mech/... und kann dadurch mit
% benannten Optionen konfiguriert werden.
%
% Grundprinzip der API:
%   \pic (name) at (<Ort>) {mech/<element>={<key>=<wert>,...}};
%
% Benannte Pics exportieren Koordinaten bzw. werden bei node-basierten Koerpern
% auf den internen Shape-Knoten aliasiert. So koennen regulaere TikZ-Anker wie
% (m1.east), (m1.north west) usw. direkt benutzt werden.
%
% Oeffentliche Pics in dieser Datei:
%   mech/mass       rechteckiger Starrkoerper bzw. konzentrierte Masse
%   mech/point-mass kreisfoermige Punktmasse
%   mech/cart       Wagen mit zwei Raedern und optionalem Boden
%   mech/beam       Balken zwischen zwei beliebigen Punkten
% =========================================================
\pgfkeys{
  /mech/mass/.is family,
  /mech/mass,
  width/.initial=1.5cm,
  height/.initial=1cm,
  label/.initial=,
  style/.initial=mech/body,
}
\tikzset{
  pics/mech/mass/.style args={#1}{
    code={
      \pgfkeys{/mech/mass/.cd,#1}
      \node[
        \pgfkeysvalueof{/mech/mass/style},
        minimum width=\pgfkeysvalueof{/mech/mass/width},
        minimum height=\pgfkeysvalueof{/mech/mass/height},
        inner sep=0pt
      ] (-shape) at (0,0) {\pgfkeysvalueof{/mech/mass/label}};

      % Named body pics behave like regular TikZ nodes.
      % Example: \pic (m1) ... then (m1.east), (m1.south east), ...
      \edef\mechPicName{\pgfkeysvalueof{/tikz/name prefix}}
      \ifx\mechPicName\pgfutil@empty\else
        \pgfnodealias{\mechPicName}{\mechPicName-shape}
      \fi

      \coordinate (-center) at (-shape.center);

      % Three standardized connection levels on each vertical side.
      % The shared \mechConnectionOffset is also used by fixed supports.
      \coordinate (-west-upper)  at ($(-shape.west)+(0,\mechConnectionOffset)$);
      \coordinate (-west-middle) at (-shape.west);
      \coordinate (-west-lower)  at ($(-shape.west)+(0,-\mechConnectionOffset)$);

      \coordinate (-east-upper)  at ($(-shape.east)+(0,\mechConnectionOffset)$);
      \coordinate (-east-middle) at (-shape.east);
      \coordinate (-east-lower)  at ($(-shape.east)+(0,-\mechConnectionOffset)$);

      % Generic connection aliases
      \coordinate (-connection-upper)  at (-east-upper);
      \coordinate (-connection-middle) at (-east-middle);
      \coordinate (-connection-lower)  at (-east-lower);

      \coordinate (-north) at (-shape.north);
      \coordinate (-south) at (-shape.south);
    }
  },
  pics/mech/mass/.default={},
}
\pgfkeys{
  /mech/point-mass/.is family,
  /mech/point-mass,
  radius/.initial=2.2mm,
  label/.initial=,
  style/.initial=mech/body,
}
\tikzset{
  pics/mech/point-mass/.style args={#1}{
    code={
      \pgfkeys{/mech/point-mass/.cd,#1}
      \edef\mechPointMassRadius{\pgfkeysvalueof{/mech/point-mass/radius}}
      \edef\mechPointMassLabel{\pgfkeysvalueof{/mech/point-mass/label}}
      \edef\mechPointMassStyle{\pgfkeysvalueof{/mech/point-mass/style}}
      \node[
        \mechPointMassStyle,
        circle,
        minimum size=2\mechPointMassRadius,
        inner sep=0pt
      ] (-shape) at (0,0) {\mechPointMassLabel};

      \edef\mechPicName{\pgfkeysvalueof{/tikz/name prefix}}
      \ifx\mechPicName\pgfutil@empty\else
        \pgfnodealias{\mechPicName}{\mechPicName-shape}
      \fi

      \coordinate (-center) at (-shape.center);
    }
  },
  pics/mech/point-mass/.default={},
}
% -------------------------------------------------------------------------
% Wagen
% -------------------------------------------------------------------------
% Der Boolean "ground" wird als TeX-if gespeichert. Dadurch kann die optionale
% Bodendarstellung ohne Stringvergleich ein- bzw. ausgeschaltet werden.
\newif\ifmechCartGround
\pgfkeys{
  /mech/cart/.is family,
  /mech/cart,
  width/.initial=2cm,
  height/.initial=1cm,
  wheel radius/.initial=2mm,
  label/.initial=,
  ground/.is if=mechCartGround,
  ground/.default=true,
  ground=true,
}
\tikzset{
  pics/mech/cart/.style args={#1}{
    code={
      \pgfkeys{/mech/cart/.cd,#1}
      \edef\mechCartWidth{\pgfkeysvalueof{/mech/cart/width}}
      \edef\mechCartHeight{\pgfkeysvalueof{/mech/cart/height}}
      \edef\mechCartWheelRadius{\pgfkeysvalueof{/mech/cart/wheel radius}}
      \node[
        mech/body,
        minimum width=\mechCartWidth,
        minimum height=\mechCartHeight,
        inner sep=0pt
      ] (-shape) at (0,0) {\pgfkeysvalueof{/mech/cart/label}};

      \edef\mechPicName{\pgfkeysvalueof{/tikz/name prefix}}
      \ifx\mechPicName\pgfutil@empty\else
        \pgfnodealias{\mechPicName}{\mechPicName-shape}
      \fi

      \coordinate (-center) at (-shape.center);
      \coordinate (-west) at (-shape.west);
      \coordinate (-east) at (-shape.east);
      \coordinate (-north) at (-shape.north);
      \coordinate (-south) at (-shape.south);
      \draw[mech/body]
        (-.3\mechCartWidth,-.5\mechCartHeight-\mechCartWheelRadius)
        circle[radius=\mechCartWheelRadius];
      \draw[mech/body]
        (.3\mechCartWidth,-.5\mechCartHeight-\mechCartWheelRadius)
        circle[radius=\mechCartWheelRadius];
      \ifmechCartGround
        \draw
          (-.5\mechCartWidth,-.5\mechCartHeight-2\mechCartWheelRadius)
          --
          (.5\mechCartWidth,-.5\mechCartHeight-2\mechCartWheelRadius);
        \fill[
          pattern={Lines[angle=\mechHatchAngle,distance=1pt]},
          pattern color=black
        ]
          (-.5\mechCartWidth,-.5\mechCartHeight-2\mechCartWheelRadius)
          rectangle
          (.5\mechCartWidth,-.5\mechCartHeight-2\mechCartWheelRadius-3pt);
      \fi
    }
  },
  pics/mech/cart/.default={},
}
\pgfkeys{
  /mech/beam/.is family,
  /mech/beam,
  from/.initial={(0,0)},
  to/.initial={(1,0)},
  height/.initial=3mm,
  style/.initial=mech/body,
}
\tikzset{
  pics/mech/beam/.style args={#1}{
    code={
      \pgfkeys{/mech/beam/.cd,#1}
      \edef\mechBeamFrom{\pgfkeysvalueof{/mech/beam/from}}
      \edef\mechBeamTo{\pgfkeysvalueof{/mech/beam/to}}
      \edef\mechBeamHeight{\pgfkeysvalueof{/mech/beam/height}}
      \edef\mechBeamStyle{\pgfkeysvalueof{/mech/beam/style}}
      \path let
        \p1=\mechBeamFrom,
        \p2=\mechBeamTo,
        \n1={veclen(\x2-\x1,\y2-\y1)},
        \n2={(\y2-\y1)/\n1*\mechBeamHeight/2},
        \n3={-(\x2-\x1)/\n1*\mechBeamHeight/2}
      in
        coordinate (-A) at (\x1+\n2,\y1+\n3)
        coordinate (-B) at (\x2+\n2,\y2+\n3)
        coordinate (-C) at (\x2-\n2,\y2-\n3)
        coordinate (-D) at (\x1-\n2,\y1-\n3)
        coordinate (-start) at (\x1,\y1)
        coordinate (-end) at (\x2,\y2)
        coordinate (-center) at ({(\x1+\x2)/2},{(\y1+\y2)/2});
      \draw[\mechBeamStyle] (-A)--(-B)--(-C)--(-D)--cycle;
    }
  },
  pics/mech/beam/.default={},
}

% =========================================================
% mechlib/elements/supports.tex
% ATOMARE LAGER: FESTLAGER, LOSLAGER UND EINSPANNUNG
% =========================================================
% Diese Datei kapselt die Lagergeometrie. Fest- und Loslager teilen sich die
% Key-Familie /mech/support-common, damit Breite, Hoehe, Schraffur und Label
% konsistent parametrisiert werden koennen.
%
% Oeffentliche Pics:
%   mech/pinned-support Festlager / gelenkiges Lager
%   mech/roller-support Loslager mit sichtbarem Abstand zum Boden
%   mech/fixed-support  Einspannung, optional gedreht
%
% Aktuelle Fachwerk-Konvention:
%   - Die Dreiecksspitze liegt exakt im Mittelpunkt des Gelenks (0,0).
%   - Das Dreieck ist gegenueber der frueheren Version kompakt skaliert.
%   - Bodenlinie und Schraffur verwenden dieselben x-Grenzen wie die
%     Dreiecksgrundseite; beim Loslager wird die sichtbare Linienbreite explizit
%     beruecksichtigt.
%   - Das Gelenk selbst ist mit mechBodyColor gefuellt.
%
% Die pgfmathsetlengthmacro-Aufrufe sind absichtlich verwendet: LaTeX-Laengen
% duerfen nicht durch Textersatz wie ".5\\macro" skaliert werden, da dies bei
% expandierenden Dimensionsmakros zu sichtbarem Resttext (z. B. ".8mm") fuehren
% kann. Alle abgeleiteten Abmessungen werden deshalb numerisch berechnet.
% =========================================================
\pgfkeys{
  /mech/support-common/.is family,
  /mech/support-common,
  width/.initial=2.72mm,
  height/.initial=2.24mm,
  hatch angle/.initial=45,
  label/.initial=,
  label position/.initial=above,
}
\tikzset{
  % -----------------------------------------------------------------------
  % Festlager
  % -----------------------------------------------------------------------
  pics/mech/pinned-support/.style args={#1}{
    code={
      \pgfkeys{/mech/support-common/.cd,#1}
      \edef\mechSupportWidth{\pgfkeysvalueof{/mech/support-common/width}}
      \edef\mechSupportHeight{\pgfkeysvalueof{/mech/support-common/height}}
      \edef\mechSupportHatch{\pgfkeysvalueof{/mech/support-common/hatch angle}}
      \pgfmathsetlengthmacro{\mechSupportHalfWidth}{0.5*\mechSupportWidth}
      % Die Dreieckskontur ragt an der Grundseite durch die halbe Linienbreite
      % minimal ueber die geometrische Breite hinaus. Fuer identische sichtbare
      % Breite wird die Bodenlinie daher um 0.3pt je Seite erweitert.
      \pgfmathsetlengthmacro{\mechSupportGroundHalfWidth}{\mechSupportHalfWidth+0.3pt}
      \pgfmathsetlengthmacro{\mechSupportApexY}{0mm}
      \pgfmathsetlengthmacro{\mechSupportHatchDepth}{0.256*\mechSupportHeight}

      \coordinate (-connection) at (0,0);
      \coordinate (-center) at (0,0);

      \draw[mech/support]
        (-\mechSupportHalfWidth,-\mechSupportHeight)
        -- (0,\mechSupportApexY)
        -- (\mechSupportHalfWidth,-\mechSupportHeight)
        -- cycle;
      \draw[mech/body] (0,0) circle[radius=.6mm];
      \draw[black,line width=.6pt,line cap=butt]
        (-\mechSupportGroundHalfWidth,-\mechSupportHeight)
        -- (\mechSupportGroundHalfWidth,-\mechSupportHeight);
      \fill[
        pattern={Lines[angle=\mechSupportHatch,distance=1pt]},
        pattern color=black
      ]
        (-\mechSupportGroundHalfWidth,-\mechSupportHeight)
        rectangle
        (\mechSupportGroundHalfWidth,{-(\mechSupportHeight+\mechSupportHatchDepth)});

      \node[\pgfkeysvalueof{/mech/support-common/label position}]
        at (0,0) {\pgfkeysvalueof{/mech/support-common/label}};
    }
  },
  pics/mech/pinned-support/.default={},
  % -----------------------------------------------------------------------
  % Loslager
  % -----------------------------------------------------------------------
  pics/mech/roller-support/.style args={#1}{
    code={
      \pgfkeys{/mech/support-common/.cd,#1}
      \edef\mechSupportWidth{\pgfkeysvalueof{/mech/support-common/width}}
      \edef\mechSupportHeight{\pgfkeysvalueof{/mech/support-common/height}}
      \edef\mechSupportHatch{\pgfkeysvalueof{/mech/support-common/hatch angle}}
      \pgfmathsetlengthmacro{\mechSupportHalfWidth}{0.5*\mechSupportWidth}
      % The triangle outline extends by half its line width at the base corners.
      % Add 0.3pt so the roller ground has the same visible total width.
      \pgfmathsetlengthmacro{\mechSupportGroundHalfWidth}{\mechSupportHalfWidth+0.3pt}
      \pgfmathsetlengthmacro{\mechSupportApexY}{0mm}
      \pgfmathsetlengthmacro{\mechSupportRollerGap}{0.272*\mechSupportHeight}
      \pgfmathsetlengthmacro{\mechSupportHatchDepth}{0.256*\mechSupportHeight}
      \pgfmathsetlengthmacro{\mechSupportGroundY}{\mechSupportHeight+\mechSupportRollerGap}

      \coordinate (-connection) at (0,0);
      \coordinate (-center) at (0,0);

      \draw[mech/support]
        (-\mechSupportHalfWidth,-\mechSupportHeight)
        -- (0,\mechSupportApexY)
        -- (\mechSupportHalfWidth,-\mechSupportHeight)
        -- cycle;
      \draw[mech/body] (0,0) circle[radius=.6mm];
      \draw[black,line width=.6pt,line cap=butt]
        (-\mechSupportGroundHalfWidth,-\mechSupportGroundY)
        -- (\mechSupportGroundHalfWidth,-\mechSupportGroundY);
      \fill[
        pattern={Lines[angle=\mechSupportHatch,distance=1pt]},
        pattern color=black
      ]
        (-\mechSupportGroundHalfWidth,-\mechSupportGroundY)
        rectangle
        (\mechSupportGroundHalfWidth,{-(\mechSupportGroundY+\mechSupportHatchDepth)});

      \node[\pgfkeysvalueof{/mech/support-common/label position}]
        at (0,0) {\pgfkeysvalueof{/mech/support-common/label}};
    }
  },
  pics/mech/roller-support/.default={},
}
\pgfkeys{
  /mech/fixed-support/.is family,
  /mech/fixed-support,
  width/.initial=1cm,
  depth/.initial=1.5mm,
  hatch angle/.initial=45,
  angle/.initial=0,
  label/.initial=,
}
\tikzset{
  pics/mech/fixed-support/.style args={#1}{
    code={
      \pgfkeys{/mech/fixed-support/.cd,#1}
      \edef\mechFixedWidth{\pgfkeysvalueof{/mech/fixed-support/width}}
      \edef\mechFixedDepth{\pgfkeysvalueof{/mech/fixed-support/depth}}
      \edef\mechFixedHatch{\pgfkeysvalueof{/mech/fixed-support/hatch angle}}
      \edef\mechFixedAngle{\pgfkeysvalueof{/mech/fixed-support/angle}}

      \coordinate (-center) at (0,0);
      \coordinate (-connection-upper)  at (0,\mechConnectionOffset);
      \coordinate (-connection-middle) at (0,0);
      \coordinate (-connection-lower)  at (0,-\mechConnectionOffset);

      \begin{scope}[rotate=\mechFixedAngle]
        \draw[mech/body,fill=none] (-.5\mechFixedWidth,0)--(.5\mechFixedWidth,0);
        \fill[
          pattern={Lines[angle=\mechFixedHatch,distance={3pt/sqrt(2)}]},
          pattern color=black
        ]
          (-.5\mechFixedWidth,0)
          rectangle
          (.5\mechFixedWidth,\mechFixedDepth);
      \end{scope}

      \node[above] at (0,0) {\pgfkeysvalueof{/mech/fixed-support/label}};
    }
  },
  pics/mech/fixed-support/.default={},
}

% =========================================================
% mechlib/elements/connectors.tex
% ATOMARE VERBINDUNGSELEMENTE: FEDER UND DAEMPFER
% =========================================================
% Beide Elemente folgen derselben grundlegenden API: "from", "to" und "label".
% Die Endpunkte koennen beliebige TikZ-Koordinaten oder exportierte Anker aus
% anderen mechlib-Pics sein.
%
% Oeffentliche Pics:
%   mech/spring  Zickzack-Feder entlang der Verbindungslinie
%   mech/damper  Daempfer mit rechteckigem Daempferkoerper in Linienmitte
%
% Die Beschriftung wird "sloped" gezeichnet und folgt daher der Orientierung
% des Verbindungselements. Ueber "label side" kann sie ober- oder unterhalb
% der lokalen Verbindungslinie platziert werden.
% =========================================================
\pgfkeys{
  /mech/spring/.is family,
  /mech/spring,
  from/.initial={(0,0)},
  to/.initial={(1,0)},
  label/.initial=,
  label side/.initial=above,
  amplitude/.initial=2mm,
  segment length/.initial=4mm,
  style/.initial=mech/connector,
}
\tikzset{
  pics/mech/spring/.style args={#1}{
    code={
      \pgfkeys{/mech/spring/.cd,#1}
      \edef\mechSpringFrom{\pgfkeysvalueof{/mech/spring/from}}
      \edef\mechSpringTo{\pgfkeysvalueof{/mech/spring/to}}
      \draw[\pgfkeysvalueof{/mech/spring/style},
        decorate,
        decoration={zigzag,
          segment length=\pgfkeysvalueof{/mech/spring/segment length},
          amplitude=\pgfkeysvalueof{/mech/spring/amplitude}}]
        \mechSpringFrom -- \mechSpringTo
        node[midway,sloped,mech/label,outer sep=2mm,\pgfkeysvalueof{/mech/spring/label side}]
        {\pgfkeysvalueof{/mech/spring/label}};
    }
  },
  pics/mech/spring/.default={},
}
\pgfkeys{
  /mech/damper/.is family,
  /mech/damper,
  from/.initial={(0,0)},
  to/.initial={(1,0)},
  label/.initial=,
  label side/.initial=below,
  body width/.initial=8mm,
  body height/.initial=3mm,
  style/.initial=mech/connector,
}
\tikzset{
  pics/mech/damper/.style args={#1}{
    code={
      \pgfkeys{/mech/damper/.cd,#1}
      \edef\mechDamperFrom{\pgfkeysvalueof{/mech/damper/from}}
      \edef\mechDamperTo{\pgfkeysvalueof{/mech/damper/to}}
      \draw[\pgfkeysvalueof{/mech/damper/style}]
        \mechDamperFrom -- \mechDamperTo
        node[midway,sloped,draw,fill=white,
          minimum width=\pgfkeysvalueof{/mech/damper/body width},
          minimum height=\pgfkeysvalueof{/mech/damper/body height},
          inner sep=0pt] {};
      \path \mechDamperFrom -- \mechDamperTo
        node[midway,sloped,mech/label,outer sep=2mm,\pgfkeysvalueof{/mech/damper/label side}]
        {\pgfkeysvalueof{/mech/damper/label}};
    }
  },
  pics/mech/damper/.default={},
}

% =========================================================
% mechlib/elements/trusses.tex
% FACHWERKE: STAB UND GEMEINSAME GELENKDARSTELLUNG
% =========================================================
% Dieses Modul enthaelt die atomaren Fachwerkelemente. Die Topologie eines
% Fachwerks bleibt bewusst im Anwendungsdokument: Knoten werden als TikZ-
% Koordinaten definiert und durch mech/rod bzw. \mechJoints verbunden.
%
% Oeffentliche API:
%   mech/rod     Stab zwischen "from" und "to" mit optionalem Label
%   \mechJoints{A,B,C,...} zeichnet alle angegebenen Gelenke
%
% Darstellungsregeln:
%   - Staebe verwenden standardmaessig den Stil mech/rod: schwarz und duenn.
%   - Staebe werden auf dem Layer "mech background" gezeichnet.
%   - Stabnummern besitzen keine Fuellung (fill=none) und bleiben transparent.
%   - Gelenke liegen optisch vor den Staeben und sind mit mechBodyColor gefuellt.
%
% Beispiel:
%   \pic {mech/rod={from=(A),to=(B),label=1,label side=above}};
%   \mechJoints{A,B}
% =========================================================
% -------------------------------------------------------------------------
% Einzelner Fachwerkstab
% -------------------------------------------------------------------------
% "style" erlaubt lokale Abweichungen, ohne den globalen mech/rod-Stil zu
% veraendern. Die Standardbeschriftung liegt an der Stabmitte.
\pgfkeys{
  /mech/rod/.is family,
  /mech/rod,
  from/.initial={(0,0)},
  to/.initial={(1,0)},
  label/.initial=,
  label side/.initial=above,
  label style/.initial={mech/label,fill=none,text=black},
  style/.initial=mech/rod,
}
\tikzset{
  pics/mech/rod/.style args={#1}{
    code={
      \pgfkeys{/mech/rod/.cd,#1}
      \edef\mechRodFrom{\pgfkeysvalueof{/mech/rod/from}}
      \edef\mechRodTo{\pgfkeysvalueof{/mech/rod/to}}
      \coordinate (-start) at \mechRodFrom;
      \coordinate (-end) at \mechRodTo;
      \coordinate (-center) at ($(-start)!.5!(-end)$);
      \begin{pgfonlayer}{mech background}
        \draw[\pgfkeysvalueof{/mech/rod/style}] (-start)--(-end);
      \end{pgfonlayer}
      \node[\pgfkeysvalueof{/mech/rod/label style},\pgfkeysvalueof{/mech/rod/label side}]
        at (-center) {\pgfkeysvalueof{/mech/rod/label}};
    }
  },
  pics/mech/rod/.default={},
}
% -------------------------------------------------------------------------
% Gelenkliste
% -------------------------------------------------------------------------
% Das Makro erwartet eine kommaseparierte Liste bereits definierter TikZ-
% Koordinaten. Jedes Gelenk wird identisch gezeichnet. Die Gelenke sollten nach
% den Staeben aufgerufen werden, damit sie die Stabenden optisch ueberdecken.
\newcommand{\mechJoints}[1]{%
  \foreach \mechJoint in {#1}{%
    \draw[fill=mechBodyColor,draw=black,line width=.2mm]
      (\mechJoint) circle[radius=.6mm];%
  }%
}

% -------------------------------------------------------------------------
% Annotationen und mechanische Symbolgruppen.
% -------------------------------------------------------------------------
% =========================================================
% mechlib/annotations/loads.tex
% LASTANNOTATIONEN: KRAFT, MOMENT UND STRECKENLAST
% =========================================================
% Diese Datei stellt die Lastsymbole der mechlib bereit. Alle Lasten verwenden
% standardmaessig den zentralen Stil mech/load und damit mechLoadColor.
%
% Oeffentliche Pics:
%   mech/force            Einzelkraft als Winkel/Laenge oder from/to
%   mech/moment           Kreisbogensymbol im oder gegen den Uhrzeigersinn
%   mech/distributed-load konstante, dreieckige oder linear veraenderliche Last
%
% Besonderheit der Streckenlast:
% Die Grundlinie wird lokal auf die x-Achse gedreht. Dadurch kann die gesamte
% Geometrie in einem einfachen lokalen Koordinatensystem berechnet werden. Sehr
% kurze Pfeile, deren Laenge kleiner als ihre Pfeilspitze ist, werden bewusst
% unterdrueckt; so entstehen am Nullende einer Dreieckslast keine Artefakte.
% =========================================================
\pgfkeys{
  /mech/force/.is family,
  /mech/force,
  from/.initial={(0,0)},
  to/.initial=,
  angle/.initial=0,
  length/.initial=1.5cm,
  label/.initial=,
  label position/.initial=above,
  direction/.initial=away,
  style/.initial=mech/load,
}
\tikzset{
  pics/mech/force/.style args={#1}{
    code={
      \pgfkeys{/mech/force/.cd,#1}
      \edef\mechForceFrom{\pgfkeysvalueof{/mech/force/from}}
      \edef\mechForceTo{\pgfkeysvalueof{/mech/force/to}}
      \edef\mechForceDirection{\pgfkeysvalueof{/mech/force/direction}}
      \def\mechForceToward{toward}
      \ifx\mechForceTo\empty
        \ifx\mechForceDirection\mechForceToward
          \draw[\pgfkeysvalueof{/mech/force/style},-{Latex[length=2.5mm]}]
            (\pgfkeysvalueof{/mech/force/angle}:\pgfkeysvalueof{/mech/force/length}) -- (0,0)
            node[midway,\pgfkeysvalueof{/mech/force/label position}]
            {\pgfkeysvalueof{/mech/force/label}};
        \else
          \draw[\pgfkeysvalueof{/mech/force/style},-{Latex[length=2.5mm]}]
            (0,0) -- (\pgfkeysvalueof{/mech/force/angle}:\pgfkeysvalueof{/mech/force/length})
            node[midway,\pgfkeysvalueof{/mech/force/label position}]
            {\pgfkeysvalueof{/mech/force/label}};
        \fi
      \else
        \ifx\mechForceDirection\mechForceToward
          \draw[\pgfkeysvalueof{/mech/force/style},-{Latex[length=2.5mm]}]
            \mechForceTo -- \mechForceFrom
            node[midway,\pgfkeysvalueof{/mech/force/label position}]
            {\pgfkeysvalueof{/mech/force/label}};
        \else
          \draw[\pgfkeysvalueof{/mech/force/style},-{Latex[length=2.5mm]}]
            \mechForceFrom -- \mechForceTo
            node[midway,\pgfkeysvalueof{/mech/force/label position}]
            {\pgfkeysvalueof{/mech/force/label}};
        \fi
      \fi
    }
  },
  pics/mech/force/.default={},
}
\pgfkeys{
  /mech/moment/.is family,
  /mech/moment,
  radius/.initial=6mm,
  label/.initial=,
  direction/.is choice,
  direction/ccw/.code={\def\mechMomentStart{-60}\def\mechMomentEnd{220}},
  direction/cw/.code={\def\mechMomentStart{240}\def\mechMomentEnd{-40}},
  direction=ccw,
  style/.initial=mech/load,
}
\tikzset{
  pics/mech/moment/.style args={#1}{
    code={
      \pgfkeys{/mech/moment/.cd,direction=ccw,#1}
      \draw[\pgfkeysvalueof{/mech/moment/style},-{Latex[length=2.5mm]}]
        (\mechMomentStart:\pgfkeysvalueof{/mech/moment/radius})
        arc[start angle=\mechMomentStart,end angle=\mechMomentEnd,
            radius=\pgfkeysvalueof{/mech/moment/radius}];
      \node[mech/label,above] at (0,\pgfkeysvalueof{/mech/moment/radius})
        {\pgfkeysvalueof{/mech/moment/label}};
    }
  },
  pics/mech/moment/.default={},
}
\pgfkeys{
  /mech/distributed-load/.is family,
  /mech/distributed-load,
  from/.initial={(0,0)},
  to/.initial={(1,0)},
  start value/.initial=1,
  end value/.initial=1,
  height/.initial=8mm,
  offset/.initial=0mm,
  arrows/.initial=6,
  arrowhead length/.initial=2mm,
  label/.initial=,
  label position/.initial=above,
  style/.initial=mech/load,
}
\tikzset{
  pics/mech/distributed-load/.style args={#1}{
    code={
      \pgfkeys{/mech/distributed-load/.cd,#1}
      \edef\mechDLFrom{\pgfkeysvalueof{/mech/distributed-load/from}}
      \edef\mechDLTo{\pgfkeysvalueof{/mech/distributed-load/to}}
      \path let
        \p1=\mechDLFrom,\p2=\mechDLTo,
        \n1={veclen(\x2-\x1,\y2-\y1)},
        \n2={atan2(\y2-\y1,\x2-\x1)}
      in
        coordinate (-start) at (\x1,\y1)
        coordinate (-end) at (\x2,\y2)
        coordinate (-helper) at (\n2,\n1);
      \path (-helper); \pgfgetlastxy{\mechDLAngle}{\mechDLLength};
      \pgfmathsetmacro{\mechDLStartValue}{\pgfkeysvalueof{/mech/distributed-load/start value}}
      \pgfmathsetmacro{\mechDLEndValue}{\pgfkeysvalueof{/mech/distributed-load/end value}}
      \pgfmathsetmacro{\mechDLArrowCount}{\pgfkeysvalueof{/mech/distributed-load/arrows}}
      \begin{scope}[shift={(-start)},rotate=\mechDLAngle]
        \pgfmathsetlengthmacro{\mechDLOffset}{\pgfkeysvalueof{/mech/distributed-load/offset}}
        \pgfmathsetlengthmacro{\mechDLArrowheadLength}{\pgfkeysvalueof{/mech/distributed-load/arrowhead length}}
        \pgfmathsetlengthmacro{\mechDLStartArrowLength}{abs(\mechDLStartValue)*\pgfkeysvalueof{/mech/distributed-load/height}}

        \fill[mech/load fill]
          (0,\mechDLOffset)
          --(0,{\mechDLOffset+\mechDLStartValue*\pgfkeysvalueof{/mech/distributed-load/height}})
          --(\mechDLLength,{\mechDLOffset+\mechDLEndValue*\pgfkeysvalueof{/mech/distributed-load/height}})
          --(\mechDLLength,\mechDLOffset)
          --cycle;

        % Keep contour and first arrow continuous whenever the first arrow
        % is long enough to accommodate its arrowhead. Otherwise draw only
        % the contour and suppress the too-short arrow.
        \ifdim\mechDLStartArrowLength<\mechDLArrowheadLength
          \draw[\pgfkeysvalueof{/mech/distributed-load/style}]
            (\mechDLLength,{\mechDLOffset+\mechDLEndValue*\pgfkeysvalueof{/mech/distributed-load/height}})
            --(0,{\mechDLOffset+\mechDLStartValue*\pgfkeysvalueof{/mech/distributed-load/height}});
        \else
          \draw[
            \pgfkeysvalueof{/mech/distributed-load/style},
            -{Latex[length=\pgfkeysvalueof{/mech/distributed-load/arrowhead length}]}
          ]
            (\mechDLLength,{\mechDLOffset+\mechDLEndValue*\pgfkeysvalueof{/mech/distributed-load/height}})
            --(0,{\mechDLOffset+\mechDLStartValue*\pgfkeysvalueof{/mech/distributed-load/height}})
            --(0,\mechDLOffset);
        \fi

        \foreach \mechDLIndex in {1,...,\mechDLArrowCount}{
          \pgfmathsetmacro{\mechDLT}{\mechDLIndex/\mechDLArrowCount}
          \pgfmathsetmacro{\mechDLValue}{\mechDLStartValue+(\mechDLEndValue-\mechDLStartValue)*\mechDLT}
          \pgfmathsetlengthmacro{\mechDLArrowLength}{abs(\mechDLValue)*\pgfkeysvalueof{/mech/distributed-load/height}}

          \ifdim\mechDLArrowLength<\mechDLArrowheadLength
            % Suppress arrows that are shorter than their arrowhead.
          \else
            \draw[
              \pgfkeysvalueof{/mech/distributed-load/style},
              -{Latex[length=\pgfkeysvalueof{/mech/distributed-load/arrowhead length}]}
            ]
              ({\mechDLT*\mechDLLength},{\mechDLOffset+\mechDLValue*\pgfkeysvalueof{/mech/distributed-load/height}})
              --({\mechDLT*\mechDLLength},\mechDLOffset);
          \fi
        }

        \node[mech/label,\pgfkeysvalueof{/mech/distributed-load/label position}]
          at (\mechDLLength,{\mechDLOffset+\mechDLEndValue*\pgfkeysvalueof{/mech/distributed-load/height}})
          {\pgfkeysvalueof{/mech/distributed-load/label}};
      \end{scope}
    }
  },
  pics/mech/distributed-load/.default={},
}

% =========================================================
% mechlib/annotations/dimensions.tex
% BEMASSUNGEN: LINEARES MASS UND WINKELMASS
% =========================================================
% Lineare Bemassungen werden geometrisch aus zwei Endpunkten abgeleitet. Der
% Offset wirkt senkrecht auf die Verbindung von "from" nach "to"; damit
% funktioniert dasselbe Pic fuer horizontale, vertikale und schraege Masse.
%
% Oeffentliche Pics:
%   mech/dimension       lineare Bemassung mit Doppelpfeil
%   mech/angle-dimension Kreisbogen zwischen zwei Winkeln
%
% Interne Koordinaten -A und -B sind die senkrecht versetzten Endpunkte der
% Masslinie. Die Vektorrechnung erfolgt in TikZ "let"-Syntax.
% =========================================================
\pgfkeys{
  /mech/dimension/.is family,
  /mech/dimension,
  from/.initial={(0,0)},
  to/.initial={(1,0)},
  label/.initial={$l$},
  offset/.initial=3mm,
  style/.initial=mech/dimension,
}
\tikzset{
  pics/mech/dimension/.style args={#1}{
    code={
      \pgfkeys{/mech/dimension/.cd,#1}
      \edef\mechDimFrom{\pgfkeysvalueof{/mech/dimension/from}}
      \edef\mechDimTo{\pgfkeysvalueof{/mech/dimension/to}}
      \path let
        \p1=\mechDimFrom,\p2=\mechDimTo,
        \n1={veclen(\x2-\x1,\y2-\y1)},
        \n2={-(\y2-\y1)/\n1*\pgfkeysvalueof{/mech/dimension/offset}},
        \n3={(\x2-\x1)/\n1*\pgfkeysvalueof{/mech/dimension/offset}}
      in
        coordinate (-A) at (\x1+\n2,\y1+\n3)
        coordinate (-B) at (\x2+\n2,\y2+\n3);
      \draw[\pgfkeysvalueof{/mech/dimension/style}]
        (-A)--(-B)
        node[midway,fill=white,inner sep=1pt]
        {\pgfkeysvalueof{/mech/dimension/label}};
    }
  },
  pics/mech/dimension/.default={},
}
\pgfkeys{
  /mech/angle-dimension/.is family,
  /mech/angle-dimension,
  start angle/.initial=0,
  end angle/.initial=45,
  radius/.initial=7mm,
  label/.initial={$\alpha$},
}
\tikzset{
  pics/mech/angle-dimension/.style args={#1}{
    code={
      \pgfkeys{/mech/angle-dimension/.cd,#1}
      \draw
        (\pgfkeysvalueof{/mech/angle-dimension/start angle}:\pgfkeysvalueof{/mech/angle-dimension/radius})
        arc[start angle=\pgfkeysvalueof{/mech/angle-dimension/start angle},
            end angle=\pgfkeysvalueof{/mech/angle-dimension/end angle},
            radius=\pgfkeysvalueof{/mech/angle-dimension/radius}];
      \pgfmathsetmacro{\mechAngleMid}{.5*(\pgfkeysvalueof{/mech/angle-dimension/start angle}+\pgfkeysvalueof{/mech/angle-dimension/end angle})}
      \node at (\mechAngleMid:{.7*\pgfkeysvalueof{/mech/angle-dimension/radius}})
        {\pgfkeysvalueof{/mech/angle-dimension/label}};
    }
  },
  pics/mech/angle-dimension/.default={},
}

% =========================================================
% mechlib/annotations/coordinates.tex
% KOORDINATENANNOTATIONEN: EINZELACHSE UND 2D-KOORDINATENSYSTEM
% =========================================================
% Die Pics dieser Datei zeichnen mathematische Bezugsachsen im Stil mech/axis.
% Winkel werden in TikZ-Grad angegeben; Laengen sind regulaere TeX-Dimensionen.
%
% Oeffentliche Pics:
%   mech/axis              einzelne gerichtete Achse
%   mech/coordinate-system zwei unabhaengig orientierbare Achsen
%
% Beide Pics exportieren den lokalen Ursprung als "-origin". Bei einem benannten
% Pic (cs) kann dieser z. B. als (cs-origin) referenziert werden.
% =========================================================
\pgfkeys{
  /mech/axis/.is family,
  /mech/axis,
  label/.initial={$x$},
  angle/.initial=0,
  length/.initial=1cm,
  direction/.initial=positive,
}
\tikzset{
  pics/mech/axis/.style args={#1}{
    code={
      \pgfkeys{/mech/axis/.cd,#1}
      \edef\mechAxisDirection{\pgfkeysvalueof{/mech/axis/direction}}
      \def\mechAxisNegative{negative}
      \ifx\mechAxisDirection\mechAxisNegative
        \draw[mech/axis,{Latex[length=2mm]}-]
          (0,0)--(\pgfkeysvalueof{/mech/axis/angle}:\pgfkeysvalueof{/mech/axis/length})
          node[pos=1,anchor=west]{\pgfkeysvalueof{/mech/axis/label}};
      \else
        \draw[mech/axis,-{Latex[length=2mm]}]
          (0,0)--(\pgfkeysvalueof{/mech/axis/angle}:\pgfkeysvalueof{/mech/axis/length})
          node[pos=1,anchor=west]{\pgfkeysvalueof{/mech/axis/label}};
      \fi
      \coordinate (-origin) at (0,0);
    }
  },
  pics/mech/axis/.default={},
}
\pgfkeys{
  /mech/coordinate-system/.is family,
  /mech/coordinate-system,
  x label/.initial={$x$},
  y label/.initial={$y$},
  x angle/.initial=0,
  y angle/.initial=90,
  x length/.initial=1cm,
  y length/.initial=1cm,
}
\tikzset{
  pics/mech/coordinate-system/.style args={#1}{
    code={
      \pgfkeys{/mech/coordinate-system/.cd,#1}
      \draw[mech/axis,-{Latex[length=2mm]}]
        (0,0)--(\pgfkeysvalueof{/mech/coordinate-system/x angle}:\pgfkeysvalueof{/mech/coordinate-system/x length})
        node[pos=1,anchor=west]{\pgfkeysvalueof{/mech/coordinate-system/x label}};
      \draw[mech/axis,-{Latex[length=2mm]}]
        (0,0)--(\pgfkeysvalueof{/mech/coordinate-system/y angle}:\pgfkeysvalueof{/mech/coordinate-system/y length})
        node[pos=1,anchor=east]{\pgfkeysvalueof{/mech/coordinate-system/y label}};
      \coordinate (-origin) at (0,0);
    }
  },
  pics/mech/coordinate-system/.default={},
}

% =========================================================
% mechlib/annotations/mechanics.tex
% MECHANISCHE ANNOTATIONEN: SCHNITTGROESSEN UND LAGERREAKTIONEN
% =========================================================
% Die Pics kombinieren mehrere elementare Lastsymbole zu haeufig benoetigten
% mechanischen Symbolgruppen. Farben und Pfeilgestaltung folgen mech/load.
%
% Oeffentliche Pics:
%   mech/internal-forces  Normalkraft N, Querkraft Q und Moment M am Schnitt
%   mech/fixed-reactions  horizontale/vertikale Reaktion und Einspannmoment
%
% Bei mech/internal-forces steuert "side=right|left" die Vorzeichenrichtung
% aller drei Schnittgroessen konsistent. Das Moment wird dazu intern als
% mech/moment mit passender Drehrichtung wiederverwendet.
% =========================================================
\pgfkeys{
  /mech/internal-forces/.is family,
  /mech/internal-forces,
  normal/.initial={$N$},
  shear/.initial={$Q$},
  moment/.initial={$M$},
  length/.initial=7mm,
  side/.is choice,
  side/right/.code={\def\mechIFSign{1}},
  side/left/.code={\def\mechIFSign{-1}},
  side=right,
}
\tikzset{
  pics/mech/internal-forces/.style args={#1}{
    code={
      \pgfkeys{/mech/internal-forces/.cd,side=right,#1}
      \edef\mechIFLength{\pgfkeysvalueof{/mech/internal-forces/length}}
      \draw[mech/load,-{Latex[length=2mm]}]
        (0,0)--({\mechIFSign*\mechIFLength},0)
        node[pos=1,anchor=west]{\pgfkeysvalueof{/mech/internal-forces/normal}};
      \draw[mech/load,-{Latex[length=2mm]}]
        (0,0)--(0,{-.8*\mechIFSign*\mechIFLength})
        node[pos=1,anchor=west]{\pgfkeysvalueof{/mech/internal-forces/shear}};
      \ifnum\mechIFSign=1
        \pic {mech/moment={direction=ccw,radius=.55\mechIFLength,label=\pgfkeysvalueof{/mech/internal-forces/moment}}};
      \else
        \pic {mech/moment={direction=cw,radius=.55\mechIFLength,label=\pgfkeysvalueof{/mech/internal-forces/moment}}};
      \fi
    }
  },
  pics/mech/internal-forces/.default={},
}
\pgfkeys{
  /mech/fixed-reactions/.is family,
  /mech/fixed-reactions,
  x/.initial={$A_x$},
  y/.initial={$A_y$},
  moment/.initial={$M_A$},
  length/.initial=7mm,
}
\tikzset{
  pics/mech/fixed-reactions/.style args={#1}{
    code={
      \pgfkeys{/mech/fixed-reactions/.cd,#1}
      \edef\mechReactionLength{\pgfkeysvalueof{/mech/fixed-reactions/length}}
      \draw[mech/load,-{Latex[length=2mm]}] (-\mechReactionLength,0)--(0,0)
        node[pos=0,anchor=east]{\pgfkeysvalueof{/mech/fixed-reactions/x}};
      \draw[mech/load,-{Latex[length=2mm]}] (0,.7\mechReactionLength)--(0,-.3\mechReactionLength)
        node[pos=0,anchor=south]{\pgfkeysvalueof{/mech/fixed-reactions/y}};
      \pic {mech/moment={direction=ccw,radius=.45\mechReactionLength,label=\pgfkeysvalueof{/mech/fixed-reactions/moment}}};
    }
  },
  pics/mech/fixed-reactions/.default={},
}


%
% Liegt die Bibliothek an einem anderen Ort, kann vor dem Laden die Variable
% \mechlibpath gesetzt werden:
%   \newcommand{\mechlibpath}{../../../support/mechlib}
%   % =========================================================
% mechlib/mechlib.tex
% ZENTRALER LOADER / OEFFENTLICHER EINSTIEGSPUNKT
% =========================================================
% Diese Datei ist die einzige Datei, die ein Anwender normalerweise direkt
% einbindet. Sie laedt die benoetigten LaTeX-/TikZ-Abhaengigkeiten und danach
% alle Teilmodule der Bibliothek in einer definierten Reihenfolge.
%
% Typische Verwendung bei einem Ordner "mechlib" neben dem Hauptdokument:
%   % =========================================================
% mechlib/mechlib.tex
% ZENTRALER LOADER / OEFFENTLICHER EINSTIEGSPUNKT
% =========================================================
% Diese Datei ist die einzige Datei, die ein Anwender normalerweise direkt
% einbindet. Sie laedt die benoetigten LaTeX-/TikZ-Abhaengigkeiten und danach
% alle Teilmodule der Bibliothek in einer definierten Reihenfolge.
%
% Typische Verwendung bei einem Ordner "mechlib" neben dem Hauptdokument:
%   % =========================================================
% mechlib/mechlib.tex
% ZENTRALER LOADER / OEFFENTLICHER EINSTIEGSPUNKT
% =========================================================
% Diese Datei ist die einzige Datei, die ein Anwender normalerweise direkt
% einbindet. Sie laedt die benoetigten LaTeX-/TikZ-Abhaengigkeiten und danach
% alle Teilmodule der Bibliothek in einer definierten Reihenfolge.
%
% Typische Verwendung bei einem Ordner "mechlib" neben dem Hauptdokument:
%   \input{mechlib/mechlib}
%
% Liegt die Bibliothek an einem anderen Ort, kann vor dem Laden die Variable
% \mechlibpath gesetzt werden:
%   \newcommand{\mechlibpath}{../../../support/mechlib}
%   \input{\mechlibpath/mechlib}
%
% Lade-Reihenfolge:
%   1. core/        Farben, Layer, Stile und globale Geometriekonventionen
%   2. elements/    mechanische Grundelemente (Koerper, Lager, Verbinder, Staebe)
%   3. annotations/ Lasten, Masse, Koordinaten und mechanische Annotationen
%
% WICHTIG:
% Die Modulreihenfolge ist absichtlich festgelegt. Elemente greifen auf Farben,
% Layer und Stile aus core/ zurueck; Annotationen koennen wiederum die bereits
% definierten Grundstile und Pics verwenden.
% =========================================================
% -------------------------------------------------------------------------
% Externe Pakete und TikZ-Bibliotheken
% -------------------------------------------------------------------------
% \RequirePackage ist bewusst statt \usepackage gewaehlt, da mechlib.tex selbst
% wie eine Bibliotheks-/Paketdatei verwendet wird.
\RequirePackage{tikz}
\RequirePackage{amsmath}
\RequirePackage{siunitx}
\RequirePackage{xcolor}
\usetikzlibrary{
  arrows.meta,
  calc,
  patterns,
  patterns.meta,
  intersections,
  angles,
  quotes,
  decorations.pathmorphing,
  backgrounds
}
% Base path: override before loading when mechlib is stored elsewhere.
% -------------------------------------------------------------------------
% Aufloesung des Bibliothekspfades
% -------------------------------------------------------------------------
% \providecommand respektiert eine bereits im Hauptdokument gesetzte Definition.
\providecommand{\mechlibpath}{mechlib}
% -------------------------------------------------------------------------
% Kernmodule: muessen vor allen Elementen geladen werden.
% -------------------------------------------------------------------------
\input{\mechlibpath/core/colors}
\input{\mechlibpath/core/layers}
\input{\mechlibpath/core/styles}
\input{\mechlibpath/core/defaults}
% -------------------------------------------------------------------------
% Atomare mechanische Elemente.
% -------------------------------------------------------------------------
\input{\mechlibpath/elements/bodies}
\input{\mechlibpath/elements/supports}
\input{\mechlibpath/elements/connectors}
\input{\mechlibpath/elements/trusses}
% -------------------------------------------------------------------------
% Annotationen und mechanische Symbolgruppen.
% -------------------------------------------------------------------------
\input{\mechlibpath/annotations/loads}
\input{\mechlibpath/annotations/dimensions}
\input{\mechlibpath/annotations/coordinates}
\input{\mechlibpath/annotations/mechanics}

%
% Liegt die Bibliothek an einem anderen Ort, kann vor dem Laden die Variable
% \mechlibpath gesetzt werden:
%   \newcommand{\mechlibpath}{../../../support/mechlib}
%   % =========================================================
% mechlib/mechlib.tex
% ZENTRALER LOADER / OEFFENTLICHER EINSTIEGSPUNKT
% =========================================================
% Diese Datei ist die einzige Datei, die ein Anwender normalerweise direkt
% einbindet. Sie laedt die benoetigten LaTeX-/TikZ-Abhaengigkeiten und danach
% alle Teilmodule der Bibliothek in einer definierten Reihenfolge.
%
% Typische Verwendung bei einem Ordner "mechlib" neben dem Hauptdokument:
%   \input{mechlib/mechlib}
%
% Liegt die Bibliothek an einem anderen Ort, kann vor dem Laden die Variable
% \mechlibpath gesetzt werden:
%   \newcommand{\mechlibpath}{../../../support/mechlib}
%   \input{\mechlibpath/mechlib}
%
% Lade-Reihenfolge:
%   1. core/        Farben, Layer, Stile und globale Geometriekonventionen
%   2. elements/    mechanische Grundelemente (Koerper, Lager, Verbinder, Staebe)
%   3. annotations/ Lasten, Masse, Koordinaten und mechanische Annotationen
%
% WICHTIG:
% Die Modulreihenfolge ist absichtlich festgelegt. Elemente greifen auf Farben,
% Layer und Stile aus core/ zurueck; Annotationen koennen wiederum die bereits
% definierten Grundstile und Pics verwenden.
% =========================================================
% -------------------------------------------------------------------------
% Externe Pakete und TikZ-Bibliotheken
% -------------------------------------------------------------------------
% \RequirePackage ist bewusst statt \usepackage gewaehlt, da mechlib.tex selbst
% wie eine Bibliotheks-/Paketdatei verwendet wird.
\RequirePackage{tikz}
\RequirePackage{amsmath}
\RequirePackage{siunitx}
\RequirePackage{xcolor}
\usetikzlibrary{
  arrows.meta,
  calc,
  patterns,
  patterns.meta,
  intersections,
  angles,
  quotes,
  decorations.pathmorphing,
  backgrounds
}
% Base path: override before loading when mechlib is stored elsewhere.
% -------------------------------------------------------------------------
% Aufloesung des Bibliothekspfades
% -------------------------------------------------------------------------
% \providecommand respektiert eine bereits im Hauptdokument gesetzte Definition.
\providecommand{\mechlibpath}{mechlib}
% -------------------------------------------------------------------------
% Kernmodule: muessen vor allen Elementen geladen werden.
% -------------------------------------------------------------------------
\input{\mechlibpath/core/colors}
\input{\mechlibpath/core/layers}
\input{\mechlibpath/core/styles}
\input{\mechlibpath/core/defaults}
% -------------------------------------------------------------------------
% Atomare mechanische Elemente.
% -------------------------------------------------------------------------
\input{\mechlibpath/elements/bodies}
\input{\mechlibpath/elements/supports}
\input{\mechlibpath/elements/connectors}
\input{\mechlibpath/elements/trusses}
% -------------------------------------------------------------------------
% Annotationen und mechanische Symbolgruppen.
% -------------------------------------------------------------------------
\input{\mechlibpath/annotations/loads}
\input{\mechlibpath/annotations/dimensions}
\input{\mechlibpath/annotations/coordinates}
\input{\mechlibpath/annotations/mechanics}

%
% Lade-Reihenfolge:
%   1. core/        Farben, Layer, Stile und globale Geometriekonventionen
%   2. elements/    mechanische Grundelemente (Koerper, Lager, Verbinder, Staebe)
%   3. annotations/ Lasten, Masse, Koordinaten und mechanische Annotationen
%
% WICHTIG:
% Die Modulreihenfolge ist absichtlich festgelegt. Elemente greifen auf Farben,
% Layer und Stile aus core/ zurueck; Annotationen koennen wiederum die bereits
% definierten Grundstile und Pics verwenden.
% =========================================================
% -------------------------------------------------------------------------
% Externe Pakete und TikZ-Bibliotheken
% -------------------------------------------------------------------------
% \RequirePackage ist bewusst statt \usepackage gewaehlt, da mechlib.tex selbst
% wie eine Bibliotheks-/Paketdatei verwendet wird.
\RequirePackage{tikz}
\RequirePackage{amsmath}
\RequirePackage{siunitx}
\RequirePackage{xcolor}
\usetikzlibrary{
  arrows.meta,
  calc,
  patterns,
  patterns.meta,
  intersections,
  angles,
  quotes,
  decorations.pathmorphing,
  backgrounds
}
% Base path: override before loading when mechlib is stored elsewhere.
% -------------------------------------------------------------------------
% Aufloesung des Bibliothekspfades
% -------------------------------------------------------------------------
% \providecommand respektiert eine bereits im Hauptdokument gesetzte Definition.
\providecommand{\mechlibpath}{mechlib}
% -------------------------------------------------------------------------
% Kernmodule: muessen vor allen Elementen geladen werden.
% -------------------------------------------------------------------------
% =========================================================
% mechlib/core/colors.tex
% ZENTRALE FARBPALETTE
% =========================================================
% Alle bibliotheksweiten Farben werden ausschliesslich hier definiert. Dadurch
% kann der visuelle Charakter der gesamten mechlib an einer Stelle angepasst
% werden, ohne einzelne Elemente zu veraendern.
%
% Konvention:
%   mechBodyColor    Fuellung allgemeiner Koerper und Fachwerkgelenke
%   mechSupportColor Fuellung von Lagerkoerpern
%   mechLoadColor    Kraefte, Momente und Streckenlasten
%   mechAxisColor    Achsen und Koordinatensysteme
%   mechRodColor     historische Stabfarbe; aktuelle Fachwerkstaebe sind schwarz
%
% Die Namen sind Teil der internen visuellen API und sollten bevorzugt statt
% harter RGB-Werte in den Elementdateien verwendet werden.
% =========================================================
% -------------------------------------------------------------------------
% Farbdefinitionen
% -------------------------------------------------------------------------
% RGB wird bewusst explizit verwendet, damit die Darstellung unabhaengig von
% Dokumentklassen und externen Farbmodellen reproduzierbar bleibt.
\definecolor{mechBodyColor}{RGB}{180,180,180}
\definecolor{mechSupportColor}{RGB}{205,205,205}
\definecolor{mechLoadColor}{RGB}{178,25,25}
\definecolor{mechAxisColor}{RGB}{6,13,141}
\definecolor{mechRodColor}{RGB}{130,130,130}

% =========================================================
% mechlib/core/layers.tex
% ZEICHENEBENEN / PGF-LAYER
% =========================================================
% Fachwerkstaebe und andere hinterlegte Geometrien koennen auf dem Layer
% "mech background" gezeichnet werden. Der normale TikZ-Layer "main" liegt
% darueber. Dadurch koennen z. B. gefuellte Gelenke sauber vor den Staeben
% erscheinen, ohne dass die Zeichenreihenfolge im Anwendungsdokument aufwendig
% kontrolliert werden muss.
% =========================================================
% Zuerst wird der benannte Hintergrundlayer deklariert. Anschliessend wird die
% globale Layerreihenfolge explizit gesetzt: Hintergrund vor Hauptlayer.
\pgfdeclarelayer{mech background}
\pgfsetlayers{mech background,main}

% =========================================================
% mechlib/core/styles.tex
% ZENTRALE, NAMENSRAUMGESCHUETZTE TIKZ-STILE
% =========================================================
% Diese Datei definiert ausschliesslich visuelle Grundstile. Die eigentliche
% Geometrie eines Elements gehoert in elements/ bzw. annotations/.
%
% Alle oeffentlichen Bibliotheksstile beginnen mit "mech/". Das reduziert
% Namenskollisionen mit dem Hauptdokument und macht im TikZ-Code sichtbar,
% welche Formatierung aus der mechlib stammt.
%
% Designregel:
%   - Farben werden aus core/colors.tex bezogen.
%   - Geometrien werden NICHT hier definiert.
%   - Linienstarken werden hier zentral gehalten, soweit sie elementuebergreifend
%     konsistent sein sollen.
% =========================================================
% -------------------------------------------------------------------------
% Visuelle Basisstile
% -------------------------------------------------------------------------
% Einzelne Pics duerfen diese Stile ueber ihre jeweiligen "style"-Keys ersetzen.
\tikzset{
  mech/body/.style={
    draw=black,
    fill=mechBodyColor,
    line width=.6pt
  },
  mech/support/.style={
    draw=black,
    fill=mechSupportColor,
    line width=.6pt
  },
  mech/connector/.style={
    draw=black,
    line width=.6pt
  },
  mech/rod/.style={
    draw=black,
    line width=.85pt
  },
  mech/load/.style={
    draw=mechLoadColor,
    text=mechLoadColor,
    line width=1pt
  },
  mech/load fill/.style={
    fill=mechLoadColor,
    fill opacity=.16,
    draw=none
  },
  mech/dimension/.style={
    {Latex[length=2mm]}-{Latex[length=2mm]},
    line width=.5pt
  },
  mech/axis/.style={
    draw=mechAxisColor,
    text=mechAxisColor,
    line width=.6pt
  },
  mech/label/.style={
    inner sep=1pt
  }
}

% =========================================================
% mechlib/core/defaults.tex
% GLOBALE STANDARDWERTE UND GEMEINSAME GEOMETRIEKONVENTIONEN
% =========================================================
% Hier stehen Werte, die von mehreren Modulen gemeinsam benutzt werden. Solche
% Konstanten verhindern, dass geometrisch zusammengehoerige Elemente lokal mit
% leicht unterschiedlichen Zahlen definiert werden.
%
% Beispiel: \mechConnectionOffset legt die vertikale Lage paralleler
% Anschlussanker fest. Masse und Einspannung verwenden exakt denselben Wert;
% dadurch liegen obere/mittlere/untere Anschluesse bei gleicher Mittelpunktlage
% garantiert auf gleicher Hoehe.
% =========================================================
% Globale Schriftvorgabe fuer jedes TikZ-Bild. "append style" ergaenzt
% vorhandene Bildstile, statt sie zu ueberschreiben.
\tikzset{
  every picture/.append style={font=\fontsize{10}{10}\selectfont},
}

% Standardwinkel fuer technische Bodenschraffuren.
\def\mechHatchAngle{45}

% Shared vertical distance of parallel connector anchors from the middle anchor.
% Wall and mass use this exact value so upper/middle/lower are aligned whenever
% their center points lie on the same horizontal line.
% Gemeinsamer vertikaler Versatz fuer parallele Anschlussanker.
\def\mechConnectionOffset{2.5mm}

% -------------------------------------------------------------------------
% Atomare mechanische Elemente.
% -------------------------------------------------------------------------
% =========================================================
% mechlib/elements/bodies.tex
% ATOMARE KOERPER: MASSE, PUNKTMASSE, WAGEN UND BALKEN
% =========================================================
% Die hier definierten TikZ-Pics stellen mechanische Grundkoerper bereit. Jedes
% Pic besitzt eine eigene pgfkeys-Familie unter /mech/... und kann dadurch mit
% benannten Optionen konfiguriert werden.
%
% Grundprinzip der API:
%   \pic (name) at (<Ort>) {mech/<element>={<key>=<wert>,...}};
%
% Benannte Pics exportieren Koordinaten bzw. werden bei node-basierten Koerpern
% auf den internen Shape-Knoten aliasiert. So koennen regulaere TikZ-Anker wie
% (m1.east), (m1.north west) usw. direkt benutzt werden.
%
% Oeffentliche Pics in dieser Datei:
%   mech/mass       rechteckiger Starrkoerper bzw. konzentrierte Masse
%   mech/point-mass kreisfoermige Punktmasse
%   mech/cart       Wagen mit zwei Raedern und optionalem Boden
%   mech/beam       Balken zwischen zwei beliebigen Punkten
% =========================================================
\pgfkeys{
  /mech/mass/.is family,
  /mech/mass,
  width/.initial=1.5cm,
  height/.initial=1cm,
  label/.initial=,
  style/.initial=mech/body,
}
\tikzset{
  pics/mech/mass/.style args={#1}{
    code={
      \pgfkeys{/mech/mass/.cd,#1}
      \node[
        \pgfkeysvalueof{/mech/mass/style},
        minimum width=\pgfkeysvalueof{/mech/mass/width},
        minimum height=\pgfkeysvalueof{/mech/mass/height},
        inner sep=0pt
      ] (-shape) at (0,0) {\pgfkeysvalueof{/mech/mass/label}};

      % Named body pics behave like regular TikZ nodes.
      % Example: \pic (m1) ... then (m1.east), (m1.south east), ...
      \edef\mechPicName{\pgfkeysvalueof{/tikz/name prefix}}
      \ifx\mechPicName\pgfutil@empty\else
        \pgfnodealias{\mechPicName}{\mechPicName-shape}
      \fi

      \coordinate (-center) at (-shape.center);

      % Three standardized connection levels on each vertical side.
      % The shared \mechConnectionOffset is also used by fixed supports.
      \coordinate (-west-upper)  at ($(-shape.west)+(0,\mechConnectionOffset)$);
      \coordinate (-west-middle) at (-shape.west);
      \coordinate (-west-lower)  at ($(-shape.west)+(0,-\mechConnectionOffset)$);

      \coordinate (-east-upper)  at ($(-shape.east)+(0,\mechConnectionOffset)$);
      \coordinate (-east-middle) at (-shape.east);
      \coordinate (-east-lower)  at ($(-shape.east)+(0,-\mechConnectionOffset)$);

      % Generic connection aliases
      \coordinate (-connection-upper)  at (-east-upper);
      \coordinate (-connection-middle) at (-east-middle);
      \coordinate (-connection-lower)  at (-east-lower);

      \coordinate (-north) at (-shape.north);
      \coordinate (-south) at (-shape.south);
    }
  },
  pics/mech/mass/.default={},
}
\pgfkeys{
  /mech/point-mass/.is family,
  /mech/point-mass,
  radius/.initial=2.2mm,
  label/.initial=,
  style/.initial=mech/body,
}
\tikzset{
  pics/mech/point-mass/.style args={#1}{
    code={
      \pgfkeys{/mech/point-mass/.cd,#1}
      \edef\mechPointMassRadius{\pgfkeysvalueof{/mech/point-mass/radius}}
      \edef\mechPointMassLabel{\pgfkeysvalueof{/mech/point-mass/label}}
      \edef\mechPointMassStyle{\pgfkeysvalueof{/mech/point-mass/style}}
      \node[
        \mechPointMassStyle,
        circle,
        minimum size=2\mechPointMassRadius,
        inner sep=0pt
      ] (-shape) at (0,0) {\mechPointMassLabel};

      \edef\mechPicName{\pgfkeysvalueof{/tikz/name prefix}}
      \ifx\mechPicName\pgfutil@empty\else
        \pgfnodealias{\mechPicName}{\mechPicName-shape}
      \fi

      \coordinate (-center) at (-shape.center);
    }
  },
  pics/mech/point-mass/.default={},
}
% -------------------------------------------------------------------------
% Wagen
% -------------------------------------------------------------------------
% Der Boolean "ground" wird als TeX-if gespeichert. Dadurch kann die optionale
% Bodendarstellung ohne Stringvergleich ein- bzw. ausgeschaltet werden.
\newif\ifmechCartGround
\pgfkeys{
  /mech/cart/.is family,
  /mech/cart,
  width/.initial=2cm,
  height/.initial=1cm,
  wheel radius/.initial=2mm,
  label/.initial=,
  ground/.is if=mechCartGround,
  ground/.default=true,
  ground=true,
}
\tikzset{
  pics/mech/cart/.style args={#1}{
    code={
      \pgfkeys{/mech/cart/.cd,#1}
      \edef\mechCartWidth{\pgfkeysvalueof{/mech/cart/width}}
      \edef\mechCartHeight{\pgfkeysvalueof{/mech/cart/height}}
      \edef\mechCartWheelRadius{\pgfkeysvalueof{/mech/cart/wheel radius}}
      \node[
        mech/body,
        minimum width=\mechCartWidth,
        minimum height=\mechCartHeight,
        inner sep=0pt
      ] (-shape) at (0,0) {\pgfkeysvalueof{/mech/cart/label}};

      \edef\mechPicName{\pgfkeysvalueof{/tikz/name prefix}}
      \ifx\mechPicName\pgfutil@empty\else
        \pgfnodealias{\mechPicName}{\mechPicName-shape}
      \fi

      \coordinate (-center) at (-shape.center);
      \coordinate (-west) at (-shape.west);
      \coordinate (-east) at (-shape.east);
      \coordinate (-north) at (-shape.north);
      \coordinate (-south) at (-shape.south);
      \draw[mech/body]
        (-.3\mechCartWidth,-.5\mechCartHeight-\mechCartWheelRadius)
        circle[radius=\mechCartWheelRadius];
      \draw[mech/body]
        (.3\mechCartWidth,-.5\mechCartHeight-\mechCartWheelRadius)
        circle[radius=\mechCartWheelRadius];
      \ifmechCartGround
        \draw
          (-.5\mechCartWidth,-.5\mechCartHeight-2\mechCartWheelRadius)
          --
          (.5\mechCartWidth,-.5\mechCartHeight-2\mechCartWheelRadius);
        \fill[
          pattern={Lines[angle=\mechHatchAngle,distance=1pt]},
          pattern color=black
        ]
          (-.5\mechCartWidth,-.5\mechCartHeight-2\mechCartWheelRadius)
          rectangle
          (.5\mechCartWidth,-.5\mechCartHeight-2\mechCartWheelRadius-3pt);
      \fi
    }
  },
  pics/mech/cart/.default={},
}
\pgfkeys{
  /mech/beam/.is family,
  /mech/beam,
  from/.initial={(0,0)},
  to/.initial={(1,0)},
  height/.initial=3mm,
  style/.initial=mech/body,
}
\tikzset{
  pics/mech/beam/.style args={#1}{
    code={
      \pgfkeys{/mech/beam/.cd,#1}
      \edef\mechBeamFrom{\pgfkeysvalueof{/mech/beam/from}}
      \edef\mechBeamTo{\pgfkeysvalueof{/mech/beam/to}}
      \edef\mechBeamHeight{\pgfkeysvalueof{/mech/beam/height}}
      \edef\mechBeamStyle{\pgfkeysvalueof{/mech/beam/style}}
      \path let
        \p1=\mechBeamFrom,
        \p2=\mechBeamTo,
        \n1={veclen(\x2-\x1,\y2-\y1)},
        \n2={(\y2-\y1)/\n1*\mechBeamHeight/2},
        \n3={-(\x2-\x1)/\n1*\mechBeamHeight/2}
      in
        coordinate (-A) at (\x1+\n2,\y1+\n3)
        coordinate (-B) at (\x2+\n2,\y2+\n3)
        coordinate (-C) at (\x2-\n2,\y2-\n3)
        coordinate (-D) at (\x1-\n2,\y1-\n3)
        coordinate (-start) at (\x1,\y1)
        coordinate (-end) at (\x2,\y2)
        coordinate (-center) at ({(\x1+\x2)/2},{(\y1+\y2)/2});
      \draw[\mechBeamStyle] (-A)--(-B)--(-C)--(-D)--cycle;
    }
  },
  pics/mech/beam/.default={},
}

% =========================================================
% mechlib/elements/supports.tex
% ATOMARE LAGER: FESTLAGER, LOSLAGER UND EINSPANNUNG
% =========================================================
% Diese Datei kapselt die Lagergeometrie. Fest- und Loslager teilen sich die
% Key-Familie /mech/support-common, damit Breite, Hoehe, Schraffur und Label
% konsistent parametrisiert werden koennen.
%
% Oeffentliche Pics:
%   mech/pinned-support Festlager / gelenkiges Lager
%   mech/roller-support Loslager mit sichtbarem Abstand zum Boden
%   mech/fixed-support  Einspannung, optional gedreht
%
% Aktuelle Fachwerk-Konvention:
%   - Die Dreiecksspitze liegt exakt im Mittelpunkt des Gelenks (0,0).
%   - Das Dreieck ist gegenueber der frueheren Version kompakt skaliert.
%   - Bodenlinie und Schraffur verwenden dieselben x-Grenzen wie die
%     Dreiecksgrundseite; beim Loslager wird die sichtbare Linienbreite explizit
%     beruecksichtigt.
%   - Das Gelenk selbst ist mit mechBodyColor gefuellt.
%
% Die pgfmathsetlengthmacro-Aufrufe sind absichtlich verwendet: LaTeX-Laengen
% duerfen nicht durch Textersatz wie ".5\\macro" skaliert werden, da dies bei
% expandierenden Dimensionsmakros zu sichtbarem Resttext (z. B. ".8mm") fuehren
% kann. Alle abgeleiteten Abmessungen werden deshalb numerisch berechnet.
% =========================================================
\pgfkeys{
  /mech/support-common/.is family,
  /mech/support-common,
  width/.initial=2.72mm,
  height/.initial=2.24mm,
  hatch angle/.initial=45,
  label/.initial=,
  label position/.initial=above,
}
\tikzset{
  % -----------------------------------------------------------------------
  % Festlager
  % -----------------------------------------------------------------------
  pics/mech/pinned-support/.style args={#1}{
    code={
      \pgfkeys{/mech/support-common/.cd,#1}
      \edef\mechSupportWidth{\pgfkeysvalueof{/mech/support-common/width}}
      \edef\mechSupportHeight{\pgfkeysvalueof{/mech/support-common/height}}
      \edef\mechSupportHatch{\pgfkeysvalueof{/mech/support-common/hatch angle}}
      \pgfmathsetlengthmacro{\mechSupportHalfWidth}{0.5*\mechSupportWidth}
      % Die Dreieckskontur ragt an der Grundseite durch die halbe Linienbreite
      % minimal ueber die geometrische Breite hinaus. Fuer identische sichtbare
      % Breite wird die Bodenlinie daher um 0.3pt je Seite erweitert.
      \pgfmathsetlengthmacro{\mechSupportGroundHalfWidth}{\mechSupportHalfWidth+0.3pt}
      \pgfmathsetlengthmacro{\mechSupportApexY}{0mm}
      \pgfmathsetlengthmacro{\mechSupportHatchDepth}{0.256*\mechSupportHeight}

      \coordinate (-connection) at (0,0);
      \coordinate (-center) at (0,0);

      \draw[mech/support]
        (-\mechSupportHalfWidth,-\mechSupportHeight)
        -- (0,\mechSupportApexY)
        -- (\mechSupportHalfWidth,-\mechSupportHeight)
        -- cycle;
      \draw[mech/body] (0,0) circle[radius=.6mm];
      \draw[black,line width=.6pt,line cap=butt]
        (-\mechSupportGroundHalfWidth,-\mechSupportHeight)
        -- (\mechSupportGroundHalfWidth,-\mechSupportHeight);
      \fill[
        pattern={Lines[angle=\mechSupportHatch,distance=1pt]},
        pattern color=black
      ]
        (-\mechSupportGroundHalfWidth,-\mechSupportHeight)
        rectangle
        (\mechSupportGroundHalfWidth,{-(\mechSupportHeight+\mechSupportHatchDepth)});

      \node[\pgfkeysvalueof{/mech/support-common/label position}]
        at (0,0) {\pgfkeysvalueof{/mech/support-common/label}};
    }
  },
  pics/mech/pinned-support/.default={},
  % -----------------------------------------------------------------------
  % Loslager
  % -----------------------------------------------------------------------
  pics/mech/roller-support/.style args={#1}{
    code={
      \pgfkeys{/mech/support-common/.cd,#1}
      \edef\mechSupportWidth{\pgfkeysvalueof{/mech/support-common/width}}
      \edef\mechSupportHeight{\pgfkeysvalueof{/mech/support-common/height}}
      \edef\mechSupportHatch{\pgfkeysvalueof{/mech/support-common/hatch angle}}
      \pgfmathsetlengthmacro{\mechSupportHalfWidth}{0.5*\mechSupportWidth}
      % The triangle outline extends by half its line width at the base corners.
      % Add 0.3pt so the roller ground has the same visible total width.
      \pgfmathsetlengthmacro{\mechSupportGroundHalfWidth}{\mechSupportHalfWidth+0.3pt}
      \pgfmathsetlengthmacro{\mechSupportApexY}{0mm}
      \pgfmathsetlengthmacro{\mechSupportRollerGap}{0.272*\mechSupportHeight}
      \pgfmathsetlengthmacro{\mechSupportHatchDepth}{0.256*\mechSupportHeight}
      \pgfmathsetlengthmacro{\mechSupportGroundY}{\mechSupportHeight+\mechSupportRollerGap}

      \coordinate (-connection) at (0,0);
      \coordinate (-center) at (0,0);

      \draw[mech/support]
        (-\mechSupportHalfWidth,-\mechSupportHeight)
        -- (0,\mechSupportApexY)
        -- (\mechSupportHalfWidth,-\mechSupportHeight)
        -- cycle;
      \draw[mech/body] (0,0) circle[radius=.6mm];
      \draw[black,line width=.6pt,line cap=butt]
        (-\mechSupportGroundHalfWidth,-\mechSupportGroundY)
        -- (\mechSupportGroundHalfWidth,-\mechSupportGroundY);
      \fill[
        pattern={Lines[angle=\mechSupportHatch,distance=1pt]},
        pattern color=black
      ]
        (-\mechSupportGroundHalfWidth,-\mechSupportGroundY)
        rectangle
        (\mechSupportGroundHalfWidth,{-(\mechSupportGroundY+\mechSupportHatchDepth)});

      \node[\pgfkeysvalueof{/mech/support-common/label position}]
        at (0,0) {\pgfkeysvalueof{/mech/support-common/label}};
    }
  },
  pics/mech/roller-support/.default={},
}
\pgfkeys{
  /mech/fixed-support/.is family,
  /mech/fixed-support,
  width/.initial=1cm,
  depth/.initial=1.5mm,
  hatch angle/.initial=45,
  angle/.initial=0,
  label/.initial=,
}
\tikzset{
  pics/mech/fixed-support/.style args={#1}{
    code={
      \pgfkeys{/mech/fixed-support/.cd,#1}
      \edef\mechFixedWidth{\pgfkeysvalueof{/mech/fixed-support/width}}
      \edef\mechFixedDepth{\pgfkeysvalueof{/mech/fixed-support/depth}}
      \edef\mechFixedHatch{\pgfkeysvalueof{/mech/fixed-support/hatch angle}}
      \edef\mechFixedAngle{\pgfkeysvalueof{/mech/fixed-support/angle}}

      \coordinate (-center) at (0,0);
      \coordinate (-connection-upper)  at (0,\mechConnectionOffset);
      \coordinate (-connection-middle) at (0,0);
      \coordinate (-connection-lower)  at (0,-\mechConnectionOffset);

      \begin{scope}[rotate=\mechFixedAngle]
        \draw[mech/body,fill=none] (-.5\mechFixedWidth,0)--(.5\mechFixedWidth,0);
        \fill[
          pattern={Lines[angle=\mechFixedHatch,distance={3pt/sqrt(2)}]},
          pattern color=black
        ]
          (-.5\mechFixedWidth,0)
          rectangle
          (.5\mechFixedWidth,\mechFixedDepth);
      \end{scope}

      \node[above] at (0,0) {\pgfkeysvalueof{/mech/fixed-support/label}};
    }
  },
  pics/mech/fixed-support/.default={},
}

% =========================================================
% mechlib/elements/connectors.tex
% ATOMARE VERBINDUNGSELEMENTE: FEDER UND DAEMPFER
% =========================================================
% Beide Elemente folgen derselben grundlegenden API: "from", "to" und "label".
% Die Endpunkte koennen beliebige TikZ-Koordinaten oder exportierte Anker aus
% anderen mechlib-Pics sein.
%
% Oeffentliche Pics:
%   mech/spring  Zickzack-Feder entlang der Verbindungslinie
%   mech/damper  Daempfer mit rechteckigem Daempferkoerper in Linienmitte
%
% Die Beschriftung wird "sloped" gezeichnet und folgt daher der Orientierung
% des Verbindungselements. Ueber "label side" kann sie ober- oder unterhalb
% der lokalen Verbindungslinie platziert werden.
% =========================================================
\pgfkeys{
  /mech/spring/.is family,
  /mech/spring,
  from/.initial={(0,0)},
  to/.initial={(1,0)},
  label/.initial=,
  label side/.initial=above,
  amplitude/.initial=2mm,
  segment length/.initial=4mm,
  style/.initial=mech/connector,
}
\tikzset{
  pics/mech/spring/.style args={#1}{
    code={
      \pgfkeys{/mech/spring/.cd,#1}
      \edef\mechSpringFrom{\pgfkeysvalueof{/mech/spring/from}}
      \edef\mechSpringTo{\pgfkeysvalueof{/mech/spring/to}}
      \draw[\pgfkeysvalueof{/mech/spring/style},
        decorate,
        decoration={zigzag,
          segment length=\pgfkeysvalueof{/mech/spring/segment length},
          amplitude=\pgfkeysvalueof{/mech/spring/amplitude}}]
        \mechSpringFrom -- \mechSpringTo
        node[midway,sloped,mech/label,outer sep=2mm,\pgfkeysvalueof{/mech/spring/label side}]
        {\pgfkeysvalueof{/mech/spring/label}};
    }
  },
  pics/mech/spring/.default={},
}
\pgfkeys{
  /mech/damper/.is family,
  /mech/damper,
  from/.initial={(0,0)},
  to/.initial={(1,0)},
  label/.initial=,
  label side/.initial=below,
  body width/.initial=8mm,
  body height/.initial=3mm,
  style/.initial=mech/connector,
}
\tikzset{
  pics/mech/damper/.style args={#1}{
    code={
      \pgfkeys{/mech/damper/.cd,#1}
      \edef\mechDamperFrom{\pgfkeysvalueof{/mech/damper/from}}
      \edef\mechDamperTo{\pgfkeysvalueof{/mech/damper/to}}
      \draw[\pgfkeysvalueof{/mech/damper/style}]
        \mechDamperFrom -- \mechDamperTo
        node[midway,sloped,draw,fill=white,
          minimum width=\pgfkeysvalueof{/mech/damper/body width},
          minimum height=\pgfkeysvalueof{/mech/damper/body height},
          inner sep=0pt] {};
      \path \mechDamperFrom -- \mechDamperTo
        node[midway,sloped,mech/label,outer sep=2mm,\pgfkeysvalueof{/mech/damper/label side}]
        {\pgfkeysvalueof{/mech/damper/label}};
    }
  },
  pics/mech/damper/.default={},
}

% =========================================================
% mechlib/elements/trusses.tex
% FACHWERKE: STAB UND GEMEINSAME GELENKDARSTELLUNG
% =========================================================
% Dieses Modul enthaelt die atomaren Fachwerkelemente. Die Topologie eines
% Fachwerks bleibt bewusst im Anwendungsdokument: Knoten werden als TikZ-
% Koordinaten definiert und durch mech/rod bzw. \mechJoints verbunden.
%
% Oeffentliche API:
%   mech/rod     Stab zwischen "from" und "to" mit optionalem Label
%   \mechJoints{A,B,C,...} zeichnet alle angegebenen Gelenke
%
% Darstellungsregeln:
%   - Staebe verwenden standardmaessig den Stil mech/rod: schwarz und duenn.
%   - Staebe werden auf dem Layer "mech background" gezeichnet.
%   - Stabnummern besitzen keine Fuellung (fill=none) und bleiben transparent.
%   - Gelenke liegen optisch vor den Staeben und sind mit mechBodyColor gefuellt.
%
% Beispiel:
%   \pic {mech/rod={from=(A),to=(B),label=1,label side=above}};
%   \mechJoints{A,B}
% =========================================================
% -------------------------------------------------------------------------
% Einzelner Fachwerkstab
% -------------------------------------------------------------------------
% "style" erlaubt lokale Abweichungen, ohne den globalen mech/rod-Stil zu
% veraendern. Die Standardbeschriftung liegt an der Stabmitte.
\pgfkeys{
  /mech/rod/.is family,
  /mech/rod,
  from/.initial={(0,0)},
  to/.initial={(1,0)},
  label/.initial=,
  label side/.initial=above,
  label style/.initial={mech/label,fill=none,text=black},
  style/.initial=mech/rod,
}
\tikzset{
  pics/mech/rod/.style args={#1}{
    code={
      \pgfkeys{/mech/rod/.cd,#1}
      \edef\mechRodFrom{\pgfkeysvalueof{/mech/rod/from}}
      \edef\mechRodTo{\pgfkeysvalueof{/mech/rod/to}}
      \coordinate (-start) at \mechRodFrom;
      \coordinate (-end) at \mechRodTo;
      \coordinate (-center) at ($(-start)!.5!(-end)$);
      \begin{pgfonlayer}{mech background}
        \draw[\pgfkeysvalueof{/mech/rod/style}] (-start)--(-end);
      \end{pgfonlayer}
      \node[\pgfkeysvalueof{/mech/rod/label style},\pgfkeysvalueof{/mech/rod/label side}]
        at (-center) {\pgfkeysvalueof{/mech/rod/label}};
    }
  },
  pics/mech/rod/.default={},
}
% -------------------------------------------------------------------------
% Gelenkliste
% -------------------------------------------------------------------------
% Das Makro erwartet eine kommaseparierte Liste bereits definierter TikZ-
% Koordinaten. Jedes Gelenk wird identisch gezeichnet. Die Gelenke sollten nach
% den Staeben aufgerufen werden, damit sie die Stabenden optisch ueberdecken.
\newcommand{\mechJoints}[1]{%
  \foreach \mechJoint in {#1}{%
    \draw[fill=mechBodyColor,draw=black,line width=.2mm]
      (\mechJoint) circle[radius=.6mm];%
  }%
}

% -------------------------------------------------------------------------
% Annotationen und mechanische Symbolgruppen.
% -------------------------------------------------------------------------
% =========================================================
% mechlib/annotations/loads.tex
% LASTANNOTATIONEN: KRAFT, MOMENT UND STRECKENLAST
% =========================================================
% Diese Datei stellt die Lastsymbole der mechlib bereit. Alle Lasten verwenden
% standardmaessig den zentralen Stil mech/load und damit mechLoadColor.
%
% Oeffentliche Pics:
%   mech/force            Einzelkraft als Winkel/Laenge oder from/to
%   mech/moment           Kreisbogensymbol im oder gegen den Uhrzeigersinn
%   mech/distributed-load konstante, dreieckige oder linear veraenderliche Last
%
% Besonderheit der Streckenlast:
% Die Grundlinie wird lokal auf die x-Achse gedreht. Dadurch kann die gesamte
% Geometrie in einem einfachen lokalen Koordinatensystem berechnet werden. Sehr
% kurze Pfeile, deren Laenge kleiner als ihre Pfeilspitze ist, werden bewusst
% unterdrueckt; so entstehen am Nullende einer Dreieckslast keine Artefakte.
% =========================================================
\pgfkeys{
  /mech/force/.is family,
  /mech/force,
  from/.initial={(0,0)},
  to/.initial=,
  angle/.initial=0,
  length/.initial=1.5cm,
  label/.initial=,
  label position/.initial=above,
  direction/.initial=away,
  style/.initial=mech/load,
}
\tikzset{
  pics/mech/force/.style args={#1}{
    code={
      \pgfkeys{/mech/force/.cd,#1}
      \edef\mechForceFrom{\pgfkeysvalueof{/mech/force/from}}
      \edef\mechForceTo{\pgfkeysvalueof{/mech/force/to}}
      \edef\mechForceDirection{\pgfkeysvalueof{/mech/force/direction}}
      \def\mechForceToward{toward}
      \ifx\mechForceTo\empty
        \ifx\mechForceDirection\mechForceToward
          \draw[\pgfkeysvalueof{/mech/force/style},-{Latex[length=2.5mm]}]
            (\pgfkeysvalueof{/mech/force/angle}:\pgfkeysvalueof{/mech/force/length}) -- (0,0)
            node[midway,\pgfkeysvalueof{/mech/force/label position}]
            {\pgfkeysvalueof{/mech/force/label}};
        \else
          \draw[\pgfkeysvalueof{/mech/force/style},-{Latex[length=2.5mm]}]
            (0,0) -- (\pgfkeysvalueof{/mech/force/angle}:\pgfkeysvalueof{/mech/force/length})
            node[midway,\pgfkeysvalueof{/mech/force/label position}]
            {\pgfkeysvalueof{/mech/force/label}};
        \fi
      \else
        \ifx\mechForceDirection\mechForceToward
          \draw[\pgfkeysvalueof{/mech/force/style},-{Latex[length=2.5mm]}]
            \mechForceTo -- \mechForceFrom
            node[midway,\pgfkeysvalueof{/mech/force/label position}]
            {\pgfkeysvalueof{/mech/force/label}};
        \else
          \draw[\pgfkeysvalueof{/mech/force/style},-{Latex[length=2.5mm]}]
            \mechForceFrom -- \mechForceTo
            node[midway,\pgfkeysvalueof{/mech/force/label position}]
            {\pgfkeysvalueof{/mech/force/label}};
        \fi
      \fi
    }
  },
  pics/mech/force/.default={},
}
\pgfkeys{
  /mech/moment/.is family,
  /mech/moment,
  radius/.initial=6mm,
  label/.initial=,
  direction/.is choice,
  direction/ccw/.code={\def\mechMomentStart{-60}\def\mechMomentEnd{220}},
  direction/cw/.code={\def\mechMomentStart{240}\def\mechMomentEnd{-40}},
  direction=ccw,
  style/.initial=mech/load,
}
\tikzset{
  pics/mech/moment/.style args={#1}{
    code={
      \pgfkeys{/mech/moment/.cd,direction=ccw,#1}
      \draw[\pgfkeysvalueof{/mech/moment/style},-{Latex[length=2.5mm]}]
        (\mechMomentStart:\pgfkeysvalueof{/mech/moment/radius})
        arc[start angle=\mechMomentStart,end angle=\mechMomentEnd,
            radius=\pgfkeysvalueof{/mech/moment/radius}];
      \node[mech/label,above] at (0,\pgfkeysvalueof{/mech/moment/radius})
        {\pgfkeysvalueof{/mech/moment/label}};
    }
  },
  pics/mech/moment/.default={},
}
\pgfkeys{
  /mech/distributed-load/.is family,
  /mech/distributed-load,
  from/.initial={(0,0)},
  to/.initial={(1,0)},
  start value/.initial=1,
  end value/.initial=1,
  height/.initial=8mm,
  offset/.initial=0mm,
  arrows/.initial=6,
  arrowhead length/.initial=2mm,
  label/.initial=,
  label position/.initial=above,
  style/.initial=mech/load,
}
\tikzset{
  pics/mech/distributed-load/.style args={#1}{
    code={
      \pgfkeys{/mech/distributed-load/.cd,#1}
      \edef\mechDLFrom{\pgfkeysvalueof{/mech/distributed-load/from}}
      \edef\mechDLTo{\pgfkeysvalueof{/mech/distributed-load/to}}
      \path let
        \p1=\mechDLFrom,\p2=\mechDLTo,
        \n1={veclen(\x2-\x1,\y2-\y1)},
        \n2={atan2(\y2-\y1,\x2-\x1)}
      in
        coordinate (-start) at (\x1,\y1)
        coordinate (-end) at (\x2,\y2)
        coordinate (-helper) at (\n2,\n1);
      \path (-helper); \pgfgetlastxy{\mechDLAngle}{\mechDLLength};
      \pgfmathsetmacro{\mechDLStartValue}{\pgfkeysvalueof{/mech/distributed-load/start value}}
      \pgfmathsetmacro{\mechDLEndValue}{\pgfkeysvalueof{/mech/distributed-load/end value}}
      \pgfmathsetmacro{\mechDLArrowCount}{\pgfkeysvalueof{/mech/distributed-load/arrows}}
      \begin{scope}[shift={(-start)},rotate=\mechDLAngle]
        \pgfmathsetlengthmacro{\mechDLOffset}{\pgfkeysvalueof{/mech/distributed-load/offset}}
        \pgfmathsetlengthmacro{\mechDLArrowheadLength}{\pgfkeysvalueof{/mech/distributed-load/arrowhead length}}
        \pgfmathsetlengthmacro{\mechDLStartArrowLength}{abs(\mechDLStartValue)*\pgfkeysvalueof{/mech/distributed-load/height}}

        \fill[mech/load fill]
          (0,\mechDLOffset)
          --(0,{\mechDLOffset+\mechDLStartValue*\pgfkeysvalueof{/mech/distributed-load/height}})
          --(\mechDLLength,{\mechDLOffset+\mechDLEndValue*\pgfkeysvalueof{/mech/distributed-load/height}})
          --(\mechDLLength,\mechDLOffset)
          --cycle;

        % Keep contour and first arrow continuous whenever the first arrow
        % is long enough to accommodate its arrowhead. Otherwise draw only
        % the contour and suppress the too-short arrow.
        \ifdim\mechDLStartArrowLength<\mechDLArrowheadLength
          \draw[\pgfkeysvalueof{/mech/distributed-load/style}]
            (\mechDLLength,{\mechDLOffset+\mechDLEndValue*\pgfkeysvalueof{/mech/distributed-load/height}})
            --(0,{\mechDLOffset+\mechDLStartValue*\pgfkeysvalueof{/mech/distributed-load/height}});
        \else
          \draw[
            \pgfkeysvalueof{/mech/distributed-load/style},
            -{Latex[length=\pgfkeysvalueof{/mech/distributed-load/arrowhead length}]}
          ]
            (\mechDLLength,{\mechDLOffset+\mechDLEndValue*\pgfkeysvalueof{/mech/distributed-load/height}})
            --(0,{\mechDLOffset+\mechDLStartValue*\pgfkeysvalueof{/mech/distributed-load/height}})
            --(0,\mechDLOffset);
        \fi

        \foreach \mechDLIndex in {1,...,\mechDLArrowCount}{
          \pgfmathsetmacro{\mechDLT}{\mechDLIndex/\mechDLArrowCount}
          \pgfmathsetmacro{\mechDLValue}{\mechDLStartValue+(\mechDLEndValue-\mechDLStartValue)*\mechDLT}
          \pgfmathsetlengthmacro{\mechDLArrowLength}{abs(\mechDLValue)*\pgfkeysvalueof{/mech/distributed-load/height}}

          \ifdim\mechDLArrowLength<\mechDLArrowheadLength
            % Suppress arrows that are shorter than their arrowhead.
          \else
            \draw[
              \pgfkeysvalueof{/mech/distributed-load/style},
              -{Latex[length=\pgfkeysvalueof{/mech/distributed-load/arrowhead length}]}
            ]
              ({\mechDLT*\mechDLLength},{\mechDLOffset+\mechDLValue*\pgfkeysvalueof{/mech/distributed-load/height}})
              --({\mechDLT*\mechDLLength},\mechDLOffset);
          \fi
        }

        \node[mech/label,\pgfkeysvalueof{/mech/distributed-load/label position}]
          at (\mechDLLength,{\mechDLOffset+\mechDLEndValue*\pgfkeysvalueof{/mech/distributed-load/height}})
          {\pgfkeysvalueof{/mech/distributed-load/label}};
      \end{scope}
    }
  },
  pics/mech/distributed-load/.default={},
}

% =========================================================
% mechlib/annotations/dimensions.tex
% BEMASSUNGEN: LINEARES MASS UND WINKELMASS
% =========================================================
% Lineare Bemassungen werden geometrisch aus zwei Endpunkten abgeleitet. Der
% Offset wirkt senkrecht auf die Verbindung von "from" nach "to"; damit
% funktioniert dasselbe Pic fuer horizontale, vertikale und schraege Masse.
%
% Oeffentliche Pics:
%   mech/dimension       lineare Bemassung mit Doppelpfeil
%   mech/angle-dimension Kreisbogen zwischen zwei Winkeln
%
% Interne Koordinaten -A und -B sind die senkrecht versetzten Endpunkte der
% Masslinie. Die Vektorrechnung erfolgt in TikZ "let"-Syntax.
% =========================================================
\pgfkeys{
  /mech/dimension/.is family,
  /mech/dimension,
  from/.initial={(0,0)},
  to/.initial={(1,0)},
  label/.initial={$l$},
  offset/.initial=3mm,
  style/.initial=mech/dimension,
}
\tikzset{
  pics/mech/dimension/.style args={#1}{
    code={
      \pgfkeys{/mech/dimension/.cd,#1}
      \edef\mechDimFrom{\pgfkeysvalueof{/mech/dimension/from}}
      \edef\mechDimTo{\pgfkeysvalueof{/mech/dimension/to}}
      \path let
        \p1=\mechDimFrom,\p2=\mechDimTo,
        \n1={veclen(\x2-\x1,\y2-\y1)},
        \n2={-(\y2-\y1)/\n1*\pgfkeysvalueof{/mech/dimension/offset}},
        \n3={(\x2-\x1)/\n1*\pgfkeysvalueof{/mech/dimension/offset}}
      in
        coordinate (-A) at (\x1+\n2,\y1+\n3)
        coordinate (-B) at (\x2+\n2,\y2+\n3);
      \draw[\pgfkeysvalueof{/mech/dimension/style}]
        (-A)--(-B)
        node[midway,fill=white,inner sep=1pt]
        {\pgfkeysvalueof{/mech/dimension/label}};
    }
  },
  pics/mech/dimension/.default={},
}
\pgfkeys{
  /mech/angle-dimension/.is family,
  /mech/angle-dimension,
  start angle/.initial=0,
  end angle/.initial=45,
  radius/.initial=7mm,
  label/.initial={$\alpha$},
}
\tikzset{
  pics/mech/angle-dimension/.style args={#1}{
    code={
      \pgfkeys{/mech/angle-dimension/.cd,#1}
      \draw
        (\pgfkeysvalueof{/mech/angle-dimension/start angle}:\pgfkeysvalueof{/mech/angle-dimension/radius})
        arc[start angle=\pgfkeysvalueof{/mech/angle-dimension/start angle},
            end angle=\pgfkeysvalueof{/mech/angle-dimension/end angle},
            radius=\pgfkeysvalueof{/mech/angle-dimension/radius}];
      \pgfmathsetmacro{\mechAngleMid}{.5*(\pgfkeysvalueof{/mech/angle-dimension/start angle}+\pgfkeysvalueof{/mech/angle-dimension/end angle})}
      \node at (\mechAngleMid:{.7*\pgfkeysvalueof{/mech/angle-dimension/radius}})
        {\pgfkeysvalueof{/mech/angle-dimension/label}};
    }
  },
  pics/mech/angle-dimension/.default={},
}

% =========================================================
% mechlib/annotations/coordinates.tex
% KOORDINATENANNOTATIONEN: EINZELACHSE UND 2D-KOORDINATENSYSTEM
% =========================================================
% Die Pics dieser Datei zeichnen mathematische Bezugsachsen im Stil mech/axis.
% Winkel werden in TikZ-Grad angegeben; Laengen sind regulaere TeX-Dimensionen.
%
% Oeffentliche Pics:
%   mech/axis              einzelne gerichtete Achse
%   mech/coordinate-system zwei unabhaengig orientierbare Achsen
%
% Beide Pics exportieren den lokalen Ursprung als "-origin". Bei einem benannten
% Pic (cs) kann dieser z. B. als (cs-origin) referenziert werden.
% =========================================================
\pgfkeys{
  /mech/axis/.is family,
  /mech/axis,
  label/.initial={$x$},
  angle/.initial=0,
  length/.initial=1cm,
  direction/.initial=positive,
}
\tikzset{
  pics/mech/axis/.style args={#1}{
    code={
      \pgfkeys{/mech/axis/.cd,#1}
      \edef\mechAxisDirection{\pgfkeysvalueof{/mech/axis/direction}}
      \def\mechAxisNegative{negative}
      \ifx\mechAxisDirection\mechAxisNegative
        \draw[mech/axis,{Latex[length=2mm]}-]
          (0,0)--(\pgfkeysvalueof{/mech/axis/angle}:\pgfkeysvalueof{/mech/axis/length})
          node[pos=1,anchor=west]{\pgfkeysvalueof{/mech/axis/label}};
      \else
        \draw[mech/axis,-{Latex[length=2mm]}]
          (0,0)--(\pgfkeysvalueof{/mech/axis/angle}:\pgfkeysvalueof{/mech/axis/length})
          node[pos=1,anchor=west]{\pgfkeysvalueof{/mech/axis/label}};
      \fi
      \coordinate (-origin) at (0,0);
    }
  },
  pics/mech/axis/.default={},
}
\pgfkeys{
  /mech/coordinate-system/.is family,
  /mech/coordinate-system,
  x label/.initial={$x$},
  y label/.initial={$y$},
  x angle/.initial=0,
  y angle/.initial=90,
  x length/.initial=1cm,
  y length/.initial=1cm,
}
\tikzset{
  pics/mech/coordinate-system/.style args={#1}{
    code={
      \pgfkeys{/mech/coordinate-system/.cd,#1}
      \draw[mech/axis,-{Latex[length=2mm]}]
        (0,0)--(\pgfkeysvalueof{/mech/coordinate-system/x angle}:\pgfkeysvalueof{/mech/coordinate-system/x length})
        node[pos=1,anchor=west]{\pgfkeysvalueof{/mech/coordinate-system/x label}};
      \draw[mech/axis,-{Latex[length=2mm]}]
        (0,0)--(\pgfkeysvalueof{/mech/coordinate-system/y angle}:\pgfkeysvalueof{/mech/coordinate-system/y length})
        node[pos=1,anchor=east]{\pgfkeysvalueof{/mech/coordinate-system/y label}};
      \coordinate (-origin) at (0,0);
    }
  },
  pics/mech/coordinate-system/.default={},
}

% =========================================================
% mechlib/annotations/mechanics.tex
% MECHANISCHE ANNOTATIONEN: SCHNITTGROESSEN UND LAGERREAKTIONEN
% =========================================================
% Die Pics kombinieren mehrere elementare Lastsymbole zu haeufig benoetigten
% mechanischen Symbolgruppen. Farben und Pfeilgestaltung folgen mech/load.
%
% Oeffentliche Pics:
%   mech/internal-forces  Normalkraft N, Querkraft Q und Moment M am Schnitt
%   mech/fixed-reactions  horizontale/vertikale Reaktion und Einspannmoment
%
% Bei mech/internal-forces steuert "side=right|left" die Vorzeichenrichtung
% aller drei Schnittgroessen konsistent. Das Moment wird dazu intern als
% mech/moment mit passender Drehrichtung wiederverwendet.
% =========================================================
\pgfkeys{
  /mech/internal-forces/.is family,
  /mech/internal-forces,
  normal/.initial={$N$},
  shear/.initial={$Q$},
  moment/.initial={$M$},
  length/.initial=7mm,
  side/.is choice,
  side/right/.code={\def\mechIFSign{1}},
  side/left/.code={\def\mechIFSign{-1}},
  side=right,
}
\tikzset{
  pics/mech/internal-forces/.style args={#1}{
    code={
      \pgfkeys{/mech/internal-forces/.cd,side=right,#1}
      \edef\mechIFLength{\pgfkeysvalueof{/mech/internal-forces/length}}
      \draw[mech/load,-{Latex[length=2mm]}]
        (0,0)--({\mechIFSign*\mechIFLength},0)
        node[pos=1,anchor=west]{\pgfkeysvalueof{/mech/internal-forces/normal}};
      \draw[mech/load,-{Latex[length=2mm]}]
        (0,0)--(0,{-.8*\mechIFSign*\mechIFLength})
        node[pos=1,anchor=west]{\pgfkeysvalueof{/mech/internal-forces/shear}};
      \ifnum\mechIFSign=1
        \pic {mech/moment={direction=ccw,radius=.55\mechIFLength,label=\pgfkeysvalueof{/mech/internal-forces/moment}}};
      \else
        \pic {mech/moment={direction=cw,radius=.55\mechIFLength,label=\pgfkeysvalueof{/mech/internal-forces/moment}}};
      \fi
    }
  },
  pics/mech/internal-forces/.default={},
}
\pgfkeys{
  /mech/fixed-reactions/.is family,
  /mech/fixed-reactions,
  x/.initial={$A_x$},
  y/.initial={$A_y$},
  moment/.initial={$M_A$},
  length/.initial=7mm,
}
\tikzset{
  pics/mech/fixed-reactions/.style args={#1}{
    code={
      \pgfkeys{/mech/fixed-reactions/.cd,#1}
      \edef\mechReactionLength{\pgfkeysvalueof{/mech/fixed-reactions/length}}
      \draw[mech/load,-{Latex[length=2mm]}] (-\mechReactionLength,0)--(0,0)
        node[pos=0,anchor=east]{\pgfkeysvalueof{/mech/fixed-reactions/x}};
      \draw[mech/load,-{Latex[length=2mm]}] (0,.7\mechReactionLength)--(0,-.3\mechReactionLength)
        node[pos=0,anchor=south]{\pgfkeysvalueof{/mech/fixed-reactions/y}};
      \pic {mech/moment={direction=ccw,radius=.45\mechReactionLength,label=\pgfkeysvalueof{/mech/fixed-reactions/moment}}};
    }
  },
  pics/mech/fixed-reactions/.default={},
}


%
% Liegt die Bibliothek an einem anderen Ort, kann vor dem Laden die Variable
% \mechlibpath gesetzt werden:
%   \newcommand{\mechlibpath}{../../../support/mechlib}
%   % =========================================================
% mechlib/mechlib.tex
% ZENTRALER LOADER / OEFFENTLICHER EINSTIEGSPUNKT
% =========================================================
% Diese Datei ist die einzige Datei, die ein Anwender normalerweise direkt
% einbindet. Sie laedt die benoetigten LaTeX-/TikZ-Abhaengigkeiten und danach
% alle Teilmodule der Bibliothek in einer definierten Reihenfolge.
%
% Typische Verwendung bei einem Ordner "mechlib" neben dem Hauptdokument:
%   % =========================================================
% mechlib/mechlib.tex
% ZENTRALER LOADER / OEFFENTLICHER EINSTIEGSPUNKT
% =========================================================
% Diese Datei ist die einzige Datei, die ein Anwender normalerweise direkt
% einbindet. Sie laedt die benoetigten LaTeX-/TikZ-Abhaengigkeiten und danach
% alle Teilmodule der Bibliothek in einer definierten Reihenfolge.
%
% Typische Verwendung bei einem Ordner "mechlib" neben dem Hauptdokument:
%   \input{mechlib/mechlib}
%
% Liegt die Bibliothek an einem anderen Ort, kann vor dem Laden die Variable
% \mechlibpath gesetzt werden:
%   \newcommand{\mechlibpath}{../../../support/mechlib}
%   \input{\mechlibpath/mechlib}
%
% Lade-Reihenfolge:
%   1. core/        Farben, Layer, Stile und globale Geometriekonventionen
%   2. elements/    mechanische Grundelemente (Koerper, Lager, Verbinder, Staebe)
%   3. annotations/ Lasten, Masse, Koordinaten und mechanische Annotationen
%
% WICHTIG:
% Die Modulreihenfolge ist absichtlich festgelegt. Elemente greifen auf Farben,
% Layer und Stile aus core/ zurueck; Annotationen koennen wiederum die bereits
% definierten Grundstile und Pics verwenden.
% =========================================================
% -------------------------------------------------------------------------
% Externe Pakete und TikZ-Bibliotheken
% -------------------------------------------------------------------------
% \RequirePackage ist bewusst statt \usepackage gewaehlt, da mechlib.tex selbst
% wie eine Bibliotheks-/Paketdatei verwendet wird.
\RequirePackage{tikz}
\RequirePackage{amsmath}
\RequirePackage{siunitx}
\RequirePackage{xcolor}
\usetikzlibrary{
  arrows.meta,
  calc,
  patterns,
  patterns.meta,
  intersections,
  angles,
  quotes,
  decorations.pathmorphing,
  backgrounds
}
% Base path: override before loading when mechlib is stored elsewhere.
% -------------------------------------------------------------------------
% Aufloesung des Bibliothekspfades
% -------------------------------------------------------------------------
% \providecommand respektiert eine bereits im Hauptdokument gesetzte Definition.
\providecommand{\mechlibpath}{mechlib}
% -------------------------------------------------------------------------
% Kernmodule: muessen vor allen Elementen geladen werden.
% -------------------------------------------------------------------------
\input{\mechlibpath/core/colors}
\input{\mechlibpath/core/layers}
\input{\mechlibpath/core/styles}
\input{\mechlibpath/core/defaults}
% -------------------------------------------------------------------------
% Atomare mechanische Elemente.
% -------------------------------------------------------------------------
\input{\mechlibpath/elements/bodies}
\input{\mechlibpath/elements/supports}
\input{\mechlibpath/elements/connectors}
\input{\mechlibpath/elements/trusses}
% -------------------------------------------------------------------------
% Annotationen und mechanische Symbolgruppen.
% -------------------------------------------------------------------------
\input{\mechlibpath/annotations/loads}
\input{\mechlibpath/annotations/dimensions}
\input{\mechlibpath/annotations/coordinates}
\input{\mechlibpath/annotations/mechanics}

%
% Liegt die Bibliothek an einem anderen Ort, kann vor dem Laden die Variable
% \mechlibpath gesetzt werden:
%   \newcommand{\mechlibpath}{../../../support/mechlib}
%   % =========================================================
% mechlib/mechlib.tex
% ZENTRALER LOADER / OEFFENTLICHER EINSTIEGSPUNKT
% =========================================================
% Diese Datei ist die einzige Datei, die ein Anwender normalerweise direkt
% einbindet. Sie laedt die benoetigten LaTeX-/TikZ-Abhaengigkeiten und danach
% alle Teilmodule der Bibliothek in einer definierten Reihenfolge.
%
% Typische Verwendung bei einem Ordner "mechlib" neben dem Hauptdokument:
%   \input{mechlib/mechlib}
%
% Liegt die Bibliothek an einem anderen Ort, kann vor dem Laden die Variable
% \mechlibpath gesetzt werden:
%   \newcommand{\mechlibpath}{../../../support/mechlib}
%   \input{\mechlibpath/mechlib}
%
% Lade-Reihenfolge:
%   1. core/        Farben, Layer, Stile und globale Geometriekonventionen
%   2. elements/    mechanische Grundelemente (Koerper, Lager, Verbinder, Staebe)
%   3. annotations/ Lasten, Masse, Koordinaten und mechanische Annotationen
%
% WICHTIG:
% Die Modulreihenfolge ist absichtlich festgelegt. Elemente greifen auf Farben,
% Layer und Stile aus core/ zurueck; Annotationen koennen wiederum die bereits
% definierten Grundstile und Pics verwenden.
% =========================================================
% -------------------------------------------------------------------------
% Externe Pakete und TikZ-Bibliotheken
% -------------------------------------------------------------------------
% \RequirePackage ist bewusst statt \usepackage gewaehlt, da mechlib.tex selbst
% wie eine Bibliotheks-/Paketdatei verwendet wird.
\RequirePackage{tikz}
\RequirePackage{amsmath}
\RequirePackage{siunitx}
\RequirePackage{xcolor}
\usetikzlibrary{
  arrows.meta,
  calc,
  patterns,
  patterns.meta,
  intersections,
  angles,
  quotes,
  decorations.pathmorphing,
  backgrounds
}
% Base path: override before loading when mechlib is stored elsewhere.
% -------------------------------------------------------------------------
% Aufloesung des Bibliothekspfades
% -------------------------------------------------------------------------
% \providecommand respektiert eine bereits im Hauptdokument gesetzte Definition.
\providecommand{\mechlibpath}{mechlib}
% -------------------------------------------------------------------------
% Kernmodule: muessen vor allen Elementen geladen werden.
% -------------------------------------------------------------------------
\input{\mechlibpath/core/colors}
\input{\mechlibpath/core/layers}
\input{\mechlibpath/core/styles}
\input{\mechlibpath/core/defaults}
% -------------------------------------------------------------------------
% Atomare mechanische Elemente.
% -------------------------------------------------------------------------
\input{\mechlibpath/elements/bodies}
\input{\mechlibpath/elements/supports}
\input{\mechlibpath/elements/connectors}
\input{\mechlibpath/elements/trusses}
% -------------------------------------------------------------------------
% Annotationen und mechanische Symbolgruppen.
% -------------------------------------------------------------------------
\input{\mechlibpath/annotations/loads}
\input{\mechlibpath/annotations/dimensions}
\input{\mechlibpath/annotations/coordinates}
\input{\mechlibpath/annotations/mechanics}

%
% Lade-Reihenfolge:
%   1. core/        Farben, Layer, Stile und globale Geometriekonventionen
%   2. elements/    mechanische Grundelemente (Koerper, Lager, Verbinder, Staebe)
%   3. annotations/ Lasten, Masse, Koordinaten und mechanische Annotationen
%
% WICHTIG:
% Die Modulreihenfolge ist absichtlich festgelegt. Elemente greifen auf Farben,
% Layer und Stile aus core/ zurueck; Annotationen koennen wiederum die bereits
% definierten Grundstile und Pics verwenden.
% =========================================================
% -------------------------------------------------------------------------
% Externe Pakete und TikZ-Bibliotheken
% -------------------------------------------------------------------------
% \RequirePackage ist bewusst statt \usepackage gewaehlt, da mechlib.tex selbst
% wie eine Bibliotheks-/Paketdatei verwendet wird.
\RequirePackage{tikz}
\RequirePackage{amsmath}
\RequirePackage{siunitx}
\RequirePackage{xcolor}
\usetikzlibrary{
  arrows.meta,
  calc,
  patterns,
  patterns.meta,
  intersections,
  angles,
  quotes,
  decorations.pathmorphing,
  backgrounds
}
% Base path: override before loading when mechlib is stored elsewhere.
% -------------------------------------------------------------------------
% Aufloesung des Bibliothekspfades
% -------------------------------------------------------------------------
% \providecommand respektiert eine bereits im Hauptdokument gesetzte Definition.
\providecommand{\mechlibpath}{mechlib}
% -------------------------------------------------------------------------
% Kernmodule: muessen vor allen Elementen geladen werden.
% -------------------------------------------------------------------------
% =========================================================
% mechlib/core/colors.tex
% ZENTRALE FARBPALETTE
% =========================================================
% Alle bibliotheksweiten Farben werden ausschliesslich hier definiert. Dadurch
% kann der visuelle Charakter der gesamten mechlib an einer Stelle angepasst
% werden, ohne einzelne Elemente zu veraendern.
%
% Konvention:
%   mechBodyColor    Fuellung allgemeiner Koerper und Fachwerkgelenke
%   mechSupportColor Fuellung von Lagerkoerpern
%   mechLoadColor    Kraefte, Momente und Streckenlasten
%   mechAxisColor    Achsen und Koordinatensysteme
%   mechRodColor     historische Stabfarbe; aktuelle Fachwerkstaebe sind schwarz
%
% Die Namen sind Teil der internen visuellen API und sollten bevorzugt statt
% harter RGB-Werte in den Elementdateien verwendet werden.
% =========================================================
% -------------------------------------------------------------------------
% Farbdefinitionen
% -------------------------------------------------------------------------
% RGB wird bewusst explizit verwendet, damit die Darstellung unabhaengig von
% Dokumentklassen und externen Farbmodellen reproduzierbar bleibt.
\definecolor{mechBodyColor}{RGB}{180,180,180}
\definecolor{mechSupportColor}{RGB}{205,205,205}
\definecolor{mechLoadColor}{RGB}{178,25,25}
\definecolor{mechAxisColor}{RGB}{6,13,141}
\definecolor{mechRodColor}{RGB}{130,130,130}

% =========================================================
% mechlib/core/layers.tex
% ZEICHENEBENEN / PGF-LAYER
% =========================================================
% Fachwerkstaebe und andere hinterlegte Geometrien koennen auf dem Layer
% "mech background" gezeichnet werden. Der normale TikZ-Layer "main" liegt
% darueber. Dadurch koennen z. B. gefuellte Gelenke sauber vor den Staeben
% erscheinen, ohne dass die Zeichenreihenfolge im Anwendungsdokument aufwendig
% kontrolliert werden muss.
% =========================================================
% Zuerst wird der benannte Hintergrundlayer deklariert. Anschliessend wird die
% globale Layerreihenfolge explizit gesetzt: Hintergrund vor Hauptlayer.
\pgfdeclarelayer{mech background}
\pgfsetlayers{mech background,main}

% =========================================================
% mechlib/core/styles.tex
% ZENTRALE, NAMENSRAUMGESCHUETZTE TIKZ-STILE
% =========================================================
% Diese Datei definiert ausschliesslich visuelle Grundstile. Die eigentliche
% Geometrie eines Elements gehoert in elements/ bzw. annotations/.
%
% Alle oeffentlichen Bibliotheksstile beginnen mit "mech/". Das reduziert
% Namenskollisionen mit dem Hauptdokument und macht im TikZ-Code sichtbar,
% welche Formatierung aus der mechlib stammt.
%
% Designregel:
%   - Farben werden aus core/colors.tex bezogen.
%   - Geometrien werden NICHT hier definiert.
%   - Linienstarken werden hier zentral gehalten, soweit sie elementuebergreifend
%     konsistent sein sollen.
% =========================================================
% -------------------------------------------------------------------------
% Visuelle Basisstile
% -------------------------------------------------------------------------
% Einzelne Pics duerfen diese Stile ueber ihre jeweiligen "style"-Keys ersetzen.
\tikzset{
  mech/body/.style={
    draw=black,
    fill=mechBodyColor,
    line width=.6pt
  },
  mech/support/.style={
    draw=black,
    fill=mechSupportColor,
    line width=.6pt
  },
  mech/connector/.style={
    draw=black,
    line width=.6pt
  },
  mech/rod/.style={
    draw=black,
    line width=.85pt
  },
  mech/load/.style={
    draw=mechLoadColor,
    text=mechLoadColor,
    line width=1pt
  },
  mech/load fill/.style={
    fill=mechLoadColor,
    fill opacity=.16,
    draw=none
  },
  mech/dimension/.style={
    {Latex[length=2mm]}-{Latex[length=2mm]},
    line width=.5pt
  },
  mech/axis/.style={
    draw=mechAxisColor,
    text=mechAxisColor,
    line width=.6pt
  },
  mech/label/.style={
    inner sep=1pt
  }
}

% =========================================================
% mechlib/core/defaults.tex
% GLOBALE STANDARDWERTE UND GEMEINSAME GEOMETRIEKONVENTIONEN
% =========================================================
% Hier stehen Werte, die von mehreren Modulen gemeinsam benutzt werden. Solche
% Konstanten verhindern, dass geometrisch zusammengehoerige Elemente lokal mit
% leicht unterschiedlichen Zahlen definiert werden.
%
% Beispiel: \mechConnectionOffset legt die vertikale Lage paralleler
% Anschlussanker fest. Masse und Einspannung verwenden exakt denselben Wert;
% dadurch liegen obere/mittlere/untere Anschluesse bei gleicher Mittelpunktlage
% garantiert auf gleicher Hoehe.
% =========================================================
% Globale Schriftvorgabe fuer jedes TikZ-Bild. "append style" ergaenzt
% vorhandene Bildstile, statt sie zu ueberschreiben.
\tikzset{
  every picture/.append style={font=\fontsize{10}{10}\selectfont},
}

% Standardwinkel fuer technische Bodenschraffuren.
\def\mechHatchAngle{45}

% Shared vertical distance of parallel connector anchors from the middle anchor.
% Wall and mass use this exact value so upper/middle/lower are aligned whenever
% their center points lie on the same horizontal line.
% Gemeinsamer vertikaler Versatz fuer parallele Anschlussanker.
\def\mechConnectionOffset{2.5mm}

% -------------------------------------------------------------------------
% Atomare mechanische Elemente.
% -------------------------------------------------------------------------
% =========================================================
% mechlib/elements/bodies.tex
% ATOMARE KOERPER: MASSE, PUNKTMASSE, WAGEN UND BALKEN
% =========================================================
% Die hier definierten TikZ-Pics stellen mechanische Grundkoerper bereit. Jedes
% Pic besitzt eine eigene pgfkeys-Familie unter /mech/... und kann dadurch mit
% benannten Optionen konfiguriert werden.
%
% Grundprinzip der API:
%   \pic (name) at (<Ort>) {mech/<element>={<key>=<wert>,...}};
%
% Benannte Pics exportieren Koordinaten bzw. werden bei node-basierten Koerpern
% auf den internen Shape-Knoten aliasiert. So koennen regulaere TikZ-Anker wie
% (m1.east), (m1.north west) usw. direkt benutzt werden.
%
% Oeffentliche Pics in dieser Datei:
%   mech/mass       rechteckiger Starrkoerper bzw. konzentrierte Masse
%   mech/point-mass kreisfoermige Punktmasse
%   mech/cart       Wagen mit zwei Raedern und optionalem Boden
%   mech/beam       Balken zwischen zwei beliebigen Punkten
% =========================================================
\pgfkeys{
  /mech/mass/.is family,
  /mech/mass,
  width/.initial=1.5cm,
  height/.initial=1cm,
  label/.initial=,
  style/.initial=mech/body,
}
\tikzset{
  pics/mech/mass/.style args={#1}{
    code={
      \pgfkeys{/mech/mass/.cd,#1}
      \node[
        \pgfkeysvalueof{/mech/mass/style},
        minimum width=\pgfkeysvalueof{/mech/mass/width},
        minimum height=\pgfkeysvalueof{/mech/mass/height},
        inner sep=0pt
      ] (-shape) at (0,0) {\pgfkeysvalueof{/mech/mass/label}};

      % Named body pics behave like regular TikZ nodes.
      % Example: \pic (m1) ... then (m1.east), (m1.south east), ...
      \edef\mechPicName{\pgfkeysvalueof{/tikz/name prefix}}
      \ifx\mechPicName\pgfutil@empty\else
        \pgfnodealias{\mechPicName}{\mechPicName-shape}
      \fi

      \coordinate (-center) at (-shape.center);

      % Three standardized connection levels on each vertical side.
      % The shared \mechConnectionOffset is also used by fixed supports.
      \coordinate (-west-upper)  at ($(-shape.west)+(0,\mechConnectionOffset)$);
      \coordinate (-west-middle) at (-shape.west);
      \coordinate (-west-lower)  at ($(-shape.west)+(0,-\mechConnectionOffset)$);

      \coordinate (-east-upper)  at ($(-shape.east)+(0,\mechConnectionOffset)$);
      \coordinate (-east-middle) at (-shape.east);
      \coordinate (-east-lower)  at ($(-shape.east)+(0,-\mechConnectionOffset)$);

      % Generic connection aliases
      \coordinate (-connection-upper)  at (-east-upper);
      \coordinate (-connection-middle) at (-east-middle);
      \coordinate (-connection-lower)  at (-east-lower);

      \coordinate (-north) at (-shape.north);
      \coordinate (-south) at (-shape.south);
    }
  },
  pics/mech/mass/.default={},
}
\pgfkeys{
  /mech/point-mass/.is family,
  /mech/point-mass,
  radius/.initial=2.2mm,
  label/.initial=,
  style/.initial=mech/body,
}
\tikzset{
  pics/mech/point-mass/.style args={#1}{
    code={
      \pgfkeys{/mech/point-mass/.cd,#1}
      \edef\mechPointMassRadius{\pgfkeysvalueof{/mech/point-mass/radius}}
      \edef\mechPointMassLabel{\pgfkeysvalueof{/mech/point-mass/label}}
      \edef\mechPointMassStyle{\pgfkeysvalueof{/mech/point-mass/style}}
      \node[
        \mechPointMassStyle,
        circle,
        minimum size=2\mechPointMassRadius,
        inner sep=0pt
      ] (-shape) at (0,0) {\mechPointMassLabel};

      \edef\mechPicName{\pgfkeysvalueof{/tikz/name prefix}}
      \ifx\mechPicName\pgfutil@empty\else
        \pgfnodealias{\mechPicName}{\mechPicName-shape}
      \fi

      \coordinate (-center) at (-shape.center);
    }
  },
  pics/mech/point-mass/.default={},
}
% -------------------------------------------------------------------------
% Wagen
% -------------------------------------------------------------------------
% Der Boolean "ground" wird als TeX-if gespeichert. Dadurch kann die optionale
% Bodendarstellung ohne Stringvergleich ein- bzw. ausgeschaltet werden.
\newif\ifmechCartGround
\pgfkeys{
  /mech/cart/.is family,
  /mech/cart,
  width/.initial=2cm,
  height/.initial=1cm,
  wheel radius/.initial=2mm,
  label/.initial=,
  ground/.is if=mechCartGround,
  ground/.default=true,
  ground=true,
}
\tikzset{
  pics/mech/cart/.style args={#1}{
    code={
      \pgfkeys{/mech/cart/.cd,#1}
      \edef\mechCartWidth{\pgfkeysvalueof{/mech/cart/width}}
      \edef\mechCartHeight{\pgfkeysvalueof{/mech/cart/height}}
      \edef\mechCartWheelRadius{\pgfkeysvalueof{/mech/cart/wheel radius}}
      \node[
        mech/body,
        minimum width=\mechCartWidth,
        minimum height=\mechCartHeight,
        inner sep=0pt
      ] (-shape) at (0,0) {\pgfkeysvalueof{/mech/cart/label}};

      \edef\mechPicName{\pgfkeysvalueof{/tikz/name prefix}}
      \ifx\mechPicName\pgfutil@empty\else
        \pgfnodealias{\mechPicName}{\mechPicName-shape}
      \fi

      \coordinate (-center) at (-shape.center);
      \coordinate (-west) at (-shape.west);
      \coordinate (-east) at (-shape.east);
      \coordinate (-north) at (-shape.north);
      \coordinate (-south) at (-shape.south);
      \draw[mech/body]
        (-.3\mechCartWidth,-.5\mechCartHeight-\mechCartWheelRadius)
        circle[radius=\mechCartWheelRadius];
      \draw[mech/body]
        (.3\mechCartWidth,-.5\mechCartHeight-\mechCartWheelRadius)
        circle[radius=\mechCartWheelRadius];
      \ifmechCartGround
        \draw
          (-.5\mechCartWidth,-.5\mechCartHeight-2\mechCartWheelRadius)
          --
          (.5\mechCartWidth,-.5\mechCartHeight-2\mechCartWheelRadius);
        \fill[
          pattern={Lines[angle=\mechHatchAngle,distance=1pt]},
          pattern color=black
        ]
          (-.5\mechCartWidth,-.5\mechCartHeight-2\mechCartWheelRadius)
          rectangle
          (.5\mechCartWidth,-.5\mechCartHeight-2\mechCartWheelRadius-3pt);
      \fi
    }
  },
  pics/mech/cart/.default={},
}
\pgfkeys{
  /mech/beam/.is family,
  /mech/beam,
  from/.initial={(0,0)},
  to/.initial={(1,0)},
  height/.initial=3mm,
  style/.initial=mech/body,
}
\tikzset{
  pics/mech/beam/.style args={#1}{
    code={
      \pgfkeys{/mech/beam/.cd,#1}
      \edef\mechBeamFrom{\pgfkeysvalueof{/mech/beam/from}}
      \edef\mechBeamTo{\pgfkeysvalueof{/mech/beam/to}}
      \edef\mechBeamHeight{\pgfkeysvalueof{/mech/beam/height}}
      \edef\mechBeamStyle{\pgfkeysvalueof{/mech/beam/style}}
      \path let
        \p1=\mechBeamFrom,
        \p2=\mechBeamTo,
        \n1={veclen(\x2-\x1,\y2-\y1)},
        \n2={(\y2-\y1)/\n1*\mechBeamHeight/2},
        \n3={-(\x2-\x1)/\n1*\mechBeamHeight/2}
      in
        coordinate (-A) at (\x1+\n2,\y1+\n3)
        coordinate (-B) at (\x2+\n2,\y2+\n3)
        coordinate (-C) at (\x2-\n2,\y2-\n3)
        coordinate (-D) at (\x1-\n2,\y1-\n3)
        coordinate (-start) at (\x1,\y1)
        coordinate (-end) at (\x2,\y2)
        coordinate (-center) at ({(\x1+\x2)/2},{(\y1+\y2)/2});
      \draw[\mechBeamStyle] (-A)--(-B)--(-C)--(-D)--cycle;
    }
  },
  pics/mech/beam/.default={},
}

% =========================================================
% mechlib/elements/supports.tex
% ATOMARE LAGER: FESTLAGER, LOSLAGER UND EINSPANNUNG
% =========================================================
% Diese Datei kapselt die Lagergeometrie. Fest- und Loslager teilen sich die
% Key-Familie /mech/support-common, damit Breite, Hoehe, Schraffur und Label
% konsistent parametrisiert werden koennen.
%
% Oeffentliche Pics:
%   mech/pinned-support Festlager / gelenkiges Lager
%   mech/roller-support Loslager mit sichtbarem Abstand zum Boden
%   mech/fixed-support  Einspannung, optional gedreht
%
% Aktuelle Fachwerk-Konvention:
%   - Die Dreiecksspitze liegt exakt im Mittelpunkt des Gelenks (0,0).
%   - Das Dreieck ist gegenueber der frueheren Version kompakt skaliert.
%   - Bodenlinie und Schraffur verwenden dieselben x-Grenzen wie die
%     Dreiecksgrundseite; beim Loslager wird die sichtbare Linienbreite explizit
%     beruecksichtigt.
%   - Das Gelenk selbst ist mit mechBodyColor gefuellt.
%
% Die pgfmathsetlengthmacro-Aufrufe sind absichtlich verwendet: LaTeX-Laengen
% duerfen nicht durch Textersatz wie ".5\\macro" skaliert werden, da dies bei
% expandierenden Dimensionsmakros zu sichtbarem Resttext (z. B. ".8mm") fuehren
% kann. Alle abgeleiteten Abmessungen werden deshalb numerisch berechnet.
% =========================================================
\pgfkeys{
  /mech/support-common/.is family,
  /mech/support-common,
  width/.initial=2.72mm,
  height/.initial=2.24mm,
  hatch angle/.initial=45,
  label/.initial=,
  label position/.initial=above,
}
\tikzset{
  % -----------------------------------------------------------------------
  % Festlager
  % -----------------------------------------------------------------------
  pics/mech/pinned-support/.style args={#1}{
    code={
      \pgfkeys{/mech/support-common/.cd,#1}
      \edef\mechSupportWidth{\pgfkeysvalueof{/mech/support-common/width}}
      \edef\mechSupportHeight{\pgfkeysvalueof{/mech/support-common/height}}
      \edef\mechSupportHatch{\pgfkeysvalueof{/mech/support-common/hatch angle}}
      \pgfmathsetlengthmacro{\mechSupportHalfWidth}{0.5*\mechSupportWidth}
      % Die Dreieckskontur ragt an der Grundseite durch die halbe Linienbreite
      % minimal ueber die geometrische Breite hinaus. Fuer identische sichtbare
      % Breite wird die Bodenlinie daher um 0.3pt je Seite erweitert.
      \pgfmathsetlengthmacro{\mechSupportGroundHalfWidth}{\mechSupportHalfWidth+0.3pt}
      \pgfmathsetlengthmacro{\mechSupportApexY}{0mm}
      \pgfmathsetlengthmacro{\mechSupportHatchDepth}{0.256*\mechSupportHeight}

      \coordinate (-connection) at (0,0);
      \coordinate (-center) at (0,0);

      \draw[mech/support]
        (-\mechSupportHalfWidth,-\mechSupportHeight)
        -- (0,\mechSupportApexY)
        -- (\mechSupportHalfWidth,-\mechSupportHeight)
        -- cycle;
      \draw[mech/body] (0,0) circle[radius=.6mm];
      \draw[black,line width=.6pt,line cap=butt]
        (-\mechSupportGroundHalfWidth,-\mechSupportHeight)
        -- (\mechSupportGroundHalfWidth,-\mechSupportHeight);
      \fill[
        pattern={Lines[angle=\mechSupportHatch,distance=1pt]},
        pattern color=black
      ]
        (-\mechSupportGroundHalfWidth,-\mechSupportHeight)
        rectangle
        (\mechSupportGroundHalfWidth,{-(\mechSupportHeight+\mechSupportHatchDepth)});

      \node[\pgfkeysvalueof{/mech/support-common/label position}]
        at (0,0) {\pgfkeysvalueof{/mech/support-common/label}};
    }
  },
  pics/mech/pinned-support/.default={},
  % -----------------------------------------------------------------------
  % Loslager
  % -----------------------------------------------------------------------
  pics/mech/roller-support/.style args={#1}{
    code={
      \pgfkeys{/mech/support-common/.cd,#1}
      \edef\mechSupportWidth{\pgfkeysvalueof{/mech/support-common/width}}
      \edef\mechSupportHeight{\pgfkeysvalueof{/mech/support-common/height}}
      \edef\mechSupportHatch{\pgfkeysvalueof{/mech/support-common/hatch angle}}
      \pgfmathsetlengthmacro{\mechSupportHalfWidth}{0.5*\mechSupportWidth}
      % The triangle outline extends by half its line width at the base corners.
      % Add 0.3pt so the roller ground has the same visible total width.
      \pgfmathsetlengthmacro{\mechSupportGroundHalfWidth}{\mechSupportHalfWidth+0.3pt}
      \pgfmathsetlengthmacro{\mechSupportApexY}{0mm}
      \pgfmathsetlengthmacro{\mechSupportRollerGap}{0.272*\mechSupportHeight}
      \pgfmathsetlengthmacro{\mechSupportHatchDepth}{0.256*\mechSupportHeight}
      \pgfmathsetlengthmacro{\mechSupportGroundY}{\mechSupportHeight+\mechSupportRollerGap}

      \coordinate (-connection) at (0,0);
      \coordinate (-center) at (0,0);

      \draw[mech/support]
        (-\mechSupportHalfWidth,-\mechSupportHeight)
        -- (0,\mechSupportApexY)
        -- (\mechSupportHalfWidth,-\mechSupportHeight)
        -- cycle;
      \draw[mech/body] (0,0) circle[radius=.6mm];
      \draw[black,line width=.6pt,line cap=butt]
        (-\mechSupportGroundHalfWidth,-\mechSupportGroundY)
        -- (\mechSupportGroundHalfWidth,-\mechSupportGroundY);
      \fill[
        pattern={Lines[angle=\mechSupportHatch,distance=1pt]},
        pattern color=black
      ]
        (-\mechSupportGroundHalfWidth,-\mechSupportGroundY)
        rectangle
        (\mechSupportGroundHalfWidth,{-(\mechSupportGroundY+\mechSupportHatchDepth)});

      \node[\pgfkeysvalueof{/mech/support-common/label position}]
        at (0,0) {\pgfkeysvalueof{/mech/support-common/label}};
    }
  },
  pics/mech/roller-support/.default={},
}
\pgfkeys{
  /mech/fixed-support/.is family,
  /mech/fixed-support,
  width/.initial=1cm,
  depth/.initial=1.5mm,
  hatch angle/.initial=45,
  angle/.initial=0,
  label/.initial=,
}
\tikzset{
  pics/mech/fixed-support/.style args={#1}{
    code={
      \pgfkeys{/mech/fixed-support/.cd,#1}
      \edef\mechFixedWidth{\pgfkeysvalueof{/mech/fixed-support/width}}
      \edef\mechFixedDepth{\pgfkeysvalueof{/mech/fixed-support/depth}}
      \edef\mechFixedHatch{\pgfkeysvalueof{/mech/fixed-support/hatch angle}}
      \edef\mechFixedAngle{\pgfkeysvalueof{/mech/fixed-support/angle}}

      \coordinate (-center) at (0,0);
      \coordinate (-connection-upper)  at (0,\mechConnectionOffset);
      \coordinate (-connection-middle) at (0,0);
      \coordinate (-connection-lower)  at (0,-\mechConnectionOffset);

      \begin{scope}[rotate=\mechFixedAngle]
        \draw[mech/body,fill=none] (-.5\mechFixedWidth,0)--(.5\mechFixedWidth,0);
        \fill[
          pattern={Lines[angle=\mechFixedHatch,distance={3pt/sqrt(2)}]},
          pattern color=black
        ]
          (-.5\mechFixedWidth,0)
          rectangle
          (.5\mechFixedWidth,\mechFixedDepth);
      \end{scope}

      \node[above] at (0,0) {\pgfkeysvalueof{/mech/fixed-support/label}};
    }
  },
  pics/mech/fixed-support/.default={},
}

% =========================================================
% mechlib/elements/connectors.tex
% ATOMARE VERBINDUNGSELEMENTE: FEDER UND DAEMPFER
% =========================================================
% Beide Elemente folgen derselben grundlegenden API: "from", "to" und "label".
% Die Endpunkte koennen beliebige TikZ-Koordinaten oder exportierte Anker aus
% anderen mechlib-Pics sein.
%
% Oeffentliche Pics:
%   mech/spring  Zickzack-Feder entlang der Verbindungslinie
%   mech/damper  Daempfer mit rechteckigem Daempferkoerper in Linienmitte
%
% Die Beschriftung wird "sloped" gezeichnet und folgt daher der Orientierung
% des Verbindungselements. Ueber "label side" kann sie ober- oder unterhalb
% der lokalen Verbindungslinie platziert werden.
% =========================================================
\pgfkeys{
  /mech/spring/.is family,
  /mech/spring,
  from/.initial={(0,0)},
  to/.initial={(1,0)},
  label/.initial=,
  label side/.initial=above,
  amplitude/.initial=2mm,
  segment length/.initial=4mm,
  style/.initial=mech/connector,
}
\tikzset{
  pics/mech/spring/.style args={#1}{
    code={
      \pgfkeys{/mech/spring/.cd,#1}
      \edef\mechSpringFrom{\pgfkeysvalueof{/mech/spring/from}}
      \edef\mechSpringTo{\pgfkeysvalueof{/mech/spring/to}}
      \draw[\pgfkeysvalueof{/mech/spring/style},
        decorate,
        decoration={zigzag,
          segment length=\pgfkeysvalueof{/mech/spring/segment length},
          amplitude=\pgfkeysvalueof{/mech/spring/amplitude}}]
        \mechSpringFrom -- \mechSpringTo
        node[midway,sloped,mech/label,outer sep=2mm,\pgfkeysvalueof{/mech/spring/label side}]
        {\pgfkeysvalueof{/mech/spring/label}};
    }
  },
  pics/mech/spring/.default={},
}
\pgfkeys{
  /mech/damper/.is family,
  /mech/damper,
  from/.initial={(0,0)},
  to/.initial={(1,0)},
  label/.initial=,
  label side/.initial=below,
  body width/.initial=8mm,
  body height/.initial=3mm,
  style/.initial=mech/connector,
}
\tikzset{
  pics/mech/damper/.style args={#1}{
    code={
      \pgfkeys{/mech/damper/.cd,#1}
      \edef\mechDamperFrom{\pgfkeysvalueof{/mech/damper/from}}
      \edef\mechDamperTo{\pgfkeysvalueof{/mech/damper/to}}
      \draw[\pgfkeysvalueof{/mech/damper/style}]
        \mechDamperFrom -- \mechDamperTo
        node[midway,sloped,draw,fill=white,
          minimum width=\pgfkeysvalueof{/mech/damper/body width},
          minimum height=\pgfkeysvalueof{/mech/damper/body height},
          inner sep=0pt] {};
      \path \mechDamperFrom -- \mechDamperTo
        node[midway,sloped,mech/label,outer sep=2mm,\pgfkeysvalueof{/mech/damper/label side}]
        {\pgfkeysvalueof{/mech/damper/label}};
    }
  },
  pics/mech/damper/.default={},
}

% =========================================================
% mechlib/elements/trusses.tex
% FACHWERKE: STAB UND GEMEINSAME GELENKDARSTELLUNG
% =========================================================
% Dieses Modul enthaelt die atomaren Fachwerkelemente. Die Topologie eines
% Fachwerks bleibt bewusst im Anwendungsdokument: Knoten werden als TikZ-
% Koordinaten definiert und durch mech/rod bzw. \mechJoints verbunden.
%
% Oeffentliche API:
%   mech/rod     Stab zwischen "from" und "to" mit optionalem Label
%   \mechJoints{A,B,C,...} zeichnet alle angegebenen Gelenke
%
% Darstellungsregeln:
%   - Staebe verwenden standardmaessig den Stil mech/rod: schwarz und duenn.
%   - Staebe werden auf dem Layer "mech background" gezeichnet.
%   - Stabnummern besitzen keine Fuellung (fill=none) und bleiben transparent.
%   - Gelenke liegen optisch vor den Staeben und sind mit mechBodyColor gefuellt.
%
% Beispiel:
%   \pic {mech/rod={from=(A),to=(B),label=1,label side=above}};
%   \mechJoints{A,B}
% =========================================================
% -------------------------------------------------------------------------
% Einzelner Fachwerkstab
% -------------------------------------------------------------------------
% "style" erlaubt lokale Abweichungen, ohne den globalen mech/rod-Stil zu
% veraendern. Die Standardbeschriftung liegt an der Stabmitte.
\pgfkeys{
  /mech/rod/.is family,
  /mech/rod,
  from/.initial={(0,0)},
  to/.initial={(1,0)},
  label/.initial=,
  label side/.initial=above,
  label style/.initial={mech/label,fill=none,text=black},
  style/.initial=mech/rod,
}
\tikzset{
  pics/mech/rod/.style args={#1}{
    code={
      \pgfkeys{/mech/rod/.cd,#1}
      \edef\mechRodFrom{\pgfkeysvalueof{/mech/rod/from}}
      \edef\mechRodTo{\pgfkeysvalueof{/mech/rod/to}}
      \coordinate (-start) at \mechRodFrom;
      \coordinate (-end) at \mechRodTo;
      \coordinate (-center) at ($(-start)!.5!(-end)$);
      \begin{pgfonlayer}{mech background}
        \draw[\pgfkeysvalueof{/mech/rod/style}] (-start)--(-end);
      \end{pgfonlayer}
      \node[\pgfkeysvalueof{/mech/rod/label style},\pgfkeysvalueof{/mech/rod/label side}]
        at (-center) {\pgfkeysvalueof{/mech/rod/label}};
    }
  },
  pics/mech/rod/.default={},
}
% -------------------------------------------------------------------------
% Gelenkliste
% -------------------------------------------------------------------------
% Das Makro erwartet eine kommaseparierte Liste bereits definierter TikZ-
% Koordinaten. Jedes Gelenk wird identisch gezeichnet. Die Gelenke sollten nach
% den Staeben aufgerufen werden, damit sie die Stabenden optisch ueberdecken.
\newcommand{\mechJoints}[1]{%
  \foreach \mechJoint in {#1}{%
    \draw[fill=mechBodyColor,draw=black,line width=.2mm]
      (\mechJoint) circle[radius=.6mm];%
  }%
}

% -------------------------------------------------------------------------
% Annotationen und mechanische Symbolgruppen.
% -------------------------------------------------------------------------
% =========================================================
% mechlib/annotations/loads.tex
% LASTANNOTATIONEN: KRAFT, MOMENT UND STRECKENLAST
% =========================================================
% Diese Datei stellt die Lastsymbole der mechlib bereit. Alle Lasten verwenden
% standardmaessig den zentralen Stil mech/load und damit mechLoadColor.
%
% Oeffentliche Pics:
%   mech/force            Einzelkraft als Winkel/Laenge oder from/to
%   mech/moment           Kreisbogensymbol im oder gegen den Uhrzeigersinn
%   mech/distributed-load konstante, dreieckige oder linear veraenderliche Last
%
% Besonderheit der Streckenlast:
% Die Grundlinie wird lokal auf die x-Achse gedreht. Dadurch kann die gesamte
% Geometrie in einem einfachen lokalen Koordinatensystem berechnet werden. Sehr
% kurze Pfeile, deren Laenge kleiner als ihre Pfeilspitze ist, werden bewusst
% unterdrueckt; so entstehen am Nullende einer Dreieckslast keine Artefakte.
% =========================================================
\pgfkeys{
  /mech/force/.is family,
  /mech/force,
  from/.initial={(0,0)},
  to/.initial=,
  angle/.initial=0,
  length/.initial=1.5cm,
  label/.initial=,
  label position/.initial=above,
  direction/.initial=away,
  style/.initial=mech/load,
}
\tikzset{
  pics/mech/force/.style args={#1}{
    code={
      \pgfkeys{/mech/force/.cd,#1}
      \edef\mechForceFrom{\pgfkeysvalueof{/mech/force/from}}
      \edef\mechForceTo{\pgfkeysvalueof{/mech/force/to}}
      \edef\mechForceDirection{\pgfkeysvalueof{/mech/force/direction}}
      \def\mechForceToward{toward}
      \ifx\mechForceTo\empty
        \ifx\mechForceDirection\mechForceToward
          \draw[\pgfkeysvalueof{/mech/force/style},-{Latex[length=2.5mm]}]
            (\pgfkeysvalueof{/mech/force/angle}:\pgfkeysvalueof{/mech/force/length}) -- (0,0)
            node[midway,\pgfkeysvalueof{/mech/force/label position}]
            {\pgfkeysvalueof{/mech/force/label}};
        \else
          \draw[\pgfkeysvalueof{/mech/force/style},-{Latex[length=2.5mm]}]
            (0,0) -- (\pgfkeysvalueof{/mech/force/angle}:\pgfkeysvalueof{/mech/force/length})
            node[midway,\pgfkeysvalueof{/mech/force/label position}]
            {\pgfkeysvalueof{/mech/force/label}};
        \fi
      \else
        \ifx\mechForceDirection\mechForceToward
          \draw[\pgfkeysvalueof{/mech/force/style},-{Latex[length=2.5mm]}]
            \mechForceTo -- \mechForceFrom
            node[midway,\pgfkeysvalueof{/mech/force/label position}]
            {\pgfkeysvalueof{/mech/force/label}};
        \else
          \draw[\pgfkeysvalueof{/mech/force/style},-{Latex[length=2.5mm]}]
            \mechForceFrom -- \mechForceTo
            node[midway,\pgfkeysvalueof{/mech/force/label position}]
            {\pgfkeysvalueof{/mech/force/label}};
        \fi
      \fi
    }
  },
  pics/mech/force/.default={},
}
\pgfkeys{
  /mech/moment/.is family,
  /mech/moment,
  radius/.initial=6mm,
  label/.initial=,
  direction/.is choice,
  direction/ccw/.code={\def\mechMomentStart{-60}\def\mechMomentEnd{220}},
  direction/cw/.code={\def\mechMomentStart{240}\def\mechMomentEnd{-40}},
  direction=ccw,
  style/.initial=mech/load,
}
\tikzset{
  pics/mech/moment/.style args={#1}{
    code={
      \pgfkeys{/mech/moment/.cd,direction=ccw,#1}
      \draw[\pgfkeysvalueof{/mech/moment/style},-{Latex[length=2.5mm]}]
        (\mechMomentStart:\pgfkeysvalueof{/mech/moment/radius})
        arc[start angle=\mechMomentStart,end angle=\mechMomentEnd,
            radius=\pgfkeysvalueof{/mech/moment/radius}];
      \node[mech/label,above] at (0,\pgfkeysvalueof{/mech/moment/radius})
        {\pgfkeysvalueof{/mech/moment/label}};
    }
  },
  pics/mech/moment/.default={},
}
\pgfkeys{
  /mech/distributed-load/.is family,
  /mech/distributed-load,
  from/.initial={(0,0)},
  to/.initial={(1,0)},
  start value/.initial=1,
  end value/.initial=1,
  height/.initial=8mm,
  offset/.initial=0mm,
  arrows/.initial=6,
  arrowhead length/.initial=2mm,
  label/.initial=,
  label position/.initial=above,
  style/.initial=mech/load,
}
\tikzset{
  pics/mech/distributed-load/.style args={#1}{
    code={
      \pgfkeys{/mech/distributed-load/.cd,#1}
      \edef\mechDLFrom{\pgfkeysvalueof{/mech/distributed-load/from}}
      \edef\mechDLTo{\pgfkeysvalueof{/mech/distributed-load/to}}
      \path let
        \p1=\mechDLFrom,\p2=\mechDLTo,
        \n1={veclen(\x2-\x1,\y2-\y1)},
        \n2={atan2(\y2-\y1,\x2-\x1)}
      in
        coordinate (-start) at (\x1,\y1)
        coordinate (-end) at (\x2,\y2)
        coordinate (-helper) at (\n2,\n1);
      \path (-helper); \pgfgetlastxy{\mechDLAngle}{\mechDLLength};
      \pgfmathsetmacro{\mechDLStartValue}{\pgfkeysvalueof{/mech/distributed-load/start value}}
      \pgfmathsetmacro{\mechDLEndValue}{\pgfkeysvalueof{/mech/distributed-load/end value}}
      \pgfmathsetmacro{\mechDLArrowCount}{\pgfkeysvalueof{/mech/distributed-load/arrows}}
      \begin{scope}[shift={(-start)},rotate=\mechDLAngle]
        \pgfmathsetlengthmacro{\mechDLOffset}{\pgfkeysvalueof{/mech/distributed-load/offset}}
        \pgfmathsetlengthmacro{\mechDLArrowheadLength}{\pgfkeysvalueof{/mech/distributed-load/arrowhead length}}
        \pgfmathsetlengthmacro{\mechDLStartArrowLength}{abs(\mechDLStartValue)*\pgfkeysvalueof{/mech/distributed-load/height}}

        \fill[mech/load fill]
          (0,\mechDLOffset)
          --(0,{\mechDLOffset+\mechDLStartValue*\pgfkeysvalueof{/mech/distributed-load/height}})
          --(\mechDLLength,{\mechDLOffset+\mechDLEndValue*\pgfkeysvalueof{/mech/distributed-load/height}})
          --(\mechDLLength,\mechDLOffset)
          --cycle;

        % Keep contour and first arrow continuous whenever the first arrow
        % is long enough to accommodate its arrowhead. Otherwise draw only
        % the contour and suppress the too-short arrow.
        \ifdim\mechDLStartArrowLength<\mechDLArrowheadLength
          \draw[\pgfkeysvalueof{/mech/distributed-load/style}]
            (\mechDLLength,{\mechDLOffset+\mechDLEndValue*\pgfkeysvalueof{/mech/distributed-load/height}})
            --(0,{\mechDLOffset+\mechDLStartValue*\pgfkeysvalueof{/mech/distributed-load/height}});
        \else
          \draw[
            \pgfkeysvalueof{/mech/distributed-load/style},
            -{Latex[length=\pgfkeysvalueof{/mech/distributed-load/arrowhead length}]}
          ]
            (\mechDLLength,{\mechDLOffset+\mechDLEndValue*\pgfkeysvalueof{/mech/distributed-load/height}})
            --(0,{\mechDLOffset+\mechDLStartValue*\pgfkeysvalueof{/mech/distributed-load/height}})
            --(0,\mechDLOffset);
        \fi

        \foreach \mechDLIndex in {1,...,\mechDLArrowCount}{
          \pgfmathsetmacro{\mechDLT}{\mechDLIndex/\mechDLArrowCount}
          \pgfmathsetmacro{\mechDLValue}{\mechDLStartValue+(\mechDLEndValue-\mechDLStartValue)*\mechDLT}
          \pgfmathsetlengthmacro{\mechDLArrowLength}{abs(\mechDLValue)*\pgfkeysvalueof{/mech/distributed-load/height}}

          \ifdim\mechDLArrowLength<\mechDLArrowheadLength
            % Suppress arrows that are shorter than their arrowhead.
          \else
            \draw[
              \pgfkeysvalueof{/mech/distributed-load/style},
              -{Latex[length=\pgfkeysvalueof{/mech/distributed-load/arrowhead length}]}
            ]
              ({\mechDLT*\mechDLLength},{\mechDLOffset+\mechDLValue*\pgfkeysvalueof{/mech/distributed-load/height}})
              --({\mechDLT*\mechDLLength},\mechDLOffset);
          \fi
        }

        \node[mech/label,\pgfkeysvalueof{/mech/distributed-load/label position}]
          at (\mechDLLength,{\mechDLOffset+\mechDLEndValue*\pgfkeysvalueof{/mech/distributed-load/height}})
          {\pgfkeysvalueof{/mech/distributed-load/label}};
      \end{scope}
    }
  },
  pics/mech/distributed-load/.default={},
}

% =========================================================
% mechlib/annotations/dimensions.tex
% BEMASSUNGEN: LINEARES MASS UND WINKELMASS
% =========================================================
% Lineare Bemassungen werden geometrisch aus zwei Endpunkten abgeleitet. Der
% Offset wirkt senkrecht auf die Verbindung von "from" nach "to"; damit
% funktioniert dasselbe Pic fuer horizontale, vertikale und schraege Masse.
%
% Oeffentliche Pics:
%   mech/dimension       lineare Bemassung mit Doppelpfeil
%   mech/angle-dimension Kreisbogen zwischen zwei Winkeln
%
% Interne Koordinaten -A und -B sind die senkrecht versetzten Endpunkte der
% Masslinie. Die Vektorrechnung erfolgt in TikZ "let"-Syntax.
% =========================================================
\pgfkeys{
  /mech/dimension/.is family,
  /mech/dimension,
  from/.initial={(0,0)},
  to/.initial={(1,0)},
  label/.initial={$l$},
  offset/.initial=3mm,
  style/.initial=mech/dimension,
}
\tikzset{
  pics/mech/dimension/.style args={#1}{
    code={
      \pgfkeys{/mech/dimension/.cd,#1}
      \edef\mechDimFrom{\pgfkeysvalueof{/mech/dimension/from}}
      \edef\mechDimTo{\pgfkeysvalueof{/mech/dimension/to}}
      \path let
        \p1=\mechDimFrom,\p2=\mechDimTo,
        \n1={veclen(\x2-\x1,\y2-\y1)},
        \n2={-(\y2-\y1)/\n1*\pgfkeysvalueof{/mech/dimension/offset}},
        \n3={(\x2-\x1)/\n1*\pgfkeysvalueof{/mech/dimension/offset}}
      in
        coordinate (-A) at (\x1+\n2,\y1+\n3)
        coordinate (-B) at (\x2+\n2,\y2+\n3);
      \draw[\pgfkeysvalueof{/mech/dimension/style}]
        (-A)--(-B)
        node[midway,fill=white,inner sep=1pt]
        {\pgfkeysvalueof{/mech/dimension/label}};
    }
  },
  pics/mech/dimension/.default={},
}
\pgfkeys{
  /mech/angle-dimension/.is family,
  /mech/angle-dimension,
  start angle/.initial=0,
  end angle/.initial=45,
  radius/.initial=7mm,
  label/.initial={$\alpha$},
}
\tikzset{
  pics/mech/angle-dimension/.style args={#1}{
    code={
      \pgfkeys{/mech/angle-dimension/.cd,#1}
      \draw
        (\pgfkeysvalueof{/mech/angle-dimension/start angle}:\pgfkeysvalueof{/mech/angle-dimension/radius})
        arc[start angle=\pgfkeysvalueof{/mech/angle-dimension/start angle},
            end angle=\pgfkeysvalueof{/mech/angle-dimension/end angle},
            radius=\pgfkeysvalueof{/mech/angle-dimension/radius}];
      \pgfmathsetmacro{\mechAngleMid}{.5*(\pgfkeysvalueof{/mech/angle-dimension/start angle}+\pgfkeysvalueof{/mech/angle-dimension/end angle})}
      \node at (\mechAngleMid:{.7*\pgfkeysvalueof{/mech/angle-dimension/radius}})
        {\pgfkeysvalueof{/mech/angle-dimension/label}};
    }
  },
  pics/mech/angle-dimension/.default={},
}

% =========================================================
% mechlib/annotations/coordinates.tex
% KOORDINATENANNOTATIONEN: EINZELACHSE UND 2D-KOORDINATENSYSTEM
% =========================================================
% Die Pics dieser Datei zeichnen mathematische Bezugsachsen im Stil mech/axis.
% Winkel werden in TikZ-Grad angegeben; Laengen sind regulaere TeX-Dimensionen.
%
% Oeffentliche Pics:
%   mech/axis              einzelne gerichtete Achse
%   mech/coordinate-system zwei unabhaengig orientierbare Achsen
%
% Beide Pics exportieren den lokalen Ursprung als "-origin". Bei einem benannten
% Pic (cs) kann dieser z. B. als (cs-origin) referenziert werden.
% =========================================================
\pgfkeys{
  /mech/axis/.is family,
  /mech/axis,
  label/.initial={$x$},
  angle/.initial=0,
  length/.initial=1cm,
  direction/.initial=positive,
}
\tikzset{
  pics/mech/axis/.style args={#1}{
    code={
      \pgfkeys{/mech/axis/.cd,#1}
      \edef\mechAxisDirection{\pgfkeysvalueof{/mech/axis/direction}}
      \def\mechAxisNegative{negative}
      \ifx\mechAxisDirection\mechAxisNegative
        \draw[mech/axis,{Latex[length=2mm]}-]
          (0,0)--(\pgfkeysvalueof{/mech/axis/angle}:\pgfkeysvalueof{/mech/axis/length})
          node[pos=1,anchor=west]{\pgfkeysvalueof{/mech/axis/label}};
      \else
        \draw[mech/axis,-{Latex[length=2mm]}]
          (0,0)--(\pgfkeysvalueof{/mech/axis/angle}:\pgfkeysvalueof{/mech/axis/length})
          node[pos=1,anchor=west]{\pgfkeysvalueof{/mech/axis/label}};
      \fi
      \coordinate (-origin) at (0,0);
    }
  },
  pics/mech/axis/.default={},
}
\pgfkeys{
  /mech/coordinate-system/.is family,
  /mech/coordinate-system,
  x label/.initial={$x$},
  y label/.initial={$y$},
  x angle/.initial=0,
  y angle/.initial=90,
  x length/.initial=1cm,
  y length/.initial=1cm,
}
\tikzset{
  pics/mech/coordinate-system/.style args={#1}{
    code={
      \pgfkeys{/mech/coordinate-system/.cd,#1}
      \draw[mech/axis,-{Latex[length=2mm]}]
        (0,0)--(\pgfkeysvalueof{/mech/coordinate-system/x angle}:\pgfkeysvalueof{/mech/coordinate-system/x length})
        node[pos=1,anchor=west]{\pgfkeysvalueof{/mech/coordinate-system/x label}};
      \draw[mech/axis,-{Latex[length=2mm]}]
        (0,0)--(\pgfkeysvalueof{/mech/coordinate-system/y angle}:\pgfkeysvalueof{/mech/coordinate-system/y length})
        node[pos=1,anchor=east]{\pgfkeysvalueof{/mech/coordinate-system/y label}};
      \coordinate (-origin) at (0,0);
    }
  },
  pics/mech/coordinate-system/.default={},
}

% =========================================================
% mechlib/annotations/mechanics.tex
% MECHANISCHE ANNOTATIONEN: SCHNITTGROESSEN UND LAGERREAKTIONEN
% =========================================================
% Die Pics kombinieren mehrere elementare Lastsymbole zu haeufig benoetigten
% mechanischen Symbolgruppen. Farben und Pfeilgestaltung folgen mech/load.
%
% Oeffentliche Pics:
%   mech/internal-forces  Normalkraft N, Querkraft Q und Moment M am Schnitt
%   mech/fixed-reactions  horizontale/vertikale Reaktion und Einspannmoment
%
% Bei mech/internal-forces steuert "side=right|left" die Vorzeichenrichtung
% aller drei Schnittgroessen konsistent. Das Moment wird dazu intern als
% mech/moment mit passender Drehrichtung wiederverwendet.
% =========================================================
\pgfkeys{
  /mech/internal-forces/.is family,
  /mech/internal-forces,
  normal/.initial={$N$},
  shear/.initial={$Q$},
  moment/.initial={$M$},
  length/.initial=7mm,
  side/.is choice,
  side/right/.code={\def\mechIFSign{1}},
  side/left/.code={\def\mechIFSign{-1}},
  side=right,
}
\tikzset{
  pics/mech/internal-forces/.style args={#1}{
    code={
      \pgfkeys{/mech/internal-forces/.cd,side=right,#1}
      \edef\mechIFLength{\pgfkeysvalueof{/mech/internal-forces/length}}
      \draw[mech/load,-{Latex[length=2mm]}]
        (0,0)--({\mechIFSign*\mechIFLength},0)
        node[pos=1,anchor=west]{\pgfkeysvalueof{/mech/internal-forces/normal}};
      \draw[mech/load,-{Latex[length=2mm]}]
        (0,0)--(0,{-.8*\mechIFSign*\mechIFLength})
        node[pos=1,anchor=west]{\pgfkeysvalueof{/mech/internal-forces/shear}};
      \ifnum\mechIFSign=1
        \pic {mech/moment={direction=ccw,radius=.55\mechIFLength,label=\pgfkeysvalueof{/mech/internal-forces/moment}}};
      \else
        \pic {mech/moment={direction=cw,radius=.55\mechIFLength,label=\pgfkeysvalueof{/mech/internal-forces/moment}}};
      \fi
    }
  },
  pics/mech/internal-forces/.default={},
}
\pgfkeys{
  /mech/fixed-reactions/.is family,
  /mech/fixed-reactions,
  x/.initial={$A_x$},
  y/.initial={$A_y$},
  moment/.initial={$M_A$},
  length/.initial=7mm,
}
\tikzset{
  pics/mech/fixed-reactions/.style args={#1}{
    code={
      \pgfkeys{/mech/fixed-reactions/.cd,#1}
      \edef\mechReactionLength{\pgfkeysvalueof{/mech/fixed-reactions/length}}
      \draw[mech/load,-{Latex[length=2mm]}] (-\mechReactionLength,0)--(0,0)
        node[pos=0,anchor=east]{\pgfkeysvalueof{/mech/fixed-reactions/x}};
      \draw[mech/load,-{Latex[length=2mm]}] (0,.7\mechReactionLength)--(0,-.3\mechReactionLength)
        node[pos=0,anchor=south]{\pgfkeysvalueof{/mech/fixed-reactions/y}};
      \pic {mech/moment={direction=ccw,radius=.45\mechReactionLength,label=\pgfkeysvalueof{/mech/fixed-reactions/moment}}};
    }
  },
  pics/mech/fixed-reactions/.default={},
}


%
% Lade-Reihenfolge:
%   1. core/        Farben, Layer, Stile und globale Geometriekonventionen
%   2. elements/    mechanische Grundelemente (Koerper, Lager, Verbinder, Staebe)
%   3. annotations/ Lasten, Masse, Koordinaten und mechanische Annotationen
%
% WICHTIG:
% Die Modulreihenfolge ist absichtlich festgelegt. Elemente greifen auf Farben,
% Layer und Stile aus core/ zurueck; Annotationen koennen wiederum die bereits
% definierten Grundstile und Pics verwenden.
% =========================================================
% -------------------------------------------------------------------------
% Externe Pakete und TikZ-Bibliotheken
% -------------------------------------------------------------------------
% \RequirePackage ist bewusst statt \usepackage gewaehlt, da mechlib.tex selbst
% wie eine Bibliotheks-/Paketdatei verwendet wird.
\RequirePackage{tikz}
\RequirePackage{amsmath}
\RequirePackage{siunitx}
\RequirePackage{xcolor}
\usetikzlibrary{
  arrows.meta,
  calc,
  patterns,
  patterns.meta,
  intersections,
  angles,
  quotes,
  decorations.pathmorphing,
  backgrounds
}
% Base path: override before loading when mechlib is stored elsewhere.
% -------------------------------------------------------------------------
% Aufloesung des Bibliothekspfades
% -------------------------------------------------------------------------
% \providecommand respektiert eine bereits im Hauptdokument gesetzte Definition.
\providecommand{\mechlibpath}{mechlib}
% -------------------------------------------------------------------------
% Kernmodule: muessen vor allen Elementen geladen werden.
% -------------------------------------------------------------------------
% =========================================================
% mechlib/core/colors.tex
% ZENTRALE FARBPALETTE
% =========================================================
% Alle bibliotheksweiten Farben werden ausschliesslich hier definiert. Dadurch
% kann der visuelle Charakter der gesamten mechlib an einer Stelle angepasst
% werden, ohne einzelne Elemente zu veraendern.
%
% Konvention:
%   mechBodyColor    Fuellung allgemeiner Koerper und Fachwerkgelenke
%   mechSupportColor Fuellung von Lagerkoerpern
%   mechLoadColor    Kraefte, Momente und Streckenlasten
%   mechAxisColor    Achsen und Koordinatensysteme
%   mechRodColor     historische Stabfarbe; aktuelle Fachwerkstaebe sind schwarz
%
% Die Namen sind Teil der internen visuellen API und sollten bevorzugt statt
% harter RGB-Werte in den Elementdateien verwendet werden.
% =========================================================
% -------------------------------------------------------------------------
% Farbdefinitionen
% -------------------------------------------------------------------------
% RGB wird bewusst explizit verwendet, damit die Darstellung unabhaengig von
% Dokumentklassen und externen Farbmodellen reproduzierbar bleibt.
\definecolor{mechBodyColor}{RGB}{180,180,180}
\definecolor{mechSupportColor}{RGB}{205,205,205}
\definecolor{mechLoadColor}{RGB}{178,25,25}
\definecolor{mechAxisColor}{RGB}{6,13,141}
\definecolor{mechRodColor}{RGB}{130,130,130}

% =========================================================
% mechlib/core/layers.tex
% ZEICHENEBENEN / PGF-LAYER
% =========================================================
% Fachwerkstaebe und andere hinterlegte Geometrien koennen auf dem Layer
% "mech background" gezeichnet werden. Der normale TikZ-Layer "main" liegt
% darueber. Dadurch koennen z. B. gefuellte Gelenke sauber vor den Staeben
% erscheinen, ohne dass die Zeichenreihenfolge im Anwendungsdokument aufwendig
% kontrolliert werden muss.
% =========================================================
% Zuerst wird der benannte Hintergrundlayer deklariert. Anschliessend wird die
% globale Layerreihenfolge explizit gesetzt: Hintergrund vor Hauptlayer.
\pgfdeclarelayer{mech background}
\pgfsetlayers{mech background,main}

% =========================================================
% mechlib/core/styles.tex
% ZENTRALE, NAMENSRAUMGESCHUETZTE TIKZ-STILE
% =========================================================
% Diese Datei definiert ausschliesslich visuelle Grundstile. Die eigentliche
% Geometrie eines Elements gehoert in elements/ bzw. annotations/.
%
% Alle oeffentlichen Bibliotheksstile beginnen mit "mech/". Das reduziert
% Namenskollisionen mit dem Hauptdokument und macht im TikZ-Code sichtbar,
% welche Formatierung aus der mechlib stammt.
%
% Designregel:
%   - Farben werden aus core/colors.tex bezogen.
%   - Geometrien werden NICHT hier definiert.
%   - Linienstarken werden hier zentral gehalten, soweit sie elementuebergreifend
%     konsistent sein sollen.
% =========================================================
% -------------------------------------------------------------------------
% Visuelle Basisstile
% -------------------------------------------------------------------------
% Einzelne Pics duerfen diese Stile ueber ihre jeweiligen "style"-Keys ersetzen.
\tikzset{
  mech/body/.style={
    draw=black,
    fill=mechBodyColor,
    line width=.6pt
  },
  mech/support/.style={
    draw=black,
    fill=mechSupportColor,
    line width=.6pt
  },
  mech/connector/.style={
    draw=black,
    line width=.6pt
  },
  mech/rod/.style={
    draw=black,
    line width=.85pt
  },
  mech/load/.style={
    draw=mechLoadColor,
    text=mechLoadColor,
    line width=1pt
  },
  mech/load fill/.style={
    fill=mechLoadColor,
    fill opacity=.16,
    draw=none
  },
  mech/dimension/.style={
    {Latex[length=2mm]}-{Latex[length=2mm]},
    line width=.5pt
  },
  mech/axis/.style={
    draw=mechAxisColor,
    text=mechAxisColor,
    line width=.6pt
  },
  mech/label/.style={
    inner sep=1pt
  }
}

% =========================================================
% mechlib/core/defaults.tex
% GLOBALE STANDARDWERTE UND GEMEINSAME GEOMETRIEKONVENTIONEN
% =========================================================
% Hier stehen Werte, die von mehreren Modulen gemeinsam benutzt werden. Solche
% Konstanten verhindern, dass geometrisch zusammengehoerige Elemente lokal mit
% leicht unterschiedlichen Zahlen definiert werden.
%
% Beispiel: \mechConnectionOffset legt die vertikale Lage paralleler
% Anschlussanker fest. Masse und Einspannung verwenden exakt denselben Wert;
% dadurch liegen obere/mittlere/untere Anschluesse bei gleicher Mittelpunktlage
% garantiert auf gleicher Hoehe.
% =========================================================
% Globale Schriftvorgabe fuer jedes TikZ-Bild. "append style" ergaenzt
% vorhandene Bildstile, statt sie zu ueberschreiben.
\tikzset{
  every picture/.append style={font=\fontsize{10}{10}\selectfont},
}

% Standardwinkel fuer technische Bodenschraffuren.
\def\mechHatchAngle{45}

% Shared vertical distance of parallel connector anchors from the middle anchor.
% Wall and mass use this exact value so upper/middle/lower are aligned whenever
% their center points lie on the same horizontal line.
% Gemeinsamer vertikaler Versatz fuer parallele Anschlussanker.
\def\mechConnectionOffset{2.5mm}

% -------------------------------------------------------------------------
% Atomare mechanische Elemente.
% -------------------------------------------------------------------------
% =========================================================
% mechlib/elements/bodies.tex
% ATOMARE KOERPER: MASSE, PUNKTMASSE, WAGEN UND BALKEN
% =========================================================
% Die hier definierten TikZ-Pics stellen mechanische Grundkoerper bereit. Jedes
% Pic besitzt eine eigene pgfkeys-Familie unter /mech/... und kann dadurch mit
% benannten Optionen konfiguriert werden.
%
% Grundprinzip der API:
%   \pic (name) at (<Ort>) {mech/<element>={<key>=<wert>,...}};
%
% Benannte Pics exportieren Koordinaten bzw. werden bei node-basierten Koerpern
% auf den internen Shape-Knoten aliasiert. So koennen regulaere TikZ-Anker wie
% (m1.east), (m1.north west) usw. direkt benutzt werden.
%
% Oeffentliche Pics in dieser Datei:
%   mech/mass       rechteckiger Starrkoerper bzw. konzentrierte Masse
%   mech/point-mass kreisfoermige Punktmasse
%   mech/cart       Wagen mit zwei Raedern und optionalem Boden
%   mech/beam       Balken zwischen zwei beliebigen Punkten
% =========================================================
\pgfkeys{
  /mech/mass/.is family,
  /mech/mass,
  width/.initial=1.5cm,
  height/.initial=1cm,
  label/.initial=,
  style/.initial=mech/body,
}
\tikzset{
  pics/mech/mass/.style args={#1}{
    code={
      \pgfkeys{/mech/mass/.cd,#1}
      \node[
        \pgfkeysvalueof{/mech/mass/style},
        minimum width=\pgfkeysvalueof{/mech/mass/width},
        minimum height=\pgfkeysvalueof{/mech/mass/height},
        inner sep=0pt
      ] (-shape) at (0,0) {\pgfkeysvalueof{/mech/mass/label}};

      % Named body pics behave like regular TikZ nodes.
      % Example: \pic (m1) ... then (m1.east), (m1.south east), ...
      \edef\mechPicName{\pgfkeysvalueof{/tikz/name prefix}}
      \ifx\mechPicName\pgfutil@empty\else
        \pgfnodealias{\mechPicName}{\mechPicName-shape}
      \fi

      \coordinate (-center) at (-shape.center);

      % Three standardized connection levels on each vertical side.
      % The shared \mechConnectionOffset is also used by fixed supports.
      \coordinate (-west-upper)  at ($(-shape.west)+(0,\mechConnectionOffset)$);
      \coordinate (-west-middle) at (-shape.west);
      \coordinate (-west-lower)  at ($(-shape.west)+(0,-\mechConnectionOffset)$);

      \coordinate (-east-upper)  at ($(-shape.east)+(0,\mechConnectionOffset)$);
      \coordinate (-east-middle) at (-shape.east);
      \coordinate (-east-lower)  at ($(-shape.east)+(0,-\mechConnectionOffset)$);

      % Generic connection aliases
      \coordinate (-connection-upper)  at (-east-upper);
      \coordinate (-connection-middle) at (-east-middle);
      \coordinate (-connection-lower)  at (-east-lower);

      \coordinate (-north) at (-shape.north);
      \coordinate (-south) at (-shape.south);
    }
  },
  pics/mech/mass/.default={},
}
\pgfkeys{
  /mech/point-mass/.is family,
  /mech/point-mass,
  radius/.initial=2.2mm,
  label/.initial=,
  style/.initial=mech/body,
}
\tikzset{
  pics/mech/point-mass/.style args={#1}{
    code={
      \pgfkeys{/mech/point-mass/.cd,#1}
      \edef\mechPointMassRadius{\pgfkeysvalueof{/mech/point-mass/radius}}
      \edef\mechPointMassLabel{\pgfkeysvalueof{/mech/point-mass/label}}
      \edef\mechPointMassStyle{\pgfkeysvalueof{/mech/point-mass/style}}
      \node[
        \mechPointMassStyle,
        circle,
        minimum size=2\mechPointMassRadius,
        inner sep=0pt
      ] (-shape) at (0,0) {\mechPointMassLabel};

      \edef\mechPicName{\pgfkeysvalueof{/tikz/name prefix}}
      \ifx\mechPicName\pgfutil@empty\else
        \pgfnodealias{\mechPicName}{\mechPicName-shape}
      \fi

      \coordinate (-center) at (-shape.center);
    }
  },
  pics/mech/point-mass/.default={},
}
% -------------------------------------------------------------------------
% Wagen
% -------------------------------------------------------------------------
% Der Boolean "ground" wird als TeX-if gespeichert. Dadurch kann die optionale
% Bodendarstellung ohne Stringvergleich ein- bzw. ausgeschaltet werden.
\newif\ifmechCartGround
\pgfkeys{
  /mech/cart/.is family,
  /mech/cart,
  width/.initial=2cm,
  height/.initial=1cm,
  wheel radius/.initial=2mm,
  label/.initial=,
  ground/.is if=mechCartGround,
  ground/.default=true,
  ground=true,
}
\tikzset{
  pics/mech/cart/.style args={#1}{
    code={
      \pgfkeys{/mech/cart/.cd,#1}
      \edef\mechCartWidth{\pgfkeysvalueof{/mech/cart/width}}
      \edef\mechCartHeight{\pgfkeysvalueof{/mech/cart/height}}
      \edef\mechCartWheelRadius{\pgfkeysvalueof{/mech/cart/wheel radius}}
      \node[
        mech/body,
        minimum width=\mechCartWidth,
        minimum height=\mechCartHeight,
        inner sep=0pt
      ] (-shape) at (0,0) {\pgfkeysvalueof{/mech/cart/label}};

      \edef\mechPicName{\pgfkeysvalueof{/tikz/name prefix}}
      \ifx\mechPicName\pgfutil@empty\else
        \pgfnodealias{\mechPicName}{\mechPicName-shape}
      \fi

      \coordinate (-center) at (-shape.center);
      \coordinate (-west) at (-shape.west);
      \coordinate (-east) at (-shape.east);
      \coordinate (-north) at (-shape.north);
      \coordinate (-south) at (-shape.south);
      \draw[mech/body]
        (-.3\mechCartWidth,-.5\mechCartHeight-\mechCartWheelRadius)
        circle[radius=\mechCartWheelRadius];
      \draw[mech/body]
        (.3\mechCartWidth,-.5\mechCartHeight-\mechCartWheelRadius)
        circle[radius=\mechCartWheelRadius];
      \ifmechCartGround
        \draw
          (-.5\mechCartWidth,-.5\mechCartHeight-2\mechCartWheelRadius)
          --
          (.5\mechCartWidth,-.5\mechCartHeight-2\mechCartWheelRadius);
        \fill[
          pattern={Lines[angle=\mechHatchAngle,distance=1pt]},
          pattern color=black
        ]
          (-.5\mechCartWidth,-.5\mechCartHeight-2\mechCartWheelRadius)
          rectangle
          (.5\mechCartWidth,-.5\mechCartHeight-2\mechCartWheelRadius-3pt);
      \fi
    }
  },
  pics/mech/cart/.default={},
}
\pgfkeys{
  /mech/beam/.is family,
  /mech/beam,
  from/.initial={(0,0)},
  to/.initial={(1,0)},
  height/.initial=3mm,
  style/.initial=mech/body,
}
\tikzset{
  pics/mech/beam/.style args={#1}{
    code={
      \pgfkeys{/mech/beam/.cd,#1}
      \edef\mechBeamFrom{\pgfkeysvalueof{/mech/beam/from}}
      \edef\mechBeamTo{\pgfkeysvalueof{/mech/beam/to}}
      \edef\mechBeamHeight{\pgfkeysvalueof{/mech/beam/height}}
      \edef\mechBeamStyle{\pgfkeysvalueof{/mech/beam/style}}
      \path let
        \p1=\mechBeamFrom,
        \p2=\mechBeamTo,
        \n1={veclen(\x2-\x1,\y2-\y1)},
        \n2={(\y2-\y1)/\n1*\mechBeamHeight/2},
        \n3={-(\x2-\x1)/\n1*\mechBeamHeight/2}
      in
        coordinate (-A) at (\x1+\n2,\y1+\n3)
        coordinate (-B) at (\x2+\n2,\y2+\n3)
        coordinate (-C) at (\x2-\n2,\y2-\n3)
        coordinate (-D) at (\x1-\n2,\y1-\n3)
        coordinate (-start) at (\x1,\y1)
        coordinate (-end) at (\x2,\y2)
        coordinate (-center) at ({(\x1+\x2)/2},{(\y1+\y2)/2});
      \draw[\mechBeamStyle] (-A)--(-B)--(-C)--(-D)--cycle;
    }
  },
  pics/mech/beam/.default={},
}

% =========================================================
% mechlib/elements/supports.tex
% ATOMARE LAGER: FESTLAGER, LOSLAGER UND EINSPANNUNG
% =========================================================
% Diese Datei kapselt die Lagergeometrie. Fest- und Loslager teilen sich die
% Key-Familie /mech/support-common, damit Breite, Hoehe, Schraffur und Label
% konsistent parametrisiert werden koennen.
%
% Oeffentliche Pics:
%   mech/pinned-support Festlager / gelenkiges Lager
%   mech/roller-support Loslager mit sichtbarem Abstand zum Boden
%   mech/fixed-support  Einspannung, optional gedreht
%
% Aktuelle Fachwerk-Konvention:
%   - Die Dreiecksspitze liegt exakt im Mittelpunkt des Gelenks (0,0).
%   - Das Dreieck ist gegenueber der frueheren Version kompakt skaliert.
%   - Bodenlinie und Schraffur verwenden dieselben x-Grenzen wie die
%     Dreiecksgrundseite; beim Loslager wird die sichtbare Linienbreite explizit
%     beruecksichtigt.
%   - Das Gelenk selbst ist mit mechBodyColor gefuellt.
%
% Die pgfmathsetlengthmacro-Aufrufe sind absichtlich verwendet: LaTeX-Laengen
% duerfen nicht durch Textersatz wie ".5\\macro" skaliert werden, da dies bei
% expandierenden Dimensionsmakros zu sichtbarem Resttext (z. B. ".8mm") fuehren
% kann. Alle abgeleiteten Abmessungen werden deshalb numerisch berechnet.
% =========================================================
\pgfkeys{
  /mech/support-common/.is family,
  /mech/support-common,
  width/.initial=2.72mm,
  height/.initial=2.24mm,
  hatch angle/.initial=45,
  label/.initial=,
  label position/.initial=above,
}
\tikzset{
  % -----------------------------------------------------------------------
  % Festlager
  % -----------------------------------------------------------------------
  pics/mech/pinned-support/.style args={#1}{
    code={
      \pgfkeys{/mech/support-common/.cd,#1}
      \edef\mechSupportWidth{\pgfkeysvalueof{/mech/support-common/width}}
      \edef\mechSupportHeight{\pgfkeysvalueof{/mech/support-common/height}}
      \edef\mechSupportHatch{\pgfkeysvalueof{/mech/support-common/hatch angle}}
      \pgfmathsetlengthmacro{\mechSupportHalfWidth}{0.5*\mechSupportWidth}
      % Die Dreieckskontur ragt an der Grundseite durch die halbe Linienbreite
      % minimal ueber die geometrische Breite hinaus. Fuer identische sichtbare
      % Breite wird die Bodenlinie daher um 0.3pt je Seite erweitert.
      \pgfmathsetlengthmacro{\mechSupportGroundHalfWidth}{\mechSupportHalfWidth+0.3pt}
      \pgfmathsetlengthmacro{\mechSupportApexY}{0mm}
      \pgfmathsetlengthmacro{\mechSupportHatchDepth}{0.256*\mechSupportHeight}

      \coordinate (-connection) at (0,0);
      \coordinate (-center) at (0,0);

      \draw[mech/support]
        (-\mechSupportHalfWidth,-\mechSupportHeight)
        -- (0,\mechSupportApexY)
        -- (\mechSupportHalfWidth,-\mechSupportHeight)
        -- cycle;
      \draw[mech/body] (0,0) circle[radius=.6mm];
      \draw[black,line width=.6pt,line cap=butt]
        (-\mechSupportGroundHalfWidth,-\mechSupportHeight)
        -- (\mechSupportGroundHalfWidth,-\mechSupportHeight);
      \fill[
        pattern={Lines[angle=\mechSupportHatch,distance=1pt]},
        pattern color=black
      ]
        (-\mechSupportGroundHalfWidth,-\mechSupportHeight)
        rectangle
        (\mechSupportGroundHalfWidth,{-(\mechSupportHeight+\mechSupportHatchDepth)});

      \node[\pgfkeysvalueof{/mech/support-common/label position}]
        at (0,0) {\pgfkeysvalueof{/mech/support-common/label}};
    }
  },
  pics/mech/pinned-support/.default={},
  % -----------------------------------------------------------------------
  % Loslager
  % -----------------------------------------------------------------------
  pics/mech/roller-support/.style args={#1}{
    code={
      \pgfkeys{/mech/support-common/.cd,#1}
      \edef\mechSupportWidth{\pgfkeysvalueof{/mech/support-common/width}}
      \edef\mechSupportHeight{\pgfkeysvalueof{/mech/support-common/height}}
      \edef\mechSupportHatch{\pgfkeysvalueof{/mech/support-common/hatch angle}}
      \pgfmathsetlengthmacro{\mechSupportHalfWidth}{0.5*\mechSupportWidth}
      % The triangle outline extends by half its line width at the base corners.
      % Add 0.3pt so the roller ground has the same visible total width.
      \pgfmathsetlengthmacro{\mechSupportGroundHalfWidth}{\mechSupportHalfWidth+0.3pt}
      \pgfmathsetlengthmacro{\mechSupportApexY}{0mm}
      \pgfmathsetlengthmacro{\mechSupportRollerGap}{0.272*\mechSupportHeight}
      \pgfmathsetlengthmacro{\mechSupportHatchDepth}{0.256*\mechSupportHeight}
      \pgfmathsetlengthmacro{\mechSupportGroundY}{\mechSupportHeight+\mechSupportRollerGap}

      \coordinate (-connection) at (0,0);
      \coordinate (-center) at (0,0);

      \draw[mech/support]
        (-\mechSupportHalfWidth,-\mechSupportHeight)
        -- (0,\mechSupportApexY)
        -- (\mechSupportHalfWidth,-\mechSupportHeight)
        -- cycle;
      \draw[mech/body] (0,0) circle[radius=.6mm];
      \draw[black,line width=.6pt,line cap=butt]
        (-\mechSupportGroundHalfWidth,-\mechSupportGroundY)
        -- (\mechSupportGroundHalfWidth,-\mechSupportGroundY);
      \fill[
        pattern={Lines[angle=\mechSupportHatch,distance=1pt]},
        pattern color=black
      ]
        (-\mechSupportGroundHalfWidth,-\mechSupportGroundY)
        rectangle
        (\mechSupportGroundHalfWidth,{-(\mechSupportGroundY+\mechSupportHatchDepth)});

      \node[\pgfkeysvalueof{/mech/support-common/label position}]
        at (0,0) {\pgfkeysvalueof{/mech/support-common/label}};
    }
  },
  pics/mech/roller-support/.default={},
}
\pgfkeys{
  /mech/fixed-support/.is family,
  /mech/fixed-support,
  width/.initial=1cm,
  depth/.initial=1.5mm,
  hatch angle/.initial=45,
  angle/.initial=0,
  label/.initial=,
}
\tikzset{
  pics/mech/fixed-support/.style args={#1}{
    code={
      \pgfkeys{/mech/fixed-support/.cd,#1}
      \edef\mechFixedWidth{\pgfkeysvalueof{/mech/fixed-support/width}}
      \edef\mechFixedDepth{\pgfkeysvalueof{/mech/fixed-support/depth}}
      \edef\mechFixedHatch{\pgfkeysvalueof{/mech/fixed-support/hatch angle}}
      \edef\mechFixedAngle{\pgfkeysvalueof{/mech/fixed-support/angle}}

      \coordinate (-center) at (0,0);
      \coordinate (-connection-upper)  at (0,\mechConnectionOffset);
      \coordinate (-connection-middle) at (0,0);
      \coordinate (-connection-lower)  at (0,-\mechConnectionOffset);

      \begin{scope}[rotate=\mechFixedAngle]
        \draw[mech/body,fill=none] (-.5\mechFixedWidth,0)--(.5\mechFixedWidth,0);
        \fill[
          pattern={Lines[angle=\mechFixedHatch,distance={3pt/sqrt(2)}]},
          pattern color=black
        ]
          (-.5\mechFixedWidth,0)
          rectangle
          (.5\mechFixedWidth,\mechFixedDepth);
      \end{scope}

      \node[above] at (0,0) {\pgfkeysvalueof{/mech/fixed-support/label}};
    }
  },
  pics/mech/fixed-support/.default={},
}

% =========================================================
% mechlib/elements/connectors.tex
% ATOMARE VERBINDUNGSELEMENTE: FEDER UND DAEMPFER
% =========================================================
% Beide Elemente folgen derselben grundlegenden API: "from", "to" und "label".
% Die Endpunkte koennen beliebige TikZ-Koordinaten oder exportierte Anker aus
% anderen mechlib-Pics sein.
%
% Oeffentliche Pics:
%   mech/spring  Zickzack-Feder entlang der Verbindungslinie
%   mech/damper  Daempfer mit rechteckigem Daempferkoerper in Linienmitte
%
% Die Beschriftung wird "sloped" gezeichnet und folgt daher der Orientierung
% des Verbindungselements. Ueber "label side" kann sie ober- oder unterhalb
% der lokalen Verbindungslinie platziert werden.
% =========================================================
\pgfkeys{
  /mech/spring/.is family,
  /mech/spring,
  from/.initial={(0,0)},
  to/.initial={(1,0)},
  label/.initial=,
  label side/.initial=above,
  amplitude/.initial=2mm,
  segment length/.initial=4mm,
  style/.initial=mech/connector,
}
\tikzset{
  pics/mech/spring/.style args={#1}{
    code={
      \pgfkeys{/mech/spring/.cd,#1}
      \edef\mechSpringFrom{\pgfkeysvalueof{/mech/spring/from}}
      \edef\mechSpringTo{\pgfkeysvalueof{/mech/spring/to}}
      \draw[\pgfkeysvalueof{/mech/spring/style},
        decorate,
        decoration={zigzag,
          segment length=\pgfkeysvalueof{/mech/spring/segment length},
          amplitude=\pgfkeysvalueof{/mech/spring/amplitude}}]
        \mechSpringFrom -- \mechSpringTo
        node[midway,sloped,mech/label,outer sep=2mm,\pgfkeysvalueof{/mech/spring/label side}]
        {\pgfkeysvalueof{/mech/spring/label}};
    }
  },
  pics/mech/spring/.default={},
}
\pgfkeys{
  /mech/damper/.is family,
  /mech/damper,
  from/.initial={(0,0)},
  to/.initial={(1,0)},
  label/.initial=,
  label side/.initial=below,
  body width/.initial=8mm,
  body height/.initial=3mm,
  style/.initial=mech/connector,
}
\tikzset{
  pics/mech/damper/.style args={#1}{
    code={
      \pgfkeys{/mech/damper/.cd,#1}
      \edef\mechDamperFrom{\pgfkeysvalueof{/mech/damper/from}}
      \edef\mechDamperTo{\pgfkeysvalueof{/mech/damper/to}}
      \draw[\pgfkeysvalueof{/mech/damper/style}]
        \mechDamperFrom -- \mechDamperTo
        node[midway,sloped,draw,fill=white,
          minimum width=\pgfkeysvalueof{/mech/damper/body width},
          minimum height=\pgfkeysvalueof{/mech/damper/body height},
          inner sep=0pt] {};
      \path \mechDamperFrom -- \mechDamperTo
        node[midway,sloped,mech/label,outer sep=2mm,\pgfkeysvalueof{/mech/damper/label side}]
        {\pgfkeysvalueof{/mech/damper/label}};
    }
  },
  pics/mech/damper/.default={},
}

% =========================================================
% mechlib/elements/trusses.tex
% FACHWERKE: STAB UND GEMEINSAME GELENKDARSTELLUNG
% =========================================================
% Dieses Modul enthaelt die atomaren Fachwerkelemente. Die Topologie eines
% Fachwerks bleibt bewusst im Anwendungsdokument: Knoten werden als TikZ-
% Koordinaten definiert und durch mech/rod bzw. \mechJoints verbunden.
%
% Oeffentliche API:
%   mech/rod     Stab zwischen "from" und "to" mit optionalem Label
%   \mechJoints{A,B,C,...} zeichnet alle angegebenen Gelenke
%
% Darstellungsregeln:
%   - Staebe verwenden standardmaessig den Stil mech/rod: schwarz und duenn.
%   - Staebe werden auf dem Layer "mech background" gezeichnet.
%   - Stabnummern besitzen keine Fuellung (fill=none) und bleiben transparent.
%   - Gelenke liegen optisch vor den Staeben und sind mit mechBodyColor gefuellt.
%
% Beispiel:
%   \pic {mech/rod={from=(A),to=(B),label=1,label side=above}};
%   \mechJoints{A,B}
% =========================================================
% -------------------------------------------------------------------------
% Einzelner Fachwerkstab
% -------------------------------------------------------------------------
% "style" erlaubt lokale Abweichungen, ohne den globalen mech/rod-Stil zu
% veraendern. Die Standardbeschriftung liegt an der Stabmitte.
\pgfkeys{
  /mech/rod/.is family,
  /mech/rod,
  from/.initial={(0,0)},
  to/.initial={(1,0)},
  label/.initial=,
  label side/.initial=above,
  label style/.initial={mech/label,fill=none,text=black},
  style/.initial=mech/rod,
}
\tikzset{
  pics/mech/rod/.style args={#1}{
    code={
      \pgfkeys{/mech/rod/.cd,#1}
      \edef\mechRodFrom{\pgfkeysvalueof{/mech/rod/from}}
      \edef\mechRodTo{\pgfkeysvalueof{/mech/rod/to}}
      \coordinate (-start) at \mechRodFrom;
      \coordinate (-end) at \mechRodTo;
      \coordinate (-center) at ($(-start)!.5!(-end)$);
      \begin{pgfonlayer}{mech background}
        \draw[\pgfkeysvalueof{/mech/rod/style}] (-start)--(-end);
      \end{pgfonlayer}
      \node[\pgfkeysvalueof{/mech/rod/label style},\pgfkeysvalueof{/mech/rod/label side}]
        at (-center) {\pgfkeysvalueof{/mech/rod/label}};
    }
  },
  pics/mech/rod/.default={},
}
% -------------------------------------------------------------------------
% Gelenkliste
% -------------------------------------------------------------------------
% Das Makro erwartet eine kommaseparierte Liste bereits definierter TikZ-
% Koordinaten. Jedes Gelenk wird identisch gezeichnet. Die Gelenke sollten nach
% den Staeben aufgerufen werden, damit sie die Stabenden optisch ueberdecken.
\newcommand{\mechJoints}[1]{%
  \foreach \mechJoint in {#1}{%
    \draw[fill=mechBodyColor,draw=black,line width=.2mm]
      (\mechJoint) circle[radius=.6mm];%
  }%
}

% -------------------------------------------------------------------------
% Annotationen und mechanische Symbolgruppen.
% -------------------------------------------------------------------------
% =========================================================
% mechlib/annotations/loads.tex
% LASTANNOTATIONEN: KRAFT, MOMENT UND STRECKENLAST
% =========================================================
% Diese Datei stellt die Lastsymbole der mechlib bereit. Alle Lasten verwenden
% standardmaessig den zentralen Stil mech/load und damit mechLoadColor.
%
% Oeffentliche Pics:
%   mech/force            Einzelkraft als Winkel/Laenge oder from/to
%   mech/moment           Kreisbogensymbol im oder gegen den Uhrzeigersinn
%   mech/distributed-load konstante, dreieckige oder linear veraenderliche Last
%
% Besonderheit der Streckenlast:
% Die Grundlinie wird lokal auf die x-Achse gedreht. Dadurch kann die gesamte
% Geometrie in einem einfachen lokalen Koordinatensystem berechnet werden. Sehr
% kurze Pfeile, deren Laenge kleiner als ihre Pfeilspitze ist, werden bewusst
% unterdrueckt; so entstehen am Nullende einer Dreieckslast keine Artefakte.
% =========================================================
\pgfkeys{
  /mech/force/.is family,
  /mech/force,
  from/.initial={(0,0)},
  to/.initial=,
  angle/.initial=0,
  length/.initial=1.5cm,
  label/.initial=,
  label position/.initial=above,
  direction/.initial=away,
  style/.initial=mech/load,
}
\tikzset{
  pics/mech/force/.style args={#1}{
    code={
      \pgfkeys{/mech/force/.cd,#1}
      \edef\mechForceFrom{\pgfkeysvalueof{/mech/force/from}}
      \edef\mechForceTo{\pgfkeysvalueof{/mech/force/to}}
      \edef\mechForceDirection{\pgfkeysvalueof{/mech/force/direction}}
      \def\mechForceToward{toward}
      \ifx\mechForceTo\empty
        \ifx\mechForceDirection\mechForceToward
          \draw[\pgfkeysvalueof{/mech/force/style},-{Latex[length=2.5mm]}]
            (\pgfkeysvalueof{/mech/force/angle}:\pgfkeysvalueof{/mech/force/length}) -- (0,0)
            node[midway,\pgfkeysvalueof{/mech/force/label position}]
            {\pgfkeysvalueof{/mech/force/label}};
        \else
          \draw[\pgfkeysvalueof{/mech/force/style},-{Latex[length=2.5mm]}]
            (0,0) -- (\pgfkeysvalueof{/mech/force/angle}:\pgfkeysvalueof{/mech/force/length})
            node[midway,\pgfkeysvalueof{/mech/force/label position}]
            {\pgfkeysvalueof{/mech/force/label}};
        \fi
      \else
        \ifx\mechForceDirection\mechForceToward
          \draw[\pgfkeysvalueof{/mech/force/style},-{Latex[length=2.5mm]}]
            \mechForceTo -- \mechForceFrom
            node[midway,\pgfkeysvalueof{/mech/force/label position}]
            {\pgfkeysvalueof{/mech/force/label}};
        \else
          \draw[\pgfkeysvalueof{/mech/force/style},-{Latex[length=2.5mm]}]
            \mechForceFrom -- \mechForceTo
            node[midway,\pgfkeysvalueof{/mech/force/label position}]
            {\pgfkeysvalueof{/mech/force/label}};
        \fi
      \fi
    }
  },
  pics/mech/force/.default={},
}
\pgfkeys{
  /mech/moment/.is family,
  /mech/moment,
  radius/.initial=6mm,
  label/.initial=,
  direction/.is choice,
  direction/ccw/.code={\def\mechMomentStart{-60}\def\mechMomentEnd{220}},
  direction/cw/.code={\def\mechMomentStart{240}\def\mechMomentEnd{-40}},
  direction=ccw,
  style/.initial=mech/load,
}
\tikzset{
  pics/mech/moment/.style args={#1}{
    code={
      \pgfkeys{/mech/moment/.cd,direction=ccw,#1}
      \draw[\pgfkeysvalueof{/mech/moment/style},-{Latex[length=2.5mm]}]
        (\mechMomentStart:\pgfkeysvalueof{/mech/moment/radius})
        arc[start angle=\mechMomentStart,end angle=\mechMomentEnd,
            radius=\pgfkeysvalueof{/mech/moment/radius}];
      \node[mech/label,above] at (0,\pgfkeysvalueof{/mech/moment/radius})
        {\pgfkeysvalueof{/mech/moment/label}};
    }
  },
  pics/mech/moment/.default={},
}
\pgfkeys{
  /mech/distributed-load/.is family,
  /mech/distributed-load,
  from/.initial={(0,0)},
  to/.initial={(1,0)},
  start value/.initial=1,
  end value/.initial=1,
  height/.initial=8mm,
  offset/.initial=0mm,
  arrows/.initial=6,
  arrowhead length/.initial=2mm,
  label/.initial=,
  label position/.initial=above,
  style/.initial=mech/load,
}
\tikzset{
  pics/mech/distributed-load/.style args={#1}{
    code={
      \pgfkeys{/mech/distributed-load/.cd,#1}
      \edef\mechDLFrom{\pgfkeysvalueof{/mech/distributed-load/from}}
      \edef\mechDLTo{\pgfkeysvalueof{/mech/distributed-load/to}}
      \path let
        \p1=\mechDLFrom,\p2=\mechDLTo,
        \n1={veclen(\x2-\x1,\y2-\y1)},
        \n2={atan2(\y2-\y1,\x2-\x1)}
      in
        coordinate (-start) at (\x1,\y1)
        coordinate (-end) at (\x2,\y2)
        coordinate (-helper) at (\n2,\n1);
      \path (-helper); \pgfgetlastxy{\mechDLAngle}{\mechDLLength};
      \pgfmathsetmacro{\mechDLStartValue}{\pgfkeysvalueof{/mech/distributed-load/start value}}
      \pgfmathsetmacro{\mechDLEndValue}{\pgfkeysvalueof{/mech/distributed-load/end value}}
      \pgfmathsetmacro{\mechDLArrowCount}{\pgfkeysvalueof{/mech/distributed-load/arrows}}
      \begin{scope}[shift={(-start)},rotate=\mechDLAngle]
        \pgfmathsetlengthmacro{\mechDLOffset}{\pgfkeysvalueof{/mech/distributed-load/offset}}
        \pgfmathsetlengthmacro{\mechDLArrowheadLength}{\pgfkeysvalueof{/mech/distributed-load/arrowhead length}}
        \pgfmathsetlengthmacro{\mechDLStartArrowLength}{abs(\mechDLStartValue)*\pgfkeysvalueof{/mech/distributed-load/height}}

        \fill[mech/load fill]
          (0,\mechDLOffset)
          --(0,{\mechDLOffset+\mechDLStartValue*\pgfkeysvalueof{/mech/distributed-load/height}})
          --(\mechDLLength,{\mechDLOffset+\mechDLEndValue*\pgfkeysvalueof{/mech/distributed-load/height}})
          --(\mechDLLength,\mechDLOffset)
          --cycle;

        % Keep contour and first arrow continuous whenever the first arrow
        % is long enough to accommodate its arrowhead. Otherwise draw only
        % the contour and suppress the too-short arrow.
        \ifdim\mechDLStartArrowLength<\mechDLArrowheadLength
          \draw[\pgfkeysvalueof{/mech/distributed-load/style}]
            (\mechDLLength,{\mechDLOffset+\mechDLEndValue*\pgfkeysvalueof{/mech/distributed-load/height}})
            --(0,{\mechDLOffset+\mechDLStartValue*\pgfkeysvalueof{/mech/distributed-load/height}});
        \else
          \draw[
            \pgfkeysvalueof{/mech/distributed-load/style},
            -{Latex[length=\pgfkeysvalueof{/mech/distributed-load/arrowhead length}]}
          ]
            (\mechDLLength,{\mechDLOffset+\mechDLEndValue*\pgfkeysvalueof{/mech/distributed-load/height}})
            --(0,{\mechDLOffset+\mechDLStartValue*\pgfkeysvalueof{/mech/distributed-load/height}})
            --(0,\mechDLOffset);
        \fi

        \foreach \mechDLIndex in {1,...,\mechDLArrowCount}{
          \pgfmathsetmacro{\mechDLT}{\mechDLIndex/\mechDLArrowCount}
          \pgfmathsetmacro{\mechDLValue}{\mechDLStartValue+(\mechDLEndValue-\mechDLStartValue)*\mechDLT}
          \pgfmathsetlengthmacro{\mechDLArrowLength}{abs(\mechDLValue)*\pgfkeysvalueof{/mech/distributed-load/height}}

          \ifdim\mechDLArrowLength<\mechDLArrowheadLength
            % Suppress arrows that are shorter than their arrowhead.
          \else
            \draw[
              \pgfkeysvalueof{/mech/distributed-load/style},
              -{Latex[length=\pgfkeysvalueof{/mech/distributed-load/arrowhead length}]}
            ]
              ({\mechDLT*\mechDLLength},{\mechDLOffset+\mechDLValue*\pgfkeysvalueof{/mech/distributed-load/height}})
              --({\mechDLT*\mechDLLength},\mechDLOffset);
          \fi
        }

        \node[mech/label,\pgfkeysvalueof{/mech/distributed-load/label position}]
          at (\mechDLLength,{\mechDLOffset+\mechDLEndValue*\pgfkeysvalueof{/mech/distributed-load/height}})
          {\pgfkeysvalueof{/mech/distributed-load/label}};
      \end{scope}
    }
  },
  pics/mech/distributed-load/.default={},
}

% =========================================================
% mechlib/annotations/dimensions.tex
% BEMASSUNGEN: LINEARES MASS UND WINKELMASS
% =========================================================
% Lineare Bemassungen werden geometrisch aus zwei Endpunkten abgeleitet. Der
% Offset wirkt senkrecht auf die Verbindung von "from" nach "to"; damit
% funktioniert dasselbe Pic fuer horizontale, vertikale und schraege Masse.
%
% Oeffentliche Pics:
%   mech/dimension       lineare Bemassung mit Doppelpfeil
%   mech/angle-dimension Kreisbogen zwischen zwei Winkeln
%
% Interne Koordinaten -A und -B sind die senkrecht versetzten Endpunkte der
% Masslinie. Die Vektorrechnung erfolgt in TikZ "let"-Syntax.
% =========================================================
\pgfkeys{
  /mech/dimension/.is family,
  /mech/dimension,
  from/.initial={(0,0)},
  to/.initial={(1,0)},
  label/.initial={$l$},
  offset/.initial=3mm,
  style/.initial=mech/dimension,
}
\tikzset{
  pics/mech/dimension/.style args={#1}{
    code={
      \pgfkeys{/mech/dimension/.cd,#1}
      \edef\mechDimFrom{\pgfkeysvalueof{/mech/dimension/from}}
      \edef\mechDimTo{\pgfkeysvalueof{/mech/dimension/to}}
      \path let
        \p1=\mechDimFrom,\p2=\mechDimTo,
        \n1={veclen(\x2-\x1,\y2-\y1)},
        \n2={-(\y2-\y1)/\n1*\pgfkeysvalueof{/mech/dimension/offset}},
        \n3={(\x2-\x1)/\n1*\pgfkeysvalueof{/mech/dimension/offset}}
      in
        coordinate (-A) at (\x1+\n2,\y1+\n3)
        coordinate (-B) at (\x2+\n2,\y2+\n3);
      \draw[\pgfkeysvalueof{/mech/dimension/style}]
        (-A)--(-B)
        node[midway,fill=white,inner sep=1pt]
        {\pgfkeysvalueof{/mech/dimension/label}};
    }
  },
  pics/mech/dimension/.default={},
}
\pgfkeys{
  /mech/angle-dimension/.is family,
  /mech/angle-dimension,
  start angle/.initial=0,
  end angle/.initial=45,
  radius/.initial=7mm,
  label/.initial={$\alpha$},
}
\tikzset{
  pics/mech/angle-dimension/.style args={#1}{
    code={
      \pgfkeys{/mech/angle-dimension/.cd,#1}
      \draw
        (\pgfkeysvalueof{/mech/angle-dimension/start angle}:\pgfkeysvalueof{/mech/angle-dimension/radius})
        arc[start angle=\pgfkeysvalueof{/mech/angle-dimension/start angle},
            end angle=\pgfkeysvalueof{/mech/angle-dimension/end angle},
            radius=\pgfkeysvalueof{/mech/angle-dimension/radius}];
      \pgfmathsetmacro{\mechAngleMid}{.5*(\pgfkeysvalueof{/mech/angle-dimension/start angle}+\pgfkeysvalueof{/mech/angle-dimension/end angle})}
      \node at (\mechAngleMid:{.7*\pgfkeysvalueof{/mech/angle-dimension/radius}})
        {\pgfkeysvalueof{/mech/angle-dimension/label}};
    }
  },
  pics/mech/angle-dimension/.default={},
}

% =========================================================
% mechlib/annotations/coordinates.tex
% KOORDINATENANNOTATIONEN: EINZELACHSE UND 2D-KOORDINATENSYSTEM
% =========================================================
% Die Pics dieser Datei zeichnen mathematische Bezugsachsen im Stil mech/axis.
% Winkel werden in TikZ-Grad angegeben; Laengen sind regulaere TeX-Dimensionen.
%
% Oeffentliche Pics:
%   mech/axis              einzelne gerichtete Achse
%   mech/coordinate-system zwei unabhaengig orientierbare Achsen
%
% Beide Pics exportieren den lokalen Ursprung als "-origin". Bei einem benannten
% Pic (cs) kann dieser z. B. als (cs-origin) referenziert werden.
% =========================================================
\pgfkeys{
  /mech/axis/.is family,
  /mech/axis,
  label/.initial={$x$},
  angle/.initial=0,
  length/.initial=1cm,
  direction/.initial=positive,
}
\tikzset{
  pics/mech/axis/.style args={#1}{
    code={
      \pgfkeys{/mech/axis/.cd,#1}
      \edef\mechAxisDirection{\pgfkeysvalueof{/mech/axis/direction}}
      \def\mechAxisNegative{negative}
      \ifx\mechAxisDirection\mechAxisNegative
        \draw[mech/axis,{Latex[length=2mm]}-]
          (0,0)--(\pgfkeysvalueof{/mech/axis/angle}:\pgfkeysvalueof{/mech/axis/length})
          node[pos=1,anchor=west]{\pgfkeysvalueof{/mech/axis/label}};
      \else
        \draw[mech/axis,-{Latex[length=2mm]}]
          (0,0)--(\pgfkeysvalueof{/mech/axis/angle}:\pgfkeysvalueof{/mech/axis/length})
          node[pos=1,anchor=west]{\pgfkeysvalueof{/mech/axis/label}};
      \fi
      \coordinate (-origin) at (0,0);
    }
  },
  pics/mech/axis/.default={},
}
\pgfkeys{
  /mech/coordinate-system/.is family,
  /mech/coordinate-system,
  x label/.initial={$x$},
  y label/.initial={$y$},
  x angle/.initial=0,
  y angle/.initial=90,
  x length/.initial=1cm,
  y length/.initial=1cm,
}
\tikzset{
  pics/mech/coordinate-system/.style args={#1}{
    code={
      \pgfkeys{/mech/coordinate-system/.cd,#1}
      \draw[mech/axis,-{Latex[length=2mm]}]
        (0,0)--(\pgfkeysvalueof{/mech/coordinate-system/x angle}:\pgfkeysvalueof{/mech/coordinate-system/x length})
        node[pos=1,anchor=west]{\pgfkeysvalueof{/mech/coordinate-system/x label}};
      \draw[mech/axis,-{Latex[length=2mm]}]
        (0,0)--(\pgfkeysvalueof{/mech/coordinate-system/y angle}:\pgfkeysvalueof{/mech/coordinate-system/y length})
        node[pos=1,anchor=east]{\pgfkeysvalueof{/mech/coordinate-system/y label}};
      \coordinate (-origin) at (0,0);
    }
  },
  pics/mech/coordinate-system/.default={},
}

% =========================================================
% mechlib/annotations/mechanics.tex
% MECHANISCHE ANNOTATIONEN: SCHNITTGROESSEN UND LAGERREAKTIONEN
% =========================================================
% Die Pics kombinieren mehrere elementare Lastsymbole zu haeufig benoetigten
% mechanischen Symbolgruppen. Farben und Pfeilgestaltung folgen mech/load.
%
% Oeffentliche Pics:
%   mech/internal-forces  Normalkraft N, Querkraft Q und Moment M am Schnitt
%   mech/fixed-reactions  horizontale/vertikale Reaktion und Einspannmoment
%
% Bei mech/internal-forces steuert "side=right|left" die Vorzeichenrichtung
% aller drei Schnittgroessen konsistent. Das Moment wird dazu intern als
% mech/moment mit passender Drehrichtung wiederverwendet.
% =========================================================
\pgfkeys{
  /mech/internal-forces/.is family,
  /mech/internal-forces,
  normal/.initial={$N$},
  shear/.initial={$Q$},
  moment/.initial={$M$},
  length/.initial=7mm,
  side/.is choice,
  side/right/.code={\def\mechIFSign{1}},
  side/left/.code={\def\mechIFSign{-1}},
  side=right,
}
\tikzset{
  pics/mech/internal-forces/.style args={#1}{
    code={
      \pgfkeys{/mech/internal-forces/.cd,side=right,#1}
      \edef\mechIFLength{\pgfkeysvalueof{/mech/internal-forces/length}}
      \draw[mech/load,-{Latex[length=2mm]}]
        (0,0)--({\mechIFSign*\mechIFLength},0)
        node[pos=1,anchor=west]{\pgfkeysvalueof{/mech/internal-forces/normal}};
      \draw[mech/load,-{Latex[length=2mm]}]
        (0,0)--(0,{-.8*\mechIFSign*\mechIFLength})
        node[pos=1,anchor=west]{\pgfkeysvalueof{/mech/internal-forces/shear}};
      \ifnum\mechIFSign=1
        \pic {mech/moment={direction=ccw,radius=.55\mechIFLength,label=\pgfkeysvalueof{/mech/internal-forces/moment}}};
      \else
        \pic {mech/moment={direction=cw,radius=.55\mechIFLength,label=\pgfkeysvalueof{/mech/internal-forces/moment}}};
      \fi
    }
  },
  pics/mech/internal-forces/.default={},
}
\pgfkeys{
  /mech/fixed-reactions/.is family,
  /mech/fixed-reactions,
  x/.initial={$A_x$},
  y/.initial={$A_y$},
  moment/.initial={$M_A$},
  length/.initial=7mm,
}
\tikzset{
  pics/mech/fixed-reactions/.style args={#1}{
    code={
      \pgfkeys{/mech/fixed-reactions/.cd,#1}
      \edef\mechReactionLength{\pgfkeysvalueof{/mech/fixed-reactions/length}}
      \draw[mech/load,-{Latex[length=2mm]}] (-\mechReactionLength,0)--(0,0)
        node[pos=0,anchor=east]{\pgfkeysvalueof{/mech/fixed-reactions/x}};
      \draw[mech/load,-{Latex[length=2mm]}] (0,.7\mechReactionLength)--(0,-.3\mechReactionLength)
        node[pos=0,anchor=south]{\pgfkeysvalueof{/mech/fixed-reactions/y}};
      \pic {mech/moment={direction=ccw,radius=.45\mechReactionLength,label=\pgfkeysvalueof{/mech/fixed-reactions/moment}}};
    }
  },
  pics/mech/fixed-reactions/.default={},
}


%
% Lade-Reihenfolge:
%   1. core/        Farben, Layer, Stile und globale Geometriekonventionen
%   2. elements/    mechanische Grundelemente (Koerper, Lager, Verbinder, Staebe)
%   3. annotations/ Lasten, Masse, Koordinaten und mechanische Annotationen
%
% WICHTIG:
% Die Modulreihenfolge ist absichtlich festgelegt. Elemente greifen auf Farben,
% Layer und Stile aus core/ zurueck; Annotationen koennen wiederum die bereits
% definierten Grundstile und Pics verwenden.
% =========================================================
% -------------------------------------------------------------------------
% Externe Pakete und TikZ-Bibliotheken
% -------------------------------------------------------------------------
% \RequirePackage ist bewusst statt \usepackage gewaehlt, da mechlib.tex selbst
% wie eine Bibliotheks-/Paketdatei verwendet wird.
\RequirePackage{tikz}
\RequirePackage{amsmath}
\RequirePackage{siunitx}
\RequirePackage{xcolor}
\usetikzlibrary{
  arrows.meta,
  calc,
  patterns,
  patterns.meta,
  intersections,
  angles,
  quotes,
  decorations.pathmorphing,
  backgrounds
}
% Base path: override before loading when mechlib is stored elsewhere.
% -------------------------------------------------------------------------
% Aufloesung des Bibliothekspfades
% -------------------------------------------------------------------------
% \providecommand respektiert eine bereits im Hauptdokument gesetzte Definition.
\providecommand{\mechlibpath}{mechlib}
% -------------------------------------------------------------------------
% Kernmodule: muessen vor allen Elementen geladen werden.
% -------------------------------------------------------------------------
% =========================================================
% mechlib/core/colors.tex
% ZENTRALE FARBPALETTE
% =========================================================
% Alle bibliotheksweiten Farben werden ausschliesslich hier definiert. Dadurch
% kann der visuelle Charakter der gesamten mechlib an einer Stelle angepasst
% werden, ohne einzelne Elemente zu veraendern.
%
% Konvention:
%   mechBodyColor    Fuellung allgemeiner Koerper und Fachwerkgelenke
%   mechSupportColor Fuellung von Lagerkoerpern
%   mechLoadColor    Kraefte, Momente und Streckenlasten
%   mechAxisColor    Achsen und Koordinatensysteme
%   mechRodColor     historische Stabfarbe; aktuelle Fachwerkstaebe sind schwarz
%
% Die Namen sind Teil der internen visuellen API und sollten bevorzugt statt
% harter RGB-Werte in den Elementdateien verwendet werden.
% =========================================================
% -------------------------------------------------------------------------
% Farbdefinitionen
% -------------------------------------------------------------------------
% RGB wird bewusst explizit verwendet, damit die Darstellung unabhaengig von
% Dokumentklassen und externen Farbmodellen reproduzierbar bleibt.
\definecolor{mechBodyColor}{RGB}{180,180,180}
\definecolor{mechSupportColor}{RGB}{205,205,205}
\definecolor{mechLoadColor}{RGB}{178,25,25}
\definecolor{mechAxisColor}{RGB}{6,13,141}
\definecolor{mechRodColor}{RGB}{130,130,130}

% =========================================================
% mechlib/core/layers.tex
% ZEICHENEBENEN / PGF-LAYER
% =========================================================
% Fachwerkstaebe und andere hinterlegte Geometrien koennen auf dem Layer
% "mech background" gezeichnet werden. Der normale TikZ-Layer "main" liegt
% darueber. Dadurch koennen z. B. gefuellte Gelenke sauber vor den Staeben
% erscheinen, ohne dass die Zeichenreihenfolge im Anwendungsdokument aufwendig
% kontrolliert werden muss.
% =========================================================
% Zuerst wird der benannte Hintergrundlayer deklariert. Anschliessend wird die
% globale Layerreihenfolge explizit gesetzt: Hintergrund vor Hauptlayer.
\pgfdeclarelayer{mech background}
\pgfsetlayers{mech background,main}

% =========================================================
% mechlib/core/styles.tex
% ZENTRALE, NAMENSRAUMGESCHUETZTE TIKZ-STILE
% =========================================================
% Diese Datei definiert ausschliesslich visuelle Grundstile. Die eigentliche
% Geometrie eines Elements gehoert in elements/ bzw. annotations/.
%
% Alle oeffentlichen Bibliotheksstile beginnen mit "mech/". Das reduziert
% Namenskollisionen mit dem Hauptdokument und macht im TikZ-Code sichtbar,
% welche Formatierung aus der mechlib stammt.
%
% Designregel:
%   - Farben werden aus core/colors.tex bezogen.
%   - Geometrien werden NICHT hier definiert.
%   - Linienstarken werden hier zentral gehalten, soweit sie elementuebergreifend
%     konsistent sein sollen.
% =========================================================
% -------------------------------------------------------------------------
% Visuelle Basisstile
% -------------------------------------------------------------------------
% Einzelne Pics duerfen diese Stile ueber ihre jeweiligen "style"-Keys ersetzen.
\tikzset{
  mech/body/.style={
    draw=black,
    fill=mechBodyColor,
    line width=.6pt
  },
  mech/support/.style={
    draw=black,
    fill=mechSupportColor,
    line width=.6pt
  },
  mech/connector/.style={
    draw=black,
    line width=.6pt
  },
  mech/rod/.style={
    draw=black,
    line width=.85pt
  },
  mech/load/.style={
    draw=mechLoadColor,
    text=mechLoadColor,
    line width=1pt
  },
  mech/load fill/.style={
    fill=mechLoadColor,
    fill opacity=.16,
    draw=none
  },
  mech/dimension/.style={
    {Latex[length=2mm]}-{Latex[length=2mm]},
    line width=.5pt
  },
  mech/axis/.style={
    draw=mechAxisColor,
    text=mechAxisColor,
    line width=.6pt
  },
  mech/label/.style={
    inner sep=1pt
  }
}

% =========================================================
% mechlib/core/defaults.tex
% GLOBALE STANDARDWERTE UND GEMEINSAME GEOMETRIEKONVENTIONEN
% =========================================================
% Hier stehen Werte, die von mehreren Modulen gemeinsam benutzt werden. Solche
% Konstanten verhindern, dass geometrisch zusammengehoerige Elemente lokal mit
% leicht unterschiedlichen Zahlen definiert werden.
%
% Beispiel: \mechConnectionOffset legt die vertikale Lage paralleler
% Anschlussanker fest. Masse und Einspannung verwenden exakt denselben Wert;
% dadurch liegen obere/mittlere/untere Anschluesse bei gleicher Mittelpunktlage
% garantiert auf gleicher Hoehe.
% =========================================================
% Globale Schriftvorgabe fuer jedes TikZ-Bild. "append style" ergaenzt
% vorhandene Bildstile, statt sie zu ueberschreiben.
\tikzset{
  every picture/.append style={font=\fontsize{10}{10}\selectfont},
}

% Standardwinkel fuer technische Bodenschraffuren.
\def\mechHatchAngle{45}

% Shared vertical distance of parallel connector anchors from the middle anchor.
% Wall and mass use this exact value so upper/middle/lower are aligned whenever
% their center points lie on the same horizontal line.
% Gemeinsamer vertikaler Versatz fuer parallele Anschlussanker.
\def\mechConnectionOffset{2.5mm}

% -------------------------------------------------------------------------
% Atomare mechanische Elemente.
% -------------------------------------------------------------------------
% =========================================================
% mechlib/elements/bodies.tex
% ATOMARE KOERPER: MASSE, PUNKTMASSE, WAGEN UND BALKEN
% =========================================================
% Die hier definierten TikZ-Pics stellen mechanische Grundkoerper bereit. Jedes
% Pic besitzt eine eigene pgfkeys-Familie unter /mech/... und kann dadurch mit
% benannten Optionen konfiguriert werden.
%
% Grundprinzip der API:
%   \pic (name) at (<Ort>) {mech/<element>={<key>=<wert>,...}};
%
% Benannte Pics exportieren Koordinaten bzw. werden bei node-basierten Koerpern
% auf den internen Shape-Knoten aliasiert. So koennen regulaere TikZ-Anker wie
% (m1.east), (m1.north west) usw. direkt benutzt werden.
%
% Oeffentliche Pics in dieser Datei:
%   mech/mass       rechteckiger Starrkoerper bzw. konzentrierte Masse
%   mech/point-mass kreisfoermige Punktmasse
%   mech/cart       Wagen mit zwei Raedern und optionalem Boden
%   mech/beam       Balken zwischen zwei beliebigen Punkten
% =========================================================
\pgfkeys{
  /mech/mass/.is family,
  /mech/mass,
  width/.initial=1.5cm,
  height/.initial=1cm,
  label/.initial=,
  style/.initial=mech/body,
}
\tikzset{
  pics/mech/mass/.style args={#1}{
    code={
      \pgfkeys{/mech/mass/.cd,#1}
      \node[
        \pgfkeysvalueof{/mech/mass/style},
        minimum width=\pgfkeysvalueof{/mech/mass/width},
        minimum height=\pgfkeysvalueof{/mech/mass/height},
        inner sep=0pt
      ] (-shape) at (0,0) {\pgfkeysvalueof{/mech/mass/label}};

      % Named body pics behave like regular TikZ nodes.
      % Example: \pic (m1) ... then (m1.east), (m1.south east), ...
      \edef\mechPicName{\pgfkeysvalueof{/tikz/name prefix}}
      \ifx\mechPicName\pgfutil@empty\else
        \pgfnodealias{\mechPicName}{\mechPicName-shape}
      \fi

      \coordinate (-center) at (-shape.center);

      % Three standardized connection levels on each vertical side.
      % The shared \mechConnectionOffset is also used by fixed supports.
      \coordinate (-west-upper)  at ($(-shape.west)+(0,\mechConnectionOffset)$);
      \coordinate (-west-middle) at (-shape.west);
      \coordinate (-west-lower)  at ($(-shape.west)+(0,-\mechConnectionOffset)$);

      \coordinate (-east-upper)  at ($(-shape.east)+(0,\mechConnectionOffset)$);
      \coordinate (-east-middle) at (-shape.east);
      \coordinate (-east-lower)  at ($(-shape.east)+(0,-\mechConnectionOffset)$);

      % Generic connection aliases
      \coordinate (-connection-upper)  at (-east-upper);
      \coordinate (-connection-middle) at (-east-middle);
      \coordinate (-connection-lower)  at (-east-lower);

      \coordinate (-north) at (-shape.north);
      \coordinate (-south) at (-shape.south);
    }
  },
  pics/mech/mass/.default={},
}
\pgfkeys{
  /mech/point-mass/.is family,
  /mech/point-mass,
  radius/.initial=2.2mm,
  label/.initial=,
  style/.initial=mech/body,
}
\tikzset{
  pics/mech/point-mass/.style args={#1}{
    code={
      \pgfkeys{/mech/point-mass/.cd,#1}
      \edef\mechPointMassRadius{\pgfkeysvalueof{/mech/point-mass/radius}}
      \edef\mechPointMassLabel{\pgfkeysvalueof{/mech/point-mass/label}}
      \edef\mechPointMassStyle{\pgfkeysvalueof{/mech/point-mass/style}}
      \node[
        \mechPointMassStyle,
        circle,
        minimum size=2\mechPointMassRadius,
        inner sep=0pt
      ] (-shape) at (0,0) {\mechPointMassLabel};

      \edef\mechPicName{\pgfkeysvalueof{/tikz/name prefix}}
      \ifx\mechPicName\pgfutil@empty\else
        \pgfnodealias{\mechPicName}{\mechPicName-shape}
      \fi

      \coordinate (-center) at (-shape.center);
    }
  },
  pics/mech/point-mass/.default={},
}
% -------------------------------------------------------------------------
% Wagen
% -------------------------------------------------------------------------
% Der Boolean "ground" wird als TeX-if gespeichert. Dadurch kann die optionale
% Bodendarstellung ohne Stringvergleich ein- bzw. ausgeschaltet werden.
\newif\ifmechCartGround
\pgfkeys{
  /mech/cart/.is family,
  /mech/cart,
  width/.initial=2cm,
  height/.initial=1cm,
  wheel radius/.initial=2mm,
  label/.initial=,
  ground/.is if=mechCartGround,
  ground/.default=true,
  ground=true,
}
\tikzset{
  pics/mech/cart/.style args={#1}{
    code={
      \pgfkeys{/mech/cart/.cd,#1}
      \edef\mechCartWidth{\pgfkeysvalueof{/mech/cart/width}}
      \edef\mechCartHeight{\pgfkeysvalueof{/mech/cart/height}}
      \edef\mechCartWheelRadius{\pgfkeysvalueof{/mech/cart/wheel radius}}
      \node[
        mech/body,
        minimum width=\mechCartWidth,
        minimum height=\mechCartHeight,
        inner sep=0pt
      ] (-shape) at (0,0) {\pgfkeysvalueof{/mech/cart/label}};

      \edef\mechPicName{\pgfkeysvalueof{/tikz/name prefix}}
      \ifx\mechPicName\pgfutil@empty\else
        \pgfnodealias{\mechPicName}{\mechPicName-shape}
      \fi

      \coordinate (-center) at (-shape.center);
      \coordinate (-west) at (-shape.west);
      \coordinate (-east) at (-shape.east);
      \coordinate (-north) at (-shape.north);
      \coordinate (-south) at (-shape.south);
      \draw[mech/body]
        (-.3\mechCartWidth,-.5\mechCartHeight-\mechCartWheelRadius)
        circle[radius=\mechCartWheelRadius];
      \draw[mech/body]
        (.3\mechCartWidth,-.5\mechCartHeight-\mechCartWheelRadius)
        circle[radius=\mechCartWheelRadius];
      \ifmechCartGround
        \draw
          (-.5\mechCartWidth,-.5\mechCartHeight-2\mechCartWheelRadius)
          --
          (.5\mechCartWidth,-.5\mechCartHeight-2\mechCartWheelRadius);
        \fill[
          pattern={Lines[angle=\mechHatchAngle,distance=1pt]},
          pattern color=black
        ]
          (-.5\mechCartWidth,-.5\mechCartHeight-2\mechCartWheelRadius)
          rectangle
          (.5\mechCartWidth,-.5\mechCartHeight-2\mechCartWheelRadius-3pt);
      \fi
    }
  },
  pics/mech/cart/.default={},
}
\pgfkeys{
  /mech/beam/.is family,
  /mech/beam,
  from/.initial={(0,0)},
  to/.initial={(1,0)},
  height/.initial=3mm,
  style/.initial=mech/body,
}
\tikzset{
  pics/mech/beam/.style args={#1}{
    code={
      \pgfkeys{/mech/beam/.cd,#1}
      \edef\mechBeamFrom{\pgfkeysvalueof{/mech/beam/from}}
      \edef\mechBeamTo{\pgfkeysvalueof{/mech/beam/to}}
      \edef\mechBeamHeight{\pgfkeysvalueof{/mech/beam/height}}
      \edef\mechBeamStyle{\pgfkeysvalueof{/mech/beam/style}}
      \path let
        \p1=\mechBeamFrom,
        \p2=\mechBeamTo,
        \n1={veclen(\x2-\x1,\y2-\y1)},
        \n2={(\y2-\y1)/\n1*\mechBeamHeight/2},
        \n3={-(\x2-\x1)/\n1*\mechBeamHeight/2}
      in
        coordinate (-A) at (\x1+\n2,\y1+\n3)
        coordinate (-B) at (\x2+\n2,\y2+\n3)
        coordinate (-C) at (\x2-\n2,\y2-\n3)
        coordinate (-D) at (\x1-\n2,\y1-\n3)
        coordinate (-start) at (\x1,\y1)
        coordinate (-end) at (\x2,\y2)
        coordinate (-center) at ({(\x1+\x2)/2},{(\y1+\y2)/2});
      \draw[\mechBeamStyle] (-A)--(-B)--(-C)--(-D)--cycle;
    }
  },
  pics/mech/beam/.default={},
}

% =========================================================
% mechlib/elements/supports.tex
% ATOMARE LAGER: FESTLAGER, LOSLAGER UND EINSPANNUNG
% =========================================================
% Diese Datei kapselt die Lagergeometrie. Fest- und Loslager teilen sich die
% Key-Familie /mech/support-common, damit Breite, Hoehe, Schraffur und Label
% konsistent parametrisiert werden koennen.
%
% Oeffentliche Pics:
%   mech/pinned-support Festlager / gelenkiges Lager
%   mech/roller-support Loslager mit sichtbarem Abstand zum Boden
%   mech/fixed-support  Einspannung, optional gedreht
%
% Aktuelle Fachwerk-Konvention:
%   - Die Dreiecksspitze liegt exakt im Mittelpunkt des Gelenks (0,0).
%   - Das Dreieck ist gegenueber der frueheren Version kompakt skaliert.
%   - Bodenlinie und Schraffur verwenden dieselben x-Grenzen wie die
%     Dreiecksgrundseite; beim Loslager wird die sichtbare Linienbreite explizit
%     beruecksichtigt.
%   - Das Gelenk selbst ist mit mechBodyColor gefuellt.
%
% Die pgfmathsetlengthmacro-Aufrufe sind absichtlich verwendet: LaTeX-Laengen
% duerfen nicht durch Textersatz wie ".5\\macro" skaliert werden, da dies bei
% expandierenden Dimensionsmakros zu sichtbarem Resttext (z. B. ".8mm") fuehren
% kann. Alle abgeleiteten Abmessungen werden deshalb numerisch berechnet.
% =========================================================
\pgfkeys{
  /mech/support-common/.is family,
  /mech/support-common,
  width/.initial=2.72mm,
  height/.initial=2.24mm,
  hatch angle/.initial=45,
  label/.initial=,
  label position/.initial=above,
}
\tikzset{
  % -----------------------------------------------------------------------
  % Festlager
  % -----------------------------------------------------------------------
  pics/mech/pinned-support/.style args={#1}{
    code={
      \pgfkeys{/mech/support-common/.cd,#1}
      \edef\mechSupportWidth{\pgfkeysvalueof{/mech/support-common/width}}
      \edef\mechSupportHeight{\pgfkeysvalueof{/mech/support-common/height}}
      \edef\mechSupportHatch{\pgfkeysvalueof{/mech/support-common/hatch angle}}
      \pgfmathsetlengthmacro{\mechSupportHalfWidth}{0.5*\mechSupportWidth}
      % Die Dreieckskontur ragt an der Grundseite durch die halbe Linienbreite
      % minimal ueber die geometrische Breite hinaus. Fuer identische sichtbare
      % Breite wird die Bodenlinie daher um 0.3pt je Seite erweitert.
      \pgfmathsetlengthmacro{\mechSupportGroundHalfWidth}{\mechSupportHalfWidth+0.3pt}
      \pgfmathsetlengthmacro{\mechSupportApexY}{0mm}
      \pgfmathsetlengthmacro{\mechSupportHatchDepth}{0.256*\mechSupportHeight}

      \coordinate (-connection) at (0,0);
      \coordinate (-center) at (0,0);

      \draw[mech/support]
        (-\mechSupportHalfWidth,-\mechSupportHeight)
        -- (0,\mechSupportApexY)
        -- (\mechSupportHalfWidth,-\mechSupportHeight)
        -- cycle;
      \draw[mech/body] (0,0) circle[radius=.6mm];
      \draw[black,line width=.6pt,line cap=butt]
        (-\mechSupportGroundHalfWidth,-\mechSupportHeight)
        -- (\mechSupportGroundHalfWidth,-\mechSupportHeight);
      \fill[
        pattern={Lines[angle=\mechSupportHatch,distance=1pt]},
        pattern color=black
      ]
        (-\mechSupportGroundHalfWidth,-\mechSupportHeight)
        rectangle
        (\mechSupportGroundHalfWidth,{-(\mechSupportHeight+\mechSupportHatchDepth)});

      \node[\pgfkeysvalueof{/mech/support-common/label position}]
        at (0,0) {\pgfkeysvalueof{/mech/support-common/label}};
    }
  },
  pics/mech/pinned-support/.default={},
  % -----------------------------------------------------------------------
  % Loslager
  % -----------------------------------------------------------------------
  pics/mech/roller-support/.style args={#1}{
    code={
      \pgfkeys{/mech/support-common/.cd,#1}
      \edef\mechSupportWidth{\pgfkeysvalueof{/mech/support-common/width}}
      \edef\mechSupportHeight{\pgfkeysvalueof{/mech/support-common/height}}
      \edef\mechSupportHatch{\pgfkeysvalueof{/mech/support-common/hatch angle}}
      \pgfmathsetlengthmacro{\mechSupportHalfWidth}{0.5*\mechSupportWidth}
      % The triangle outline extends by half its line width at the base corners.
      % Add 0.3pt so the roller ground has the same visible total width.
      \pgfmathsetlengthmacro{\mechSupportGroundHalfWidth}{\mechSupportHalfWidth+0.3pt}
      \pgfmathsetlengthmacro{\mechSupportApexY}{0mm}
      \pgfmathsetlengthmacro{\mechSupportRollerGap}{0.272*\mechSupportHeight}
      \pgfmathsetlengthmacro{\mechSupportHatchDepth}{0.256*\mechSupportHeight}
      \pgfmathsetlengthmacro{\mechSupportGroundY}{\mechSupportHeight+\mechSupportRollerGap}

      \coordinate (-connection) at (0,0);
      \coordinate (-center) at (0,0);

      \draw[mech/support]
        (-\mechSupportHalfWidth,-\mechSupportHeight)
        -- (0,\mechSupportApexY)
        -- (\mechSupportHalfWidth,-\mechSupportHeight)
        -- cycle;
      \draw[mech/body] (0,0) circle[radius=.6mm];
      \draw[black,line width=.6pt,line cap=butt]
        (-\mechSupportGroundHalfWidth,-\mechSupportGroundY)
        -- (\mechSupportGroundHalfWidth,-\mechSupportGroundY);
      \fill[
        pattern={Lines[angle=\mechSupportHatch,distance=1pt]},
        pattern color=black
      ]
        (-\mechSupportGroundHalfWidth,-\mechSupportGroundY)
        rectangle
        (\mechSupportGroundHalfWidth,{-(\mechSupportGroundY+\mechSupportHatchDepth)});

      \node[\pgfkeysvalueof{/mech/support-common/label position}]
        at (0,0) {\pgfkeysvalueof{/mech/support-common/label}};
    }
  },
  pics/mech/roller-support/.default={},
}
\pgfkeys{
  /mech/fixed-support/.is family,
  /mech/fixed-support,
  width/.initial=1cm,
  depth/.initial=1.5mm,
  hatch angle/.initial=45,
  angle/.initial=0,
  label/.initial=,
}
\tikzset{
  pics/mech/fixed-support/.style args={#1}{
    code={
      \pgfkeys{/mech/fixed-support/.cd,#1}
      \edef\mechFixedWidth{\pgfkeysvalueof{/mech/fixed-support/width}}
      \edef\mechFixedDepth{\pgfkeysvalueof{/mech/fixed-support/depth}}
      \edef\mechFixedHatch{\pgfkeysvalueof{/mech/fixed-support/hatch angle}}
      \edef\mechFixedAngle{\pgfkeysvalueof{/mech/fixed-support/angle}}

      \coordinate (-center) at (0,0);
      \coordinate (-connection-upper)  at (0,\mechConnectionOffset);
      \coordinate (-connection-middle) at (0,0);
      \coordinate (-connection-lower)  at (0,-\mechConnectionOffset);

      \begin{scope}[rotate=\mechFixedAngle]
        \draw[mech/body,fill=none] (-.5\mechFixedWidth,0)--(.5\mechFixedWidth,0);
        \fill[
          pattern={Lines[angle=\mechFixedHatch,distance={3pt/sqrt(2)}]},
          pattern color=black
        ]
          (-.5\mechFixedWidth,0)
          rectangle
          (.5\mechFixedWidth,\mechFixedDepth);
      \end{scope}

      \node[above] at (0,0) {\pgfkeysvalueof{/mech/fixed-support/label}};
    }
  },
  pics/mech/fixed-support/.default={},
}

% =========================================================
% mechlib/elements/connectors.tex
% ATOMARE VERBINDUNGSELEMENTE: FEDER UND DAEMPFER
% =========================================================
% Beide Elemente folgen derselben grundlegenden API: "from", "to" und "label".
% Die Endpunkte koennen beliebige TikZ-Koordinaten oder exportierte Anker aus
% anderen mechlib-Pics sein.
%
% Oeffentliche Pics:
%   mech/spring  Zickzack-Feder entlang der Verbindungslinie
%   mech/damper  Daempfer mit rechteckigem Daempferkoerper in Linienmitte
%
% Die Beschriftung wird "sloped" gezeichnet und folgt daher der Orientierung
% des Verbindungselements. Ueber "label side" kann sie ober- oder unterhalb
% der lokalen Verbindungslinie platziert werden.
% =========================================================
\pgfkeys{
  /mech/spring/.is family,
  /mech/spring,
  from/.initial={(0,0)},
  to/.initial={(1,0)},
  label/.initial=,
  label side/.initial=above,
  amplitude/.initial=2mm,
  segment length/.initial=4mm,
  style/.initial=mech/connector,
}
\tikzset{
  pics/mech/spring/.style args={#1}{
    code={
      \pgfkeys{/mech/spring/.cd,#1}
      \edef\mechSpringFrom{\pgfkeysvalueof{/mech/spring/from}}
      \edef\mechSpringTo{\pgfkeysvalueof{/mech/spring/to}}
      \draw[\pgfkeysvalueof{/mech/spring/style},
        decorate,
        decoration={zigzag,
          segment length=\pgfkeysvalueof{/mech/spring/segment length},
          amplitude=\pgfkeysvalueof{/mech/spring/amplitude}}]
        \mechSpringFrom -- \mechSpringTo
        node[midway,sloped,mech/label,outer sep=2mm,\pgfkeysvalueof{/mech/spring/label side}]
        {\pgfkeysvalueof{/mech/spring/label}};
    }
  },
  pics/mech/spring/.default={},
}
\pgfkeys{
  /mech/damper/.is family,
  /mech/damper,
  from/.initial={(0,0)},
  to/.initial={(1,0)},
  label/.initial=,
  label side/.initial=below,
  body width/.initial=8mm,
  body height/.initial=3mm,
  style/.initial=mech/connector,
}
\tikzset{
  pics/mech/damper/.style args={#1}{
    code={
      \pgfkeys{/mech/damper/.cd,#1}
      \edef\mechDamperFrom{\pgfkeysvalueof{/mech/damper/from}}
      \edef\mechDamperTo{\pgfkeysvalueof{/mech/damper/to}}
      \draw[\pgfkeysvalueof{/mech/damper/style}]
        \mechDamperFrom -- \mechDamperTo
        node[midway,sloped,draw,fill=white,
          minimum width=\pgfkeysvalueof{/mech/damper/body width},
          minimum height=\pgfkeysvalueof{/mech/damper/body height},
          inner sep=0pt] {};
      \path \mechDamperFrom -- \mechDamperTo
        node[midway,sloped,mech/label,outer sep=2mm,\pgfkeysvalueof{/mech/damper/label side}]
        {\pgfkeysvalueof{/mech/damper/label}};
    }
  },
  pics/mech/damper/.default={},
}

% =========================================================
% mechlib/elements/trusses.tex
% FACHWERKE: STAB UND GEMEINSAME GELENKDARSTELLUNG
% =========================================================
% Dieses Modul enthaelt die atomaren Fachwerkelemente. Die Topologie eines
% Fachwerks bleibt bewusst im Anwendungsdokument: Knoten werden als TikZ-
% Koordinaten definiert und durch mech/rod bzw. \mechJoints verbunden.
%
% Oeffentliche API:
%   mech/rod     Stab zwischen "from" und "to" mit optionalem Label
%   \mechJoints{A,B,C,...} zeichnet alle angegebenen Gelenke
%
% Darstellungsregeln:
%   - Staebe verwenden standardmaessig den Stil mech/rod: schwarz und duenn.
%   - Staebe werden auf dem Layer "mech background" gezeichnet.
%   - Stabnummern besitzen keine Fuellung (fill=none) und bleiben transparent.
%   - Gelenke liegen optisch vor den Staeben und sind mit mechBodyColor gefuellt.
%
% Beispiel:
%   \pic {mech/rod={from=(A),to=(B),label=1,label side=above}};
%   \mechJoints{A,B}
% =========================================================
% -------------------------------------------------------------------------
% Einzelner Fachwerkstab
% -------------------------------------------------------------------------
% "style" erlaubt lokale Abweichungen, ohne den globalen mech/rod-Stil zu
% veraendern. Die Standardbeschriftung liegt an der Stabmitte.
\pgfkeys{
  /mech/rod/.is family,
  /mech/rod,
  from/.initial={(0,0)},
  to/.initial={(1,0)},
  label/.initial=,
  label side/.initial=above,
  label style/.initial={mech/label,fill=none,text=black},
  style/.initial=mech/rod,
}
\tikzset{
  pics/mech/rod/.style args={#1}{
    code={
      \pgfkeys{/mech/rod/.cd,#1}
      \edef\mechRodFrom{\pgfkeysvalueof{/mech/rod/from}}
      \edef\mechRodTo{\pgfkeysvalueof{/mech/rod/to}}
      \coordinate (-start) at \mechRodFrom;
      \coordinate (-end) at \mechRodTo;
      \coordinate (-center) at ($(-start)!.5!(-end)$);
      \begin{pgfonlayer}{mech background}
        \draw[\pgfkeysvalueof{/mech/rod/style}] (-start)--(-end);
      \end{pgfonlayer}
      \node[\pgfkeysvalueof{/mech/rod/label style},\pgfkeysvalueof{/mech/rod/label side}]
        at (-center) {\pgfkeysvalueof{/mech/rod/label}};
    }
  },
  pics/mech/rod/.default={},
}
% -------------------------------------------------------------------------
% Gelenkliste
% -------------------------------------------------------------------------
% Das Makro erwartet eine kommaseparierte Liste bereits definierter TikZ-
% Koordinaten. Jedes Gelenk wird identisch gezeichnet. Die Gelenke sollten nach
% den Staeben aufgerufen werden, damit sie die Stabenden optisch ueberdecken.
\newcommand{\mechJoints}[1]{%
  \foreach \mechJoint in {#1}{%
    \draw[fill=mechBodyColor,draw=black,line width=.2mm]
      (\mechJoint) circle[radius=.6mm];%
  }%
}

% -------------------------------------------------------------------------
% Annotationen und mechanische Symbolgruppen.
% -------------------------------------------------------------------------
% =========================================================
% mechlib/annotations/loads.tex
% LASTANNOTATIONEN: KRAFT, MOMENT UND STRECKENLAST
% =========================================================
% Diese Datei stellt die Lastsymbole der mechlib bereit. Alle Lasten verwenden
% standardmaessig den zentralen Stil mech/load und damit mechLoadColor.
%
% Oeffentliche Pics:
%   mech/force            Einzelkraft als Winkel/Laenge oder from/to
%   mech/moment           Kreisbogensymbol im oder gegen den Uhrzeigersinn
%   mech/distributed-load konstante, dreieckige oder linear veraenderliche Last
%
% Besonderheit der Streckenlast:
% Die Grundlinie wird lokal auf die x-Achse gedreht. Dadurch kann die gesamte
% Geometrie in einem einfachen lokalen Koordinatensystem berechnet werden. Sehr
% kurze Pfeile, deren Laenge kleiner als ihre Pfeilspitze ist, werden bewusst
% unterdrueckt; so entstehen am Nullende einer Dreieckslast keine Artefakte.
% =========================================================
\pgfkeys{
  /mech/force/.is family,
  /mech/force,
  from/.initial={(0,0)},
  to/.initial=,
  angle/.initial=0,
  length/.initial=1.5cm,
  label/.initial=,
  label position/.initial=above,
  direction/.initial=away,
  style/.initial=mech/load,
}
\tikzset{
  pics/mech/force/.style args={#1}{
    code={
      \pgfkeys{/mech/force/.cd,#1}
      \edef\mechForceFrom{\pgfkeysvalueof{/mech/force/from}}
      \edef\mechForceTo{\pgfkeysvalueof{/mech/force/to}}
      \edef\mechForceDirection{\pgfkeysvalueof{/mech/force/direction}}
      \def\mechForceToward{toward}
      \ifx\mechForceTo\empty
        \ifx\mechForceDirection\mechForceToward
          \draw[\pgfkeysvalueof{/mech/force/style},-{Latex[length=2.5mm]}]
            (\pgfkeysvalueof{/mech/force/angle}:\pgfkeysvalueof{/mech/force/length}) -- (0,0)
            node[midway,\pgfkeysvalueof{/mech/force/label position}]
            {\pgfkeysvalueof{/mech/force/label}};
        \else
          \draw[\pgfkeysvalueof{/mech/force/style},-{Latex[length=2.5mm]}]
            (0,0) -- (\pgfkeysvalueof{/mech/force/angle}:\pgfkeysvalueof{/mech/force/length})
            node[midway,\pgfkeysvalueof{/mech/force/label position}]
            {\pgfkeysvalueof{/mech/force/label}};
        \fi
      \else
        \ifx\mechForceDirection\mechForceToward
          \draw[\pgfkeysvalueof{/mech/force/style},-{Latex[length=2.5mm]}]
            \mechForceTo -- \mechForceFrom
            node[midway,\pgfkeysvalueof{/mech/force/label position}]
            {\pgfkeysvalueof{/mech/force/label}};
        \else
          \draw[\pgfkeysvalueof{/mech/force/style},-{Latex[length=2.5mm]}]
            \mechForceFrom -- \mechForceTo
            node[midway,\pgfkeysvalueof{/mech/force/label position}]
            {\pgfkeysvalueof{/mech/force/label}};
        \fi
      \fi
    }
  },
  pics/mech/force/.default={},
}
\pgfkeys{
  /mech/moment/.is family,
  /mech/moment,
  radius/.initial=6mm,
  label/.initial=,
  direction/.is choice,
  direction/ccw/.code={\def\mechMomentStart{-60}\def\mechMomentEnd{220}},
  direction/cw/.code={\def\mechMomentStart{240}\def\mechMomentEnd{-40}},
  direction=ccw,
  style/.initial=mech/load,
}
\tikzset{
  pics/mech/moment/.style args={#1}{
    code={
      \pgfkeys{/mech/moment/.cd,direction=ccw,#1}
      \draw[\pgfkeysvalueof{/mech/moment/style},-{Latex[length=2.5mm]}]
        (\mechMomentStart:\pgfkeysvalueof{/mech/moment/radius})
        arc[start angle=\mechMomentStart,end angle=\mechMomentEnd,
            radius=\pgfkeysvalueof{/mech/moment/radius}];
      \node[mech/label,above] at (0,\pgfkeysvalueof{/mech/moment/radius})
        {\pgfkeysvalueof{/mech/moment/label}};
    }
  },
  pics/mech/moment/.default={},
}
\pgfkeys{
  /mech/distributed-load/.is family,
  /mech/distributed-load,
  from/.initial={(0,0)},
  to/.initial={(1,0)},
  start value/.initial=1,
  end value/.initial=1,
  height/.initial=8mm,
  offset/.initial=0mm,
  arrows/.initial=6,
  arrowhead length/.initial=2mm,
  label/.initial=,
  label position/.initial=above,
  style/.initial=mech/load,
}
\tikzset{
  pics/mech/distributed-load/.style args={#1}{
    code={
      \pgfkeys{/mech/distributed-load/.cd,#1}
      \edef\mechDLFrom{\pgfkeysvalueof{/mech/distributed-load/from}}
      \edef\mechDLTo{\pgfkeysvalueof{/mech/distributed-load/to}}
      \path let
        \p1=\mechDLFrom,\p2=\mechDLTo,
        \n1={veclen(\x2-\x1,\y2-\y1)},
        \n2={atan2(\y2-\y1,\x2-\x1)}
      in
        coordinate (-start) at (\x1,\y1)
        coordinate (-end) at (\x2,\y2)
        coordinate (-helper) at (\n2,\n1);
      \path (-helper); \pgfgetlastxy{\mechDLAngle}{\mechDLLength};
      \pgfmathsetmacro{\mechDLStartValue}{\pgfkeysvalueof{/mech/distributed-load/start value}}
      \pgfmathsetmacro{\mechDLEndValue}{\pgfkeysvalueof{/mech/distributed-load/end value}}
      \pgfmathsetmacro{\mechDLArrowCount}{\pgfkeysvalueof{/mech/distributed-load/arrows}}
      \begin{scope}[shift={(-start)},rotate=\mechDLAngle]
        \pgfmathsetlengthmacro{\mechDLOffset}{\pgfkeysvalueof{/mech/distributed-load/offset}}
        \pgfmathsetlengthmacro{\mechDLArrowheadLength}{\pgfkeysvalueof{/mech/distributed-load/arrowhead length}}
        \pgfmathsetlengthmacro{\mechDLStartArrowLength}{abs(\mechDLStartValue)*\pgfkeysvalueof{/mech/distributed-load/height}}

        \fill[mech/load fill]
          (0,\mechDLOffset)
          --(0,{\mechDLOffset+\mechDLStartValue*\pgfkeysvalueof{/mech/distributed-load/height}})
          --(\mechDLLength,{\mechDLOffset+\mechDLEndValue*\pgfkeysvalueof{/mech/distributed-load/height}})
          --(\mechDLLength,\mechDLOffset)
          --cycle;

        % Keep contour and first arrow continuous whenever the first arrow
        % is long enough to accommodate its arrowhead. Otherwise draw only
        % the contour and suppress the too-short arrow.
        \ifdim\mechDLStartArrowLength<\mechDLArrowheadLength
          \draw[\pgfkeysvalueof{/mech/distributed-load/style}]
            (\mechDLLength,{\mechDLOffset+\mechDLEndValue*\pgfkeysvalueof{/mech/distributed-load/height}})
            --(0,{\mechDLOffset+\mechDLStartValue*\pgfkeysvalueof{/mech/distributed-load/height}});
        \else
          \draw[
            \pgfkeysvalueof{/mech/distributed-load/style},
            -{Latex[length=\pgfkeysvalueof{/mech/distributed-load/arrowhead length}]}
          ]
            (\mechDLLength,{\mechDLOffset+\mechDLEndValue*\pgfkeysvalueof{/mech/distributed-load/height}})
            --(0,{\mechDLOffset+\mechDLStartValue*\pgfkeysvalueof{/mech/distributed-load/height}})
            --(0,\mechDLOffset);
        \fi

        \foreach \mechDLIndex in {1,...,\mechDLArrowCount}{
          \pgfmathsetmacro{\mechDLT}{\mechDLIndex/\mechDLArrowCount}
          \pgfmathsetmacro{\mechDLValue}{\mechDLStartValue+(\mechDLEndValue-\mechDLStartValue)*\mechDLT}
          \pgfmathsetlengthmacro{\mechDLArrowLength}{abs(\mechDLValue)*\pgfkeysvalueof{/mech/distributed-load/height}}

          \ifdim\mechDLArrowLength<\mechDLArrowheadLength
            % Suppress arrows that are shorter than their arrowhead.
          \else
            \draw[
              \pgfkeysvalueof{/mech/distributed-load/style},
              -{Latex[length=\pgfkeysvalueof{/mech/distributed-load/arrowhead length}]}
            ]
              ({\mechDLT*\mechDLLength},{\mechDLOffset+\mechDLValue*\pgfkeysvalueof{/mech/distributed-load/height}})
              --({\mechDLT*\mechDLLength},\mechDLOffset);
          \fi
        }

        \node[mech/label,\pgfkeysvalueof{/mech/distributed-load/label position}]
          at (\mechDLLength,{\mechDLOffset+\mechDLEndValue*\pgfkeysvalueof{/mech/distributed-load/height}})
          {\pgfkeysvalueof{/mech/distributed-load/label}};
      \end{scope}
    }
  },
  pics/mech/distributed-load/.default={},
}

% =========================================================
% mechlib/annotations/dimensions.tex
% BEMASSUNGEN: LINEARES MASS UND WINKELMASS
% =========================================================
% Lineare Bemassungen werden geometrisch aus zwei Endpunkten abgeleitet. Der
% Offset wirkt senkrecht auf die Verbindung von "from" nach "to"; damit
% funktioniert dasselbe Pic fuer horizontale, vertikale und schraege Masse.
%
% Oeffentliche Pics:
%   mech/dimension       lineare Bemassung mit Doppelpfeil
%   mech/angle-dimension Kreisbogen zwischen zwei Winkeln
%
% Interne Koordinaten -A und -B sind die senkrecht versetzten Endpunkte der
% Masslinie. Die Vektorrechnung erfolgt in TikZ "let"-Syntax.
% =========================================================
\pgfkeys{
  /mech/dimension/.is family,
  /mech/dimension,
  from/.initial={(0,0)},
  to/.initial={(1,0)},
  label/.initial={$l$},
  offset/.initial=3mm,
  style/.initial=mech/dimension,
}
\tikzset{
  pics/mech/dimension/.style args={#1}{
    code={
      \pgfkeys{/mech/dimension/.cd,#1}
      \edef\mechDimFrom{\pgfkeysvalueof{/mech/dimension/from}}
      \edef\mechDimTo{\pgfkeysvalueof{/mech/dimension/to}}
      \path let
        \p1=\mechDimFrom,\p2=\mechDimTo,
        \n1={veclen(\x2-\x1,\y2-\y1)},
        \n2={-(\y2-\y1)/\n1*\pgfkeysvalueof{/mech/dimension/offset}},
        \n3={(\x2-\x1)/\n1*\pgfkeysvalueof{/mech/dimension/offset}}
      in
        coordinate (-A) at (\x1+\n2,\y1+\n3)
        coordinate (-B) at (\x2+\n2,\y2+\n3);
      \draw[\pgfkeysvalueof{/mech/dimension/style}]
        (-A)--(-B)
        node[midway,fill=white,inner sep=1pt]
        {\pgfkeysvalueof{/mech/dimension/label}};
    }
  },
  pics/mech/dimension/.default={},
}
\pgfkeys{
  /mech/angle-dimension/.is family,
  /mech/angle-dimension,
  start angle/.initial=0,
  end angle/.initial=45,
  radius/.initial=7mm,
  label/.initial={$\alpha$},
}
\tikzset{
  pics/mech/angle-dimension/.style args={#1}{
    code={
      \pgfkeys{/mech/angle-dimension/.cd,#1}
      \draw
        (\pgfkeysvalueof{/mech/angle-dimension/start angle}:\pgfkeysvalueof{/mech/angle-dimension/radius})
        arc[start angle=\pgfkeysvalueof{/mech/angle-dimension/start angle},
            end angle=\pgfkeysvalueof{/mech/angle-dimension/end angle},
            radius=\pgfkeysvalueof{/mech/angle-dimension/radius}];
      \pgfmathsetmacro{\mechAngleMid}{.5*(\pgfkeysvalueof{/mech/angle-dimension/start angle}+\pgfkeysvalueof{/mech/angle-dimension/end angle})}
      \node at (\mechAngleMid:{.7*\pgfkeysvalueof{/mech/angle-dimension/radius}})
        {\pgfkeysvalueof{/mech/angle-dimension/label}};
    }
  },
  pics/mech/angle-dimension/.default={},
}

% =========================================================
% mechlib/annotations/coordinates.tex
% KOORDINATENANNOTATIONEN: EINZELACHSE UND 2D-KOORDINATENSYSTEM
% =========================================================
% Die Pics dieser Datei zeichnen mathematische Bezugsachsen im Stil mech/axis.
% Winkel werden in TikZ-Grad angegeben; Laengen sind regulaere TeX-Dimensionen.
%
% Oeffentliche Pics:
%   mech/axis              einzelne gerichtete Achse
%   mech/coordinate-system zwei unabhaengig orientierbare Achsen
%
% Beide Pics exportieren den lokalen Ursprung als "-origin". Bei einem benannten
% Pic (cs) kann dieser z. B. als (cs-origin) referenziert werden.
% =========================================================
\pgfkeys{
  /mech/axis/.is family,
  /mech/axis,
  label/.initial={$x$},
  angle/.initial=0,
  length/.initial=1cm,
  direction/.initial=positive,
}
\tikzset{
  pics/mech/axis/.style args={#1}{
    code={
      \pgfkeys{/mech/axis/.cd,#1}
      \edef\mechAxisDirection{\pgfkeysvalueof{/mech/axis/direction}}
      \def\mechAxisNegative{negative}
      \ifx\mechAxisDirection\mechAxisNegative
        \draw[mech/axis,{Latex[length=2mm]}-]
          (0,0)--(\pgfkeysvalueof{/mech/axis/angle}:\pgfkeysvalueof{/mech/axis/length})
          node[pos=1,anchor=west]{\pgfkeysvalueof{/mech/axis/label}};
      \else
        \draw[mech/axis,-{Latex[length=2mm]}]
          (0,0)--(\pgfkeysvalueof{/mech/axis/angle}:\pgfkeysvalueof{/mech/axis/length})
          node[pos=1,anchor=west]{\pgfkeysvalueof{/mech/axis/label}};
      \fi
      \coordinate (-origin) at (0,0);
    }
  },
  pics/mech/axis/.default={},
}
\pgfkeys{
  /mech/coordinate-system/.is family,
  /mech/coordinate-system,
  x label/.initial={$x$},
  y label/.initial={$y$},
  x angle/.initial=0,
  y angle/.initial=90,
  x length/.initial=1cm,
  y length/.initial=1cm,
}
\tikzset{
  pics/mech/coordinate-system/.style args={#1}{
    code={
      \pgfkeys{/mech/coordinate-system/.cd,#1}
      \draw[mech/axis,-{Latex[length=2mm]}]
        (0,0)--(\pgfkeysvalueof{/mech/coordinate-system/x angle}:\pgfkeysvalueof{/mech/coordinate-system/x length})
        node[pos=1,anchor=west]{\pgfkeysvalueof{/mech/coordinate-system/x label}};
      \draw[mech/axis,-{Latex[length=2mm]}]
        (0,0)--(\pgfkeysvalueof{/mech/coordinate-system/y angle}:\pgfkeysvalueof{/mech/coordinate-system/y length})
        node[pos=1,anchor=east]{\pgfkeysvalueof{/mech/coordinate-system/y label}};
      \coordinate (-origin) at (0,0);
    }
  },
  pics/mech/coordinate-system/.default={},
}

% =========================================================
% mechlib/annotations/mechanics.tex
% MECHANISCHE ANNOTATIONEN: SCHNITTGROESSEN UND LAGERREAKTIONEN
% =========================================================
% Die Pics kombinieren mehrere elementare Lastsymbole zu haeufig benoetigten
% mechanischen Symbolgruppen. Farben und Pfeilgestaltung folgen mech/load.
%
% Oeffentliche Pics:
%   mech/internal-forces  Normalkraft N, Querkraft Q und Moment M am Schnitt
%   mech/fixed-reactions  horizontale/vertikale Reaktion und Einspannmoment
%
% Bei mech/internal-forces steuert "side=right|left" die Vorzeichenrichtung
% aller drei Schnittgroessen konsistent. Das Moment wird dazu intern als
% mech/moment mit passender Drehrichtung wiederverwendet.
% =========================================================
\pgfkeys{
  /mech/internal-forces/.is family,
  /mech/internal-forces,
  normal/.initial={$N$},
  shear/.initial={$Q$},
  moment/.initial={$M$},
  length/.initial=7mm,
  side/.is choice,
  side/right/.code={\def\mechIFSign{1}},
  side/left/.code={\def\mechIFSign{-1}},
  side=right,
}
\tikzset{
  pics/mech/internal-forces/.style args={#1}{
    code={
      \pgfkeys{/mech/internal-forces/.cd,side=right,#1}
      \edef\mechIFLength{\pgfkeysvalueof{/mech/internal-forces/length}}
      \draw[mech/load,-{Latex[length=2mm]}]
        (0,0)--({\mechIFSign*\mechIFLength},0)
        node[pos=1,anchor=west]{\pgfkeysvalueof{/mech/internal-forces/normal}};
      \draw[mech/load,-{Latex[length=2mm]}]
        (0,0)--(0,{-.8*\mechIFSign*\mechIFLength})
        node[pos=1,anchor=west]{\pgfkeysvalueof{/mech/internal-forces/shear}};
      \ifnum\mechIFSign=1
        \pic {mech/moment={direction=ccw,radius=.55\mechIFLength,label=\pgfkeysvalueof{/mech/internal-forces/moment}}};
      \else
        \pic {mech/moment={direction=cw,radius=.55\mechIFLength,label=\pgfkeysvalueof{/mech/internal-forces/moment}}};
      \fi
    }
  },
  pics/mech/internal-forces/.default={},
}
\pgfkeys{
  /mech/fixed-reactions/.is family,
  /mech/fixed-reactions,
  x/.initial={$A_x$},
  y/.initial={$A_y$},
  moment/.initial={$M_A$},
  length/.initial=7mm,
}
\tikzset{
  pics/mech/fixed-reactions/.style args={#1}{
    code={
      \pgfkeys{/mech/fixed-reactions/.cd,#1}
      \edef\mechReactionLength{\pgfkeysvalueof{/mech/fixed-reactions/length}}
      \draw[mech/load,-{Latex[length=2mm]}] (-\mechReactionLength,0)--(0,0)
        node[pos=0,anchor=east]{\pgfkeysvalueof{/mech/fixed-reactions/x}};
      \draw[mech/load,-{Latex[length=2mm]}] (0,.7\mechReactionLength)--(0,-.3\mechReactionLength)
        node[pos=0,anchor=south]{\pgfkeysvalueof{/mech/fixed-reactions/y}};
      \pic {mech/moment={direction=ccw,radius=.45\mechReactionLength,label=\pgfkeysvalueof{/mech/fixed-reactions/moment}}};
    }
  },
  pics/mech/fixed-reactions/.default={},
}

